% Modes de transfert de chaleur et calorimétrie — Physique 1ère E — Cours signé OnBuch+
% Programme officiel MINESEC. Compiler avec : tectonic N06-transferts-thermiques-calorimetrie.tex
\newcommand{\DOCMATIERE}{Physique}
\newcommand{\DOCNIVEAU}{1ère E}
\newcommand{\DOCMODULE}{Module 2 : Mouvements et interactions mécaniques}
\newcommand{\DOCLECON}{N06}
\newcommand{\DOCTITRE}{Modes de transfert de chaleur et calorimétrie}
% OnBuch+ — préambule « cours signés » ENRICHI (compilé avec tectonic / XeLaTeX)
% Système visuel « Pop » d'OnBuch+ : fond crème (jamais blanc), bordures encre
% épaisses, ombres « dures » décalées, titres Archivo Black en capitales.
% Version MATHÉMATIQUES : graphes (pgfplots/tikz), boîtes pédagogiques
% (définition, propriété, méthode, attention, le savais-tu, exercice résolu,
% exercice, TP, bilan), page de garde et signature OnBuch+ ancrée en bas.
%
% Macros attendues dans main.tex AVANT \input{preamble} :
%   \DOCMATIERE  \DOCNIVEAU  \DOCMODULE  \DOCLECON (numéro)  \DOCTITRE
\documentclass[11pt]{article}

\usepackage{fontspec}
\usepackage[french]{babel}
\usepackage[a4paper,margin=2cm,top=3.4cm,bottom=2.8cm,headheight=1.4cm,footskip=1.3cm]{geometry}
\usepackage{amsmath,amssymb,amsfonts}
\usepackage{mathtools} % \xmapsto, \coloneqq... souvent utilisés par les modèles
\usepackage{cancel} % simplifications barrées (bilans de forces)
% \abs{z} (module, valeur absolue) et \norm{u} (norme), utilisés par les modèles
\providecommand{\abs}{}\let\abs\relax\DeclarePairedDelimiter{\abs}{\lvert}{\rvert}
\providecommand{\norm}{}\let\norm\relax\DeclarePairedDelimiter{\norm}{\lVert}{\rVert}
\usepackage[table]{xcolor} % [table] : fournit aussi \rowcolors (alternance de lignes)
\usepackage{pagecolor}
\usepackage{tikz}
\usetikzlibrary{arrows.meta,positioning,calc,shapes.geometric,shapes.misc,shapes.arrows,decorations.pathmorphing,decorations.pathreplacing,decorations.markings,patterns,shadows,mindmap,trees,fit,backgrounds,matrix,intersections,angles,quotes}
\usepackage{pgfplots}
\usepackage[version=4]{mhchem} % équations nucléaires \ce{^{235}_{92}U}
\usepackage[european,siunitx]{circuitikz} % schémas électriques
\pgfplotsset{compat=1.18}
\usepgfplotslibrary{fillbetween,groupplots} % aires entre courbes (calcul intégral)
\usepackage{enumitem}
\usepackage{fancyhdr}
% [most] charge toutes les bibliothèques tcolorbox (listings, minted, raster,
% documentation...) et gonfle inutilement la mémoire TeX sur un document long
% (~300 boîtes) : on ne charge que ce qui est réellement utilisé.
\usepackage[skins,breakable]{tcolorbox}
\usepackage{graphicx}
\usepackage{colortbl}
\usepackage{booktabs}
\usepackage{tabularx}
\usepackage{diagbox} % case d'en-tête diagonale des tableaux à double entrée
% Colonnes extensibles centrée / gauche / droite (C, L, R), souvent utilisées par les modèles
\newcolumntype{C}{>{\centering\arraybackslash}X}
\newcolumntype{L}{>{\raggedright\arraybackslash}X}
\newcolumntype{R}{>{\raggedleft\arraybackslash}X}
\usepackage{array}
\usepackage{multicol}
\usepackage{siunitx}
% parse-numbers=false : accepte aussi \SI{2,5 \times 10^{-3}}{...}
\sisetup{locale=FR,output-decimal-marker={,},group-minimum-digits=4,parse-numbers=false}
\DeclareSIUnit\ohm{\ensuremath{\symup{\Omega}}} % U+2126 absent de la police
\DeclareSIUnit\division{div} % réglages d'oscilloscope (ms/div, V/div)
\DeclareSIUnit\tr{tr} % tours (vitesse de rotation, ex. \SI{1500}{\tr\per\minute})
\DeclareSIUnit\ch{ch} % cheval-vapeur (puissance moteur, unité usuelle hors SI)
\DeclareSIUnit\franccfa{FCFA} % franc CFA (coûts, contexte camerounais)
\DeclareSIUnit\dioptre{\ensuremath{\delta}} % vergence d'une lentille (dioptrie)
\usepackage{qrcode}
% \degree : commande fréquemment utilisée par les modèles pour un angle ou une
% température en dehors de \SI{}{} (non fournie par les paquets chargés).
\providecommand{\degree}{^{\circ}}
% \qed (fin de démonstration), utilisé par les modèles sans amsthm
\providecommand{\qed}{\hfill$\square$}
% \barre : double barre des valeurs interdites dans un tableau de variations (tkz-tab)
\providecommand{\barre}{\ensuremath{\|}}
% environnement proof (amsthm non chargé) utilisé par les modèles
\ifdefined\proof\else\newenvironment{proof}[1][Démonstration]{\par\noindent\textit{#1.}\ }{\qed\par}\fi
\usepackage{lastpage}
\usepackage{hyperref}
\hypersetup{hidelinks}

% ============================================================
% Palette « Pop » OnBuch+ — mêmes hex que la classe P (plus_ui.dart)
% ============================================================
\definecolor{poporange}{HTML}{F59321}
\definecolor{popdark}{HTML}{E07A0C}
\definecolor{popcream}{HTML}{FFF6E8}
\definecolor{popink}{HTML}{1C1714}
\definecolor{poppurple}{HTML}{7A5AE0}
\definecolor{popgreen}{HTML}{2E9E6A}
\definecolor{poppink}{HTML}{E0457B}
\definecolor{popblue}{HTML}{2F7DE1}
\definecolor{popgold}{HTML}{C98A12}
\definecolor{popmuted}{HTML}{8A7F76}
\definecolor{popline}{HTML}{EADFD2}
\definecolor{popgray}{HTML}{8A7F76}
\colorlet{popgrey}{popgray}
% Alias tolérants (le modèle nomme parfois les couleurs différemment)
\colorlet{popwhite}{white}
\colorlet{popblack}{popink}
\colorlet{popred}{poppink}
\colorlet{popyellow}{popgold}
% Teintes claires pour fonds de tableaux / zones
\colorlet{popblueL}{popblue!10!white}
\colorlet{poporangeL}{poporange!14!white}
\colorlet{popgreenL}{popgreen!12!white}
\colorlet{poppurpleL}{poppurple!12!white}
\colorlet{poppinkL}{poppink!12!white}
\colorlet{popgoldL}{popgold!14!white}

\pagecolor{popcream}

% ============================================================
% Typographie
% ============================================================
\defaultfontfeatures{Ligatures=TeX}
\setmainfont{PlusJakartaSans-Medium.ttf}[Path=../../../fonts/,BoldFont=PlusJakartaSans-Bold.ttf]
% Maths en sans-serif (Fira Math) avec chiffres et lettres droites de Plus
% Jakarta, pour que les formules se fondent dans le texte.
\usepackage[math-style=ISO,bold-style=ISO]{unicode-math}
\setmathfont{FiraMath-Regular.otf}[Path=../../../fonts/]
\setmathfont{PlusJakartaSans-Medium.ttf}[Path=../../../fonts/,range={up/num,up/{Latin,latin},\mathup}]
\usepackage{icomma} % virgule décimale sans espace en mode math
% Caractères mathématiques Unicode absents des polices de texte (Plus Jakarta,
% Archivo Black) : rendus par leur équivalent mathématique, en texte comme en
% formule — sinon ils s'affichent en carré vide (ex. « Arithmétique dans ℤ »).
\usepackage{newunicodechar}
\newunicodechar{ℝ}{\ensuremath{\mathbb{R}}}
\newunicodechar{ℤ}{\ensuremath{\mathbb{Z}}}
\newunicodechar{ℂ}{\ensuremath{\mathbb{C}}}
\newunicodechar{ℕ}{\ensuremath{\mathbb{N}}}
\newunicodechar{ℚ}{\ensuremath{\mathbb{Q}}}
\newunicodechar{𝔻}{\ensuremath{\mathbb{D}}}
\newunicodechar{α}{\ensuremath{\alpha}}
\newunicodechar{β}{\ensuremath{\beta}}
\newunicodechar{γ}{\ensuremath{\gamma}}
\newunicodechar{δ}{\ensuremath{\delta}}
\newunicodechar{ε}{\ensuremath{\varepsilon}}
\newunicodechar{θ}{\ensuremath{\theta}}
\newunicodechar{λ}{\ensuremath{\lambda}}
\newunicodechar{μ}{\ensuremath{\mu}}
\newunicodechar{σ}{\ensuremath{\sigma}}
\newunicodechar{τ}{\ensuremath{\tau}}
\newunicodechar{φ}{\ensuremath{\varphi}}
\newunicodechar{ψ}{\ensuremath{\psi}}
\newunicodechar{ω}{\ensuremath{\omega}}
\newunicodechar{Σ}{\ensuremath{\Sigma}}
% majuscules produites par \MakeUppercase dans les titres (α-aminés -> Α)
\newunicodechar{Α}{\ensuremath{\alpha}}
\newunicodechar{Β}{\ensuremath{\beta}}
\newunicodechar{Γ}{\ensuremath{\Gamma}}
\newunicodechar{Δ}{\ensuremath{\Delta}}
\newunicodechar{Ε}{\ensuremath{\varepsilon}}
\newunicodechar{Θ}{\ensuremath{\Theta}}
\newunicodechar{Λ}{\ensuremath{\Lambda}}
\newunicodechar{Μ}{\ensuremath{\mu}}
\newunicodechar{Τ}{\ensuremath{\tau}}
\newunicodechar{Φ}{\ensuremath{\Phi}}
\newunicodechar{Ψ}{\ensuremath{\Psi}}
\newunicodechar{Ω}{\ensuremath{\Omega}}
\newunicodechar{↦}{\ensuremath{\mapsto}}
\newunicodechar{∈}{\ensuremath{\in}}
\newunicodechar{∉}{\ensuremath{\notin}}
\newunicodechar{∧}{\ensuremath{\wedge}}
\newunicodechar{⇔}{\ensuremath{\Leftrightarrow}}
\newunicodechar{⟺}{\ensuremath{\Longleftrightarrow}}
\newunicodechar{⇒}{\ensuremath{\Rightarrow}}
\newunicodechar{⟹}{\ensuremath{\Longrightarrow}}
\newunicodechar{∘}{\ensuremath{\circ}}
\newunicodechar{∪}{\ensuremath{\cup}}
\newunicodechar{∩}{\ensuremath{\cap}}
\newunicodechar{≡}{\ensuremath{\equiv}}
\newunicodechar{⊂}{\ensuremath{\subset}}
\newunicodechar{⊥}{\ensuremath{\perp}}
\newunicodechar{∃}{\ensuremath{\exists}}
\newunicodechar{∀}{\ensuremath{\forall}}
\newunicodechar{∛}{\ensuremath{\sqrt[3]{\,}}}
\newunicodechar{∜}{\ensuremath{\sqrt[4]{\,}}}
\newunicodechar{⃗}{\ensuremath{{}^{\to}}}
\newunicodechar{ⁿ}{\ifmmode^{n}\else\textsuperscript{\ensuremath{n}}\fi}
\newunicodechar{ˣ}{\ifmmode^{x}\else\textsuperscript{\ensuremath{x}}\fi}
\newunicodechar{ᵇ}{\ifmmode^{b}\else\textsuperscript{\ensuremath{b}}\fi}
\newunicodechar{ʳ}{\ifmmode^{r}\else\textsuperscript{\ensuremath{r}}\fi}
\newunicodechar{ᵘ}{\ifmmode^{u}\else\textsuperscript{\ensuremath{u}}\fi}
\newunicodechar{ʸ}{\ifmmode^{y}\else\textsuperscript{\ensuremath{y}}\fi}
\newunicodechar{ᵃ}{\ifmmode^{a}\else\textsuperscript{\ensuremath{a}}\fi}
\newunicodechar{ᵉ}{\ifmmode^{e}\else\textsuperscript{\ensuremath{e}}\fi}
\newunicodechar{ᵏ}{\ifmmode^{k}\else\textsuperscript{\ensuremath{k}}\fi}
\newunicodechar{ᵖ}{\ifmmode^{p}\else\textsuperscript{\ensuremath{p}}\fi}
\newunicodechar{ᴺ}{\ifmmode^{N}\else\textsuperscript{\ensuremath{N}}\fi}
\newunicodechar{ᶜ}{\ifmmode^{c}\else\textsuperscript{\ensuremath{c}}\fi}
\newunicodechar{ʰ}{\ifmmode^{h}\else\textsuperscript{\ensuremath{h}}\fi}
\newunicodechar{ᵠ}{\ifmmode^{q}\else\textsuperscript{\ensuremath{q}}\fi}
\newunicodechar{ₙ}{\ifmmode_{n}\else\textsubscript{\ensuremath{n}}\fi}
\newunicodechar{ₐ}{\ifmmode_{a}\else\textsubscript{\ensuremath{a}}\fi}
\newunicodechar{ᵢ}{\ifmmode_{i}\else\textsubscript{\ensuremath{i}}\fi}
\newunicodechar{ᵣ}{\ifmmode_{r}\else\textsubscript{\ensuremath{r}}\fi}
\newunicodechar{ₖ}{\ifmmode_{k}\else\textsubscript{\ensuremath{k}}\fi}
\newunicodechar{ᵦ}{\ifmmode_{\beta}\else\textsubscript{\ensuremath{\beta}}\fi}
\newunicodechar{ₓ}{\ifmmode_{x}\else\textsubscript{\ensuremath{x}}\fi}
\newfontfamily\popdisplay{ArchivoBlack-Regular.ttf}[Path=../../../fonts/]
\newfontfamily\poplogo{SpaceGrotesk-Bold.ttf}[Path=../../../fonts/]
\newfontfamily\popsemi{PlusJakartaSans-SemiBold.ttf}[Path=../../../fonts/]
\newfontfamily\popxbold{PlusJakartaSans-ExtraBold.ttf}[Path=../../../fonts/]
\newfontfamily\popxboldls{PlusJakartaSans-ExtraBold.ttf}[Path=../../../fonts/,LetterSpace=3.0]

\setlist[enumerate]{leftmargin=1.7em,itemsep=5pt,topsep=4pt}
\setlist[itemize]{leftmargin=1.7em,itemsep=4pt,topsep=4pt,label={\color{poporange}\textbullet}}
\setlength{\parindent}{0pt}
\setlength{\parskip}{5pt}
\renewcommand{\arraystretch}{1.25}

% pgfplots : style Pop par défaut
\pgfplotsset{
  every axis/.append style={axis line style={popink,line width=1pt},tick style={popink},
    grid style={popline,line width=0.5pt},label style={font=\small},tick label style={font=\footnotesize}},
  popcurve/.style={poporange,line width=1.8pt,no markers,smooth},
}
% Même style disponible dans un \draw TikZ simple (hors pgfplots)
\tikzset{popcurve/.style={poporange,line width=1.8pt}}
\tikzset{popgrid/.style={popline,very thin}}
\colorlet{popcurve}{poporange} % le nom du style est parfois utilisé comme couleur

% ============================================================
% Pill / squiggle
% ============================================================
\newcommand{\poppill}[3]{%
  \tikz[baseline=-0.55ex]\node[draw=popink,line width=1.2pt,rounded corners=2.6pt,%
    fill=#2,inner xsep=6.5pt,inner ysep=3pt]{%
    \popxboldls\fontsize{7}{7}\selectfont\color{#3}\MakeUppercase{#1}};%
}
\newcommand{\popsquiggle}[1]{%
  \tikz[baseline=-0.4ex]{
    \draw[#1,line width=2.6pt,line cap=round]
      plot [smooth] coordinates {(0,0.09) (0.34,0.01) (0.68,0.21) (1.02,0.01) (1.36,0.10)};
  }%
}
% Mot-clé surligné (marqueur orange sous le texte)
\newcommand{\cle}[1]{{\popxbold\color{popdark}#1}}

% ============================================================
% En-tête / pied
% ============================================================
\pagestyle{fancy}
\fancyhf{}
\renewcommand{\headrulewidth}{0pt}
\renewcommand{\footrulewidth}{0pt}
\lhead{\raisebox{-0.15ex}{\poplogo\fontsize{13}{13}\selectfont\color{popink}OnBuch\color{poporange}+}}
\rhead{\poppill{\DOCMATIERE\ \textbullet\ \DOCNIVEAU\ \textbullet\ Leçon \DOCLECON}{white}{popink}}
\lfoot{\popxbold\fontsize{7.6}{7.6}\selectfont\color{popmuted}\MakeUppercase{Cours signé OnBuch+}\ \textbullet\ {\popsemi\DOCTITRE}}
\rfoot{\tikz[baseline=-0.6ex]\node[rounded corners=8.5pt,fill=popink,minimum height=17pt,minimum width=17pt,inner xsep=5pt,inner sep=0]{\popxbold\fontsize{8}{8}\selectfont\color{popcream}\thepage\,/\,\pageref*{LastPage}};}

% ============================================================
% Titres
% ============================================================
\newcommand{\chaptitre}[2]{%
  \begin{tcolorbox}[enhanced,colback=poporange,colframe=popink,boxrule=2.6pt,arc=11pt,
      drop shadow=popink,left=15pt,right=15pt,top=12pt,bottom=11pt,before skip=3pt,after skip=18pt]
    \poppill{#2}{popink}{popcream}\par\vspace{9pt}
    {\raggedright\hyphenpenalty=10000\exhyphenpenalty=10000\popdisplay\fontsize{22}{24}\selectfont\color{popcream}\MakeUppercase{#1}\par}
  \end{tcolorbox}
}
% Section numérotée : pastille orange avec le numéro + titre Archivo
\newcounter{popsec}
\newcommand{\coursec}[1]{%
  \stepcounter{popsec}\setcounter{popsub}{0}%
  \par\vspace{16pt}\noindent
  \begin{minipage}{\linewidth}
  \tikz[baseline=-0.7ex]\node[draw=popink,line width=1.6pt,fill=poporange,rounded corners=4pt,minimum size=22pt,inner sep=0]{\popdisplay\fontsize{12}{12}\selectfont\color{popcream}\thepopsec};%
  \hspace{8pt}{\popdisplay\fontsize{15}{17}\selectfont\color{popink}\MakeUppercase{#1}}\par
  \vspace{4pt}\noindent{\color{poporange}\rule{36pt}{3.6pt}}
  \end{minipage}\par\nopagebreak\vspace{8pt}
}
\newcounter{popsub}
\newcommand{\courssub}[1]{%
  \stepcounter{popsub}%
  \par\vspace{9pt}\noindent{\popxbold\fontsize{11.5}{13}\selectfont\color{popink}{\color{poporange}\thepopsec.\thepopsub}\hspace{6pt}#1}\par\nopagebreak\vspace{4pt}
}

% ============================================================
% Boîtes pédagogiques — même recette : fond blanc, bordure épaisse colorée,
% ombre dure encre, badge plein. Titre optionnel [#1].
% ============================================================
\tcbset{popbase/.style={breakable,enhanced,colback=white,boxrule=2.3pt,arc=9pt,
  shadow={3pt}{-3pt}{0pt}{popink},left=14pt,right=14pt,top=11pt,bottom=11pt,before skip=11pt,after skip=15pt,
  fontupper=\popsemi\fontsize{10.6}{14.5}\selectfont\color{popink},
  fontlower=\popsemi\fontsize{10.6}{14.5}\selectfont\color{popink}}}
% Coche dessinée (le glyphe ✓ n'existe pas dans les polices du document)
\newcommand{\popcheck}[1]{\tikz[baseline=-0.5ex]\draw[#1,line width=1.6pt,line cap=round,line join=round](-0.13,0.0)--(-0.04,-0.09)--(0.13,0.1);}
\providecommand{\coche}{\popcheck{popgreen}}
\providecommand{\resultat}[1]{\par\smallskip\noindent\fcolorbox{popgreen}{popgreenL}{\parbox{\dimexpr\linewidth-2\fboxsep-2\fboxrule\relax}{\textbf{Résultat.} #1}}\par\smallskip}
\newcommand{\popboxhead}[3]{\poppill{#1}{#2}{white}\ifx\relax#3\relax\else\hspace{6pt}{\popxbold\fontsize{10.6}{12}\selectfont\color{popink}#3}\fi\par\vspace{7pt}}

\newtcolorbox{aretenir}[1][]{popbase,colframe=popgreen,before upper={\popboxhead{À retenir}{popgreen}{#1}}}
\newtcolorbox{exemplebox}[1][]{popbase,colframe=poporange,before upper={\popboxhead{Exemple}{poporange}{#1}}}
\newtcolorbox{definition}[1][]{popbase,colframe=popblue,before upper={\popboxhead{Définition}{popblue}{#1}}}
\newtcolorbox{propriete}[1][]{popbase,colframe=poppurple,before upper={\popboxhead{Propriété}{poppurple}{#1}}}
\newtcolorbox{methode}[1][]{popbase,colframe=popgold,before upper={\popboxhead{Méthode}{popgold}{#1}}}
\newtcolorbox{attention}[1][]{popbase,colframe=poppink,before upper={\popboxhead{Attention}{poppink}{#1}}}
\newtcolorbox{experience}[1][]{popbase,colframe=popblue,colback=popblueL,before upper={\popboxhead{Expérience}{popblue}{#1}}}
% « Le savais-tu ? » : carte violette pleine (contexte, culture, Cameroun)
\newtcolorbox{savaistu}[1][]{popbase,colback=poppurple,colframe=popink,
  fontupper=\popsemi\fontsize{10.4}{14.2}\selectfont\color{white},
  before upper={\poppill{Le savais-tu\,?}{white}{poppurple}\ifx\relax#1\relax\else\hspace{6pt}{\popxbold\color{white}#1}\fi\par\vspace{7pt}}}
% Exercice résolu : énoncé en haut, \tcblower puis la solution sur fond crème
\newcounter{popexo}\newcounter{popexr}
\newtcolorbox{exoresolu}[1][]{popbase,colframe=poporange,colbacklower=popcream,
  segmentation style={popink,line width=1.2pt,dashed},
  before upper={\stepcounter{popexr}\popboxhead{Exercice résolu \thepopexr}{poporange}{#1}},
  before lower={\poppill{Solution}{popink}{popcream}\par\vspace{6pt}}}
% Exercice (non résolu en place) — la correction est dans la partie Corrigés
\newtcolorbox{exercice}[1][]{popbase,colframe=popink,
  before upper={\stepcounter{popexo}\popboxhead{Exercice \thepopexo}{popink}{#1}}}
% Corrigé
\newtcolorbox{corrige}[1][]{popbase,colframe=popgreen,colback=popgreenL,
  before upper={\popboxhead{Corrigé}{popgreen}{#1}}}
% Objectifs / prérequis / bilan
\newtcolorbox{objectifs}{popbase,colframe=popink,colback=poporangeL,
  before upper={\popboxhead{Objectifs de la leçon}{popink}{}}}
\newtcolorbox{prerequis}{popbase,colframe=popink,colback=popblueL,
  before upper={\popboxhead{Prérequis}{popblue}{}}}
\newtcolorbox{bilan}{popbase,colframe=popink,colback=white,boxrule=2.8pt,
  before upper={\popboxhead{Fiche bilan}{poporange}{L'essentiel à connaître pour l'examen}}}
% Figure encadrée avec légende
\newtcolorbox{popfigure}[1][]{enhanced,breakable=false,colback=white,colframe=popline,boxrule=1.4pt,arc=8pt,
  left=10pt,right=10pt,top=10pt,bottom=8pt,before skip=10pt,after skip=12pt,halign=center,
  lower separated=false,halign lower=center,
  fontlower=\popsemi\fontsize{9}{11.5}\selectfont\color{popmuted},
  #1}
\newcommand{\legende}[1]{\par\vspace{6pt}{\popsemi\fontsize{9}{11.5}\selectfont\color{popmuted}#1}}

% ============================================================
% Signature OnBuch+
% ============================================================
\newcommand{\poplogoinline}[1]{{\poplogo\fontsize{#1}{#1}\selectfont\color{popink}OnBuch\color{poporange}+}}
\newcommand{\signatureonbuch}{%
  \begin{tcolorbox}[enhanced,colback=popink,colframe=popink,boxrule=2.2pt,arc=10pt,
      drop shadow=poporange,left=14pt,right=14pt,top=10pt,bottom=10pt,before skip=0pt,after skip=0pt]
    \tikz[baseline=-0.7ex]\node[circle,fill=popgreen,draw=popcream,line width=1.3pt,minimum size=22pt,inner sep=0]{\popcheck{white}};%
    \hspace{9pt}\begin{minipage}[c]{0.62\linewidth}
      {\popxbold\fontsize{12}{13}\selectfont\color{popcream}Signé\ \textbullet\ {\poplogo OnBuch\color{poporange}+}}\par\vspace{2pt}
      {\raggedright\popsemi\fontsize{8}{10.5}\selectfont\color{popline}\MakeUppercase{Équipe pédagogique OnBuch+}\\\MakeUppercase{Conforme au programme officiel MINESEC}\par}
    \end{minipage}\hfill
    {\poplogo\fontsize{22}{22}\selectfont\color{popcream}OnBuch\color{poporange}+}
  \end{tcolorbox}%
}

% ============================================================
% Page de garde
% #1 titre · #2 matière · #3 niveau · #4 module · #5 n° leçon · #6 durée ·
% #7 description / objectifs (texte ou itemize)
% ============================================================
\newcommand{\pagegarde}[7]{%
  \thispagestyle{empty}
  \newgeometry{a4paper,margin=1.7cm,top=1.6cm,bottom=1.4cm}%
  \begin{tikzpicture}[remember picture,overlay]
    \fill[poporange!18] ([xshift=-3.2cm,yshift=-3.6cm]current page.north east) circle (4.6cm);
    \fill[poppurple!14] ([xshift=2.4cm,yshift=7.6cm]current page.south west) circle (3.8cm);
  \end{tikzpicture}%
  {\poplogo\fontsize{30}{30}\selectfont\color{popink}OnBuch\color{poporange}+}\hfill
  \raisebox{0.6ex}{\poppill{Cours signé \textbullet\ Premium}{popink}{popcream}}\par
  \vspace{1pt}\popsquiggle{poppurple}\par
  \vspace{8pt}
  \begin{tcolorbox}[enhanced,colback=poporange,colframe=popink,boxrule=2.8pt,arc=13pt,
      drop shadow=popink,left=18pt,right=18pt,top=12pt,bottom=13pt,after skip=10pt]
    \poppill{#2 \textbullet\ #3}{popink}{popcream}\hspace{5pt}\poppill{#4}{white}{popink}\par\vspace{8pt}
    {\popdisplay\fontsize{14}{14}\selectfont\color{popink}LEÇON #5}\par\vspace{4pt}
    {\raggedright\hyphenpenalty=10000\exhyphenpenalty=10000\popdisplay\fontsize{25}{27}\selectfont\color{popcream}\MakeUppercase{#1}\par}
    \vspace{8pt}
    {\popxbold\fontsize{9.5}{9.5}\selectfont\color{popink}\MakeUppercase{Durée conseillée : #6}}
  \end{tcolorbox}
  \begin{tcolorbox}[enhanced,colback=white,colframe=popink,boxrule=2.2pt,arc=10pt,
      drop shadow=popink,left=16pt,right=16pt,top=10pt,bottom=10pt,after skip=8pt]
    \poppill{Dans cette leçon}{poppurple}{white}\par\vspace{5pt}
    \popsemi\fontsize{9.3}{12.6}\selectfont\color{popink}#7
  \end{tcolorbox}
  \noindent\begin{minipage}[t]{0.487\linewidth}
    \begin{tcolorbox}[enhanced,colback=popgreen,colframe=popink,boxrule=2.2pt,arc=10pt,
        drop shadow=popink,left=11pt,right=11pt,top=10pt,bottom=10pt,height=3.1cm,valign=center]
      \tcbox[colback=white,colframe=popink,boxrule=1.3pt,arc=5pt,boxsep=0pt,nobeforeafter,
        left=3pt,right=3pt,top=3pt,bottom=3pt,valign=center]{\qrcode[height=1.9cm]{https://play.google.com/store/apps/details?id=com.luvvix.onbuch&hl=fr}}%
      \hspace{8pt}\begin{minipage}[c]{0.5\linewidth}
        {\popdisplay\fontsize{11}{12.5}\selectfont\color{white}\MakeUppercase{Télécharge OnBuch}}\par\vspace{4pt}
        {\raggedright\popsemi\fontsize{8.4}{11}\selectfont\color{white}Gratuit \textbullet\ Google Play\\Tuteur IA, annales, concours, résultats\par}
      \end{minipage}
    \end{tcolorbox}
  \end{minipage}\hfill
  \begin{minipage}[t]{0.487\linewidth}
    \begin{tcolorbox}[enhanced,colback=poppurple,colframe=popink,boxrule=2.2pt,arc=10pt,
        drop shadow=popink,left=11pt,right=11pt,top=10pt,bottom=10pt,height=3.1cm,valign=center]
      \tikz[baseline=-0.7ex]\node[circle,fill=white,draw=popink,line width=1.3pt,minimum size=1.6cm,inner sep=0]{\popdisplay\fontsize{24}{24}\selectfont\color{poppurple}+};%
      \hspace{8pt}\begin{minipage}[c]{0.6\linewidth}
        {\popdisplay\fontsize{11}{12.5}\selectfont\color{white}\MakeUppercase{Passe à OnBuch+}}\par\vspace{4pt}
        {\raggedright\popsemi\fontsize{8.4}{11}\selectfont\color{white}Tous les cours signés, Tuteur IA illimité, fiches PDF — onglet OnBuch+ de l'app\par}
      \end{minipage}
    \end{tcolorbox}
  \end{minipage}
  \vspace{8pt}
  \popcontenu
  \vfill
  \signatureonbuch
  \restoregeometry
  \newpage
}

% Rangée « Ce cours contient » (page de garde)
\newcommand{\poptile}[3]{% #1 couleur #2 gros texte #3 légende
  \begin{tcolorbox}[enhanced,colback=#1,colframe=popink,boxrule=2pt,arc=9pt,shadow={2.5pt}{-2.5pt}{0pt}{popink},
      left=6pt,right=6pt,top=9pt,bottom=9pt,halign=center,valign=center,height=1.75cm,width=\linewidth,nobeforeafter]
    {\popdisplay\fontsize{12.5}{13}\selectfont\color{white}\MakeUppercase{#2}}\par\vspace{3pt}
    {\centering\popsemi\fontsize{7.8}{9.5}\selectfont\color{white}#3\par}
  \end{tcolorbox}}
\newcommand{\popcontenu}{%
  \par\noindent{\popxboldls\fontsize{7.5}{7.5}\selectfont\color{popmuted}CE COURS SIGNÉ CONTIENT}\par\vspace{7pt}
  \noindent\begin{minipage}[t]{0.235\linewidth}\poptile{popblue}{Cours}{complet, illustré, pas à pas}\end{minipage}\hfill
  \begin{minipage}[t]{0.235\linewidth}\poptile{poporange}{Méthodes}{et exercices résolus}\end{minipage}\hfill
  \begin{minipage}[t]{0.235\linewidth}\poptile{poppink}{Exercices}{type examen + corrigés}\end{minipage}\hfill
  \begin{minipage}[t]{0.235\linewidth}\poptile{popgreen}{Fiche bilan}{l'essentiel à retenir}\end{minipage}\par}

% Sommaire
\newtcolorbox{sommaire}{enhanced,colback=white,colframe=popink,boxrule=2.2pt,arc=10pt,
  drop shadow=popink,left=18pt,right=18pt,top=13pt,bottom=13pt,before skip=6pt,after skip=20pt,
  before upper={\poppill{Sommaire}{poppurple}{white}\par\vspace{9pt}\popsemi\fontsize{10.6}{14.5}\selectfont\color{popink}}}

% Signature de fin de cours — poussée tout en bas de la dernière page
\newcommand{\findecours}{%
  \par\vfill\nopagebreak
  \begin{center}\popsquiggle{poporange}\end{center}\vspace{4pt}
  \signatureonbuch
}
\AtBeginDocument{\def\llbracket{\mathopen{[\mkern-3.5mu[}}\def\rrbracket{\mathclose{]\mkern-3.5mu]}}} % intervalles d'entiers (glyphe ⟦ absent de la police)
% Glyphes absents de Fira Math : redéfinis à partir de glyphes disponibles
\AtBeginDocument{%
  \def\setminus{\mathbin{\backslash}}\let\smallsetminus\setminus
  \def\vdots{\mathord{\vbox{\baselineskip3.2pt\lineskiplimit0pt\kern4pt\hbox{.}\hbox{.}\hbox{.}}}}%
  \def\ddots{\mathinner{\mkern1mu\raise6.5pt\hbox{.}\mkern2mu\raise3.6pt\hbox{.}\mkern2mu\raise0.7pt\hbox{.}\mkern1mu}}%
  \def\triangle{\mathord{\scalebox{0.9}{$\symup{\Delta}$}}}%
  \def\nabla{\mathord{\raisebox{\depth}{\scalebox{1}[-1]{$\symup{\Delta}$}}}}%
  \def\checkmark{\popcheck{popdark}}%
  \def\star{\mathbin{\ast}}\def\bigstar{\mathord{\ast}}%
  \def\diamond{\mathbin{\scalebox{0.6}{$\Diamond$}}}%
  \def\bigcup{\mathop{\mathchoice{\raisebox{-0.3em}{\scalebox{1.7}{$\cup$}}}{\scalebox{1.25}{$\cup$}}{\cup}{\cup}}\displaylimits}%
  \def\bigcap{\mathop{\mathchoice{\raisebox{-0.3em}{\scalebox{1.7}{$\cap$}}}{\scalebox{1.25}{$\cap$}}{\cap}{\cap}}\displaylimits}%
  \def\lesseqqgtr{\mathrel{\lessgtr}}\def\bigoplus{\mathop{\oplus}\displaylimits}%
}
\providecommand{\ssi}{si et seulement si} % abréviation fréquente
\AtBeginDocument{\providecommand{\wideparen}{\overparen}} % arc de cercle

%% FIGFIX-BEGIN v5 — correctifs globaux des figures (docs/onbuch-generation/tools/figfix)
\usepackage{adjustbox}\usepackage{etoolbox}\tcbuselibrary{raster}
% toute figure TikZ/pgfplots plus large que la ligne est réduite à la largeur disponible
\BeforeBeginEnvironment{tikzpicture}{\begin{adjustbox}{max width=\linewidth}}
\AfterEndEnvironment{tikzpicture}{\end{adjustbox}}
% légendes pgfplots lisibles par-dessus les courbes
\pgfplotsset{every axis/.append style={legend style={fill=white,fill opacity=0.92,text opacity=1,draw=black!15,inner sep=3pt}}}
% étiquettes d'arêtes (syntaxe edge["..."]) avec fond blanc
\tikzset{every edge quotes/.append style={fill=white,inner sep=1.2pt}}
%% FIGFIX-END
\begin{document}

\pagegarde{Modes de transfert de chaleur et calorimétrie}{Physique}{1ère E}{Module 2 : Mouvements et interactions mécaniques}{N06}{8 h}{Cette leçon étudie les trois modes de transfert de chaleur (conduction, convection, rayonnement) et introduit la calorimétrie : mesure et calcul des quantités de chaleur échangées entre corps, avec ou sans changement d'état. Elle fournit les outils indispensables pour résoudre les problèmes de type Probatoire sur les mélanges, les calorimètres et les changements d'état.
\vspace{4pt}
\begin{multicols}{2}\small
\begin{itemize}
\item À la fin de cette leçon, je sais citer des situations du quotidien où l'énergie se manifeste sous forme thermique
\item identifier les phénomènes mettant en jeu des échanges thermiques
\item distinguer et nommer le mode de transfert de chaleur (conduction, convection, rayonnement) dans une situation donnée
\item énoncer et appliquer le principe des échanges de chaleur dans un système isolé
\item définir la chaleur massique, la capacité calorifique, la valeur en eau d'un calorimètre et la chaleur latente de changement d'état
\item réaliser le schéma annoté d'un calorimètre avec ses accessoires et expliquer sa fonction
\end{itemize}\end{multicols}}

\begin{sommaire}
\begin{multicols}{2}
\begin{enumerate}
\item Chaleur, température et échanges thermiques
\item Les trois modes de transfert de chaleur
\item Quantité de chaleur sans changement d'état
\item Chaleur latente et changements d'état
\item Le calorimètre et le principe des échanges de chaleur
\item Synthèse et résolution de problèmes de calorimétrie
\item Activité d'intégration
\item Méthodes et exercices résolus
\item Exercices
\item Corrigés des exercices
\item Fiche bilan
\end{enumerate}
\end{multicols}
\end{sommaire}

\begin{objectifs}
\begin{itemize}
\item À la fin de cette leçon, je sais citer des situations du quotidien où l'énergie se manifeste sous forme thermique
\item identifier les phénomènes mettant en jeu des échanges thermiques
\item distinguer et nommer le mode de transfert de chaleur (conduction, convection, rayonnement) dans une situation donnée
\item énoncer et appliquer le principe des échanges de chaleur dans un système isolé
\item définir la chaleur massique, la capacité calorifique, la valeur en eau d'un calorimètre et la chaleur latente de changement d'état
\item réaliser le schéma annoté d'un calorimètre avec ses accessoires et expliquer sa fonction
\item calculer la quantité de chaleur échangée par un corps sans changement d'état (Q = m·c·Δθ)
\item calculer la quantité de chaleur échangée lors d'un changement d'état (Q = m·L)
\item reconnaître dans l'environnement des appareils analogues au calorimètre (glacière, thermos, marmite norvégienne)
\end{itemize}
\end{objectifs}

\begin{prerequis}
\begin{itemize}
\item Notion de température et échelle Celsius (classes de collège)
\item États de la matière et changements d'état (fusion, vaporisation, solidification, liquéfaction)
\item Notion d'énergie et unité joule (début d'année de Première)
\item Système isolé et conservation de l'énergie (notions vues en mécanique)
\item Calculs algébriques simples : isoler une grandeur dans une formule
\end{itemize}
\end{prerequis}

\newpage

\coursec{Chaleur, température et échanges thermiques}

\textbf{Situation d'accroche.} Il est 11 h à Ebolowa. Mama Ngo Bell prépare le ndolé dans une grande marmite en aluminium posée sur un feu de bois. La flamme ne touche que le fond de la marmite, pourtant, quelques minutes plus tard, les anses sont si brûlantes qu'il faut un torchon pour les saisir. Au même moment, à Douala, un élève qui marche sous le soleil sent la chaleur sur sa peau \emph{sans aucun contact} avec le Soleil, situé à 150 millions de kilomètres. Et dans la casserole, on observe que l'eau chaude monte pendant que l'eau froide descend, créant un lent tourbillon. Enfin, à Yaoundé, un vendeur de glaces garde ses produits froids toute la journée dans une simple glacière, alors que la température extérieure dépasse \SI{30}{\celsius}.

\textbf{Problématique.} Comment la chaleur voyage-t-elle du feu jusqu'à nos mains ? Qu'est-ce que la chaleur, exactement, et en quoi diffère-t-elle de la température ? Comment calculer la quantité de chaleur échangée lorsqu'on mélange de l'eau chaude et de l'eau froide, ou lorsque la glace fond dans un jus de foléré ? Cette leçon répond à toutes ces questions. Dans cette première section, nous posons les fondations : la distinction rigoureuse entre \cle{température} et \cle{chaleur}, le sens des échanges thermiques et la notion d'\cle{équilibre thermique}.

\courssub{Température et chaleur : deux grandeurs à ne pas confondre}

\courssub{La température : une grandeur d'état}

Plonge ta main dans un seau d'eau de puits un matin à Bafoussam : tu dis « l'eau est froide ». Plonge-la dans l'eau de rinçage du ndolé : tu dis « elle est chaude ». Ces sensations sont liées à une grandeur physique appelée \cle{température}.

\begin{definition}[Température]
La \cle{température} d'un corps est une grandeur physique qui caractérise son \textbf{état thermique}, c'est-à-dire son degré d'échauffement. Elle se mesure à l'aide d'un \textbf{thermomètre}.
\begin{itemize}
\item Dans la vie courante et au laboratoire, on utilise le \textbf{degré Celsius} (\si{\celsius}) : l'eau gèle à \SI{0}{\celsius} et bout à \SI{100}{\celsius} sous la pression atmosphérique normale.
\item Dans le Système international (SI), l'unité de température est le \textbf{kelvin} (\si{K}). La relation entre les deux échelles est :
\[ T(\si{K}) = \theta(\si{\celsius}) + 273{,}15 \]
\end{itemize}
\end{definition}

Microscopiquement, la température traduit l'\textbf{agitation des particules} (atomes, molécules) qui constituent le corps : plus les molécules s'agitent, plus le corps est chaud. Le zéro absolu, $T = \SI{0}{K}$ soit $\theta = \SI{-273,15}{\celsius}$, correspond à l'état où cette agitation est minimale : on ne peut pas descendre en dessous.

\begin{attention}[Une différence de température a la même valeur dans les deux échelles]
Un écart de température s'exprime indifféremment en \si{\celsius} ou en \si{K} :
\[ \Delta T = T_2 - T_1 = \theta_2 - \theta_1 \]
Par exemple, passer de \SI{20}{\celsius} à \SI{80}{\celsius} correspond à une élévation de température de \SI{60}{\celsius}, c'est-à-dire exactement \SI{60}{K}. En revanche, une \emph{valeur} de température doit être convertie : $\SI{20}{\celsius} = \SI{293,15}{K}$ et non \SI{20}{K} !
\end{attention}

\courssub{La chaleur : une énergie en transfert}

Reprenons la marmite de Mama Ngo Bell. Le fond de la marmite, en contact avec les braises, voit sa température s'élever ; puis la partie haute s'échauffe à son tour. Quelque chose « circule » du feu vers la marmite, puis à l'intérieur du métal. Ce quelque chose n'est pas de la température — la température ne se déplace pas — c'est de l'\textbf{énergie}.

\begin{definition}[Chaleur]
La \cle{chaleur}, notée $Q$, est l'\textbf{énergie transférée} entre deux corps (ou deux parties d'un même corps) qui sont à des \textbf{températures différentes}. Ce transfert est \textbf{spontané} et s'effectue \textbf{toujours du corps le plus chaud vers le corps le plus froid}, jamais l'inverse.
\begin{itemize}
\item L'unité SI de la chaleur est le \textbf{joule} (\si{J}), comme toutes les formes d'énergie.
\item On utilise encore parfois la \textbf{calorie} (\si{cal}) : c'est la quantité de chaleur nécessaire pour élever la température de \SI{1}{g} d'eau de \SI{1}{\celsius} (plus précisément de \num{14,5} à \SI{15,5}{\celsius}). La conversion à connaître pour le Probatoire est :
\[ \boxed{\ \SI{1}{cal} = \SI{4,18}{J} \qquad \text{et} \qquad \SI{1}{kcal} = \SI{4180}{J}\ } \]
\end{itemize}
\end{definition}

\begin{aretenir}[Le point essentiel de toute la leçon]
\begin{itemize}
\item La \textbf{température} caractérise l'\emph{état} d'un corps à un instant donné ; elle se mesure en \si{\celsius} ou en \si{K}.
\item La \textbf{chaleur} est une \emph{énergie en transit} ; elle n'existe que pendant un échange entre un corps chaud et un corps froid ; elle se mesure en joules (\si{J}).
\item On ne dit jamais « un corps possède de la chaleur ». Un corps possède de l'\textbf{énergie interne} (liée à l'agitation de ses particules) ; la chaleur est le nom de l'énergie \emph{lorsqu'elle passe} d'un corps à l'autre sous l'effet d'une différence de température.
\end{itemize}
\end{aretenir}

\begin{attention}[Erreur fréquente n°1 : confondre chaleur et température]
Voici le piège classique du Probatoire. Compare une \textbf{tasse d'eau bouillante} (\SI{95}{\celsius}) et le \textbf{lac de Barombi} à \SI{25}{\celsius}. La tasse est à plus haute température, certes. Mais le lac contient des milliards de fois plus de molécules d'eau : son énergie thermique totale est gigantesque comparée à celle de la tasse. La température mesure l'agitation \emph{moyenne} des particules ; l'énergie thermique dépend en plus de la \emph{quantité de matière}. Conclusion : un corps froid peut contenir beaucoup plus d'énergie thermique qu'un corps chaud. Retenir : température élevée $\neq$ grande quantité de chaleur disponible.
\end{attention}

\begin{exemplebox}[Conversion calorie $\leftrightarrow$ joule]
Une étiquette de beignet-haricot indique un apport énergétique de \SI{250}{kcal}. Convertissons en joules :
\[ Q = 250 \times \SI{4,18}{kJ} = \SI{1045}{kJ} = \SI{1,045 \times 10^6}{J} \]
De même, \SI{250}{cal} correspondent à $250 \times 4{,}18 = \SI{1045}{J}$. Attention à ne pas confondre \si{cal} et \si{kcal} : les « calories » des nutritionnistes sont en réalité des kilocalories !
\end{exemplebox}

\courssub{Le sens des échanges thermiques et l'équilibre thermique}

\courssub{Observation : la chaleur ne « remonte » jamais spontanément}

Verse de l'eau chaude de la marmite dans une calebasse d'eau froide : le mélange devient tiède. Jamais, au grand jamais, tu n'observeras l'inverse : un mélange tiède qui se séparerait spontanément en une partie chaude et une partie froide. C'est un fait d'expérience fondamental.

\begin{propriete}[Sens spontané du transfert de chaleur]
Lorsque deux corps de températures différentes $T_1 > T_2$ sont mis en contact (ou peuvent échanger de l'énergie), la chaleur s'écoule \textbf{spontanément et irréversiblement du corps chaud vers le corps froid}. Le corps chaud \textbf{cède} de la chaleur ($Q_{\text{cédée}} < 0$ pour lui), le corps froid \textbf{reçoit} de la chaleur ($Q_{\text{reçue}} > 0$ pour lui).
\end{propriete}

\courssub{L'équilibre thermique}

Que se passe-t-il au bout d'un certain temps ? Le corps chaud refroidit, le corps froid se réchauffe, et les deux températures se rapprochent jusqu'à devenir \textbf{égales}. On dit alors que les deux corps sont en \cle{équilibre thermique}.

\begin{definition}[Équilibre thermique]
Deux corps sont en \cle{équilibre thermique} lorsqu'ils sont à la \textbf{même température}. À l'équilibre thermique, les échanges de chaleur macroscopiques cessent : il n'y a plus de transfert net d'énergie entre les deux corps.
\end{definition}

C'est d'ailleurs le principe du thermomètre : mis en contact avec un corps, il échange de la chaleur avec lui jusqu'à l'équilibre thermique ; il indique alors sa propre température, qui est devenue celle du corps.

\begin{popfigure}
\begin{center}
\begin{tikzpicture}[scale=0.9, font=\small]
% Bloc chaud
\fill[poporange!70] (0,0) rectangle (2.2,1.6);
\draw[popdark, thick] (0,0) rectangle (2.2,1.6);
\node at (1.1,0.8) {\textbf{Corps chaud}};
\node at (1.1,-0.4) {$T_1 = \SI{80}{\celsius}$};
% Bloc froid
\fill[popblue!40] (2.2,0) rectangle (4.4,1.6);
\draw[popdark, thick] (2.2,0) rectangle (4.4,1.6);
\node at (3.3,0.8) {\textbf{Corps froid}};
\node at (3.3,-0.4) {$T_2 = \SI{20}{\celsius}$};
% Flèche de chaleur
\draw[-{Stealth[length=3mm]}, poppink, line width=2pt] (1.5,2.1) -- (2.9,2.1) node[midway, above, poppink] {$Q$};
\node[popdark] at (2.2,-1.1) {\emph{Contact thermique : $Q$ s'écoule du chaud vers le froid}};
% Évolution vers l'équilibre
\draw[-{Stealth[length=3mm]}, popmuted, thick] (5.0,0.8) -- (6.0,0.8);
% État final
\fill[popgold!45] (6.4,0) rectangle (10.8,1.6);
\draw[popdark, thick] (6.4,0) rectangle (10.8,1.6);
\draw[popdark, dashed] (8.6,0) -- (8.6,1.6);
\node at (8.6,0.8) {\textbf{Équilibre thermique}};
\node at (8.6,-0.4) {$T_f$ identique pour les deux corps};
\node[popdark] at (8.6,-1.1) {\emph{Les échanges cessent : $20 < T_f < \SI{80}{\celsius}$}};
\end{tikzpicture}
\end{center}
\legende{Deux blocs de températures différentes mis en contact : la chaleur $Q$ passe du corps chaud vers le corps froid jusqu'à l'équilibre thermique, où les deux corps atteignent la même température finale $T_f$.}
\end{popfigure}

\begin{popfigure}
\begin{center}
\begin{tikzpicture}
\begin{axis}[
width=0.85\linewidth, height=5.2cm,
xlabel={Temps $t$ (min)}, ylabel={Température $\theta$ (\si{\celsius})},
xmin=0, xmax=10, ymin=0, ymax=90,
grid=both, grid style={popline!40},
legend style={at={(0.98,0.97)}, anchor=north east, font=\small},
]
\addplot[popcurve, poporange, domain=0:10, samples=100] {45 + 35*exp(-0.45*x)};
\addlegendentry{Corps chaud (se refroidit)}
\addplot[popcurve, popblue, domain=0:10, samples=100] {45 - 25*exp(-0.45*x)};
\addlegendentry{Corps froid (se réchauffe)}
\addplot[popmuted, dashed, domain=0:10] {45};
\addlegendentry{Température d'équilibre $T_f$}
\end{axis}
\end{tikzpicture}
\end{center}
\legende{Évolution des températures de deux corps en contact thermique : les courbes convergent vers la même température finale, l'équilibre thermique.}
\end{popfigure}

\courssub{La chaleur dans l'environnement : exemples d'énergie thermique}

L'énergie se manifeste sous forme thermique partout autour de nous. Sache citer ces exemples, c'est un savoir-faire exigible :

\begin{itemize}
\item \textbf{Le feu de bois} de la cuisine de Mama Ngo Bell : la combustion du bois libère de l'énergie chimique transformée en chaleur.
\item \textbf{Le Soleil} : il envoie chaque jour sur le Cameroun une énorme quantité de chaleur, source de toute vie et du séchage du manioc et du cacao.
\item \textbf{Le moteur d'une mototaxi} : après une course Douala--Akwa, le pot d'échappement est brûlant ; une partie de l'énergie du carburant est dégradée en chaleur.
\item \textbf{Le fer à repasser} : l'énergie électrique d'ENEO est convertie en chaleur dans la résistance du fer.
\item \textbf{Le corps humain} : nos cellules « brûlent » les aliments et maintiennent notre température à environ \SI{37}{\celsius}, dégageant en permanence de la chaleur vers l'extérieur.
\end{itemize}

\begin{savaistu}[Ton corps est un radiateur de 100 watts]
Un être humain au repos dégage une puissance thermique d'environ \SI{100}{W}, comme une vieille ampoule à incandescence ! En une journée, cela représente $Q = 100 \times 24 \times 3600 \approx \SI{8,6 \times 10^6}{J}$, soit environ \SI{2000}{kcal} : c'est précisément l'apport alimentaire quotidien recommandé. C'est pourquoi une salle de classe pleine se réchauffe vite, et pourquoi on a chaud dans un taxi-brousse bondé même sans soleil direct.
\end{savaistu}

\courssub{Identifier les phénomènes mettant en jeu des échanges thermiques}

Un savoir-faire du programme consiste à \textbf{reconnaître} les situations où de la chaleur est échangée, et à identifier \emph{qui cède} et \emph{qui reçoit}. Analysons trois situations quotidiennes.

\begin{exemplebox}[Trois situations à analyser]
\begin{enumerate}
\item \textbf{Chauffage de l'eau sur le feu.} Les braises et les gaz brûlants (corps chaud) cèdent de la chaleur à la marmite puis à l'eau (corps froids) : la température de l'eau s'élève jusqu'à l'ébullition.
\item \textbf{Refroidissement d'un plat de eru servi trop chaud.} Le plat (corps chaud) cède de la chaleur à l'air ambiant et à l'assiette (corps froids) ; sa température descend jusqu'à celle de la pièce : c'est l'équilibre thermique avec l'environnement.
\item \textbf{Glaçons dans un jus de foléré.} Le jus, plus chaud, cède de la chaleur aux glaçons ; ceux-ci l'utilisent pour \emph{fondre} (changement d'état, étudié en section 4) pendant que le jus se rafraîchit. C'est le principe du vendeur de glaces : la glacière ralentit les échanges avec l'extérieur chaud.
\end{enumerate}
\end{exemplebox}

\begin{methode}[Analyser une situation d'échange thermique en 4 étapes]
\begin{enumerate}
\item \textbf{Repérer les corps} en présence et leurs températures initiales.
\item \textbf{Identifier le corps chaud et le corps froid} : le transfert ira du chaud vers le froid.
\item \textbf{Écrire le sens du transfert} : « $Q$ passe de \ldots vers \ldots » ; le corps chaud cède, le corps froid reçoit.
\item \textbf{Prévoir l'état final} : si les corps restent en contact assez longtemps, ils atteignent une température commune $T_f$ (équilibre thermique), comprise entre les deux températures initiales.
\end{enumerate}
\end{methode}

\courssub{Unités de chaleur : joule et calorie, le pont de conversion}

Puisque la chaleur est une énergie, son unité SI est le joule, déjà rencontré en mécanique (travail d'une force). La calorie est une unité historique, encore très utilisée en calorimétrie et en nutrition ; le passage de l'une à l'autre est un réflexe à acquérir.

\begin{popfigure}
\begin{center}
\begin{tikzpicture}[font=\small]
\node[draw=popblue, fill=popblueL, rounded corners, thick, minimum width=3.4cm, minimum height=1.2cm] (cal) at (0,0) {\textbf{Calories (cal)}};
\node[draw=poporange, fill=poporangeL, rounded corners, thick, minimum width=3.4cm, minimum height=1.2cm] (J) at (8,0) {\textbf{Joules (J)}};
\draw[-{Stealth[length=3mm]}, popgreen, line width=1.6pt] (cal.east) -- ++(1.1,0.35) -- ++(2.4,0) -- (J.west |- 0,0.35);
\node[popgreen] at (4,0.85) {$\times 4{,}18$};
\draw[-{Stealth[length=3mm]}, poppurple, line width=1.6pt] (J.west) ++(0,-0.35) -- ++(-2.4,0) -- ++(-1.1,0) -- (cal.east |- 0,-0.35);
\node[poppurple] at (4,-0.85) {$\div 4{,}18$};
\node[popdark] at (4,-1.7) {Exemples : $\SI{250}{cal} = \SI{1045}{J}$ \quad ; \quad $\SI{836}{J} = \SI{200}{cal}$};
\end{tikzpicture}
\end{center}
\legende{Diagramme de conversion entre calories et joules : multiplier par \num{4,18} pour passer des calories aux joules, diviser par \num{4,18} pour le chemin inverse.}
\end{popfigure}

\begin{exoresolu}[Qui cède, qui reçoit, et combien de joules ?]
\textbf{Énoncé.} Pour préparer le bain de bébé, une maman à Bafoussam mélange de l'eau chaude à $\theta_1 = \SI{60}{\celsius}$ avec de l'eau froide à $\theta_2 = \SI{18}{\celsius}$. Le thermomètre du bain indique finalement $\theta_f = \SI{32}{\celsius}$.
\begin{enumerate}
\item Identifier le corps qui cède de la chaleur et celui qui en reçoit. Justifier.
\item L'énoncé d'un exercice indique que l'eau chaude a cédé \SI{1500}{cal}. Exprimer cette quantité de chaleur en joules.
\item Pourquoi la température finale est-elle comprise entre \SI{18}{\celsius} et \SI{60}{\celsius} ?
\end{enumerate}
\tcblower
\textbf{Solution détaillée.}
\begin{enumerate}
\item L'eau chaude est à la température la plus élevée ($\SI{60}{\celsius} > \SI{18}{\celsius}$). Or la chaleur s'écoule spontanément du corps chaud vers le corps froid. Donc \textbf{l'eau chaude cède de la chaleur} (elle se refroidit de 60 à \SI{32}{\celsius}) et \textbf{l'eau froide reçoit de la chaleur} (elle se réchauffe de 18 à \SI{32}{\celsius}).
\item Conversion : $Q = 1500 \times 4{,}18 = \SI{6270}{J} = \SI{6,27}{kJ}$.
\item Le transfert s'arrête à l'\textbf{équilibre thermique}, lorsque les deux corps atteignent la même température. Cette température finale est nécessairement \emph{intermédiaire} : l'eau chaude ne peut pas descendre en dessous de la température de l'eau froide (elle cesserait de lui céder de la chaleur), et réciproquement. D'où $\SI{18}{\celsius} < \theta_f < \SI{60}{\celsius}$, ce qui est bien vérifié : $\theta_f = \SI{32}{\celsius}$.
\end{enumerate}
\end{exoresolu}

\begin{attention}[Pièges de rédaction au Probatoire]
\begin{itemize}
\item Ne jamais écrire « la température passe du corps chaud au corps froid » : c'est la \textbf{chaleur} (une énergie) qui se transfère ; les températures, elles, \emph{évoluent}.
\item Ne pas oublier l'unité : une chaleur s'exprime en \si{J} (ou \si{kJ}, \si{cal}), jamais en \si{\celsius}.
\item Dans une conversion, vérifier le sens : les joules sont « plus petits » que les calories, donc la valeur en joules doit être \emph{plus grande} que la valeur en calories ($\times 4{,}18$).
\item Une température finale commune inférieure aux deux températures initiales (ou supérieure aux deux) est \textbf{physiquement impossible} : c'est le signe d'une erreur de calcul.
\end{itemize}
\end{attention}

\begin{aretenir}[Bilan de la section]
\begin{itemize}
\item La \cle{température} (en \si{\celsius} ou \si{K}) décrit l'état thermique d'un corps ; la \cle{chaleur} $Q$ (en \si{J}) est une énergie \emph{en transfert} entre corps de températures différentes.
\item Le transfert de chaleur est spontané et s'effectue \textbf{toujours du corps chaud vers le corps froid}, jusqu'à l'\cle{équilibre thermique} (températures égales).
\item Conversion à retenir : $\SI{1}{cal} = \SI{4,18}{J}$.
\item Identifier un échange thermique, c'est répondre à trois questions : qui est chaud ? qui est froid ? dans quel sens circule $Q$ ?
\end{itemize}
Dans la section suivante, nous répondrons à la question posée par la marmite de Mama Ngo Bell et le soleil de Douala : \emph{par quels mécanismes} la chaleur se déplace-t-elle ? Ce sont les trois modes de transfert : conduction, convection et rayonnement.
\end{aretenir}

\coursec{Les trois modes de transfert de chaleur}

Dans la section précédente, nous avons distingué la \cle{température} (grandeur d'état, liée à l'agitation microscopique) de la \cle{chaleur} (énergie en transit, qui passe spontanément du corps chaud vers le corps froid). Une question se pose alors naturellement : \emph{comment} cette énergie voyage-t-elle concrètement ? Quand tu tiens le manche métallique d'une marmite sur le feu, quand l'eau bout dans la casserole, quand tu sens la chaleur du soleil sur ta peau à Kribi, trois mécanismes physiques \textbf{différents} sont à l'œuvre. Les identifier, c'est la clé de toute la calorimétrie — et de bien des gestes de la vie quotidienne, du choix d'une casserole à la construction d'une case fraîche.

\courssub{La conduction thermique}

\textbf{Observation.} Plonge une cuillère métallique dans un pot d'eau chaude et une cuillère en bois dans un autre pot identique. Au bout d'une minute, le manche de la cuillère métallique est chaud, celui de la cuillère en bois est resté tiède. Pourtant, aucune goutte de métal ne s'est déplacée le long du manche : la chaleur a « voyagé » sans que la matière ne bouge.

\begin{definition}[Conduction thermique]
La \cle{conduction} est un mode de transfert de chaleur qui s'effectue \textbf{par contact direct}, \textbf{de proche en proche}, \textbf{sans déplacement macroscopique de matière}. Au niveau microscopique, les particules de la zone chaude, plus agitées, communiquent leur agitation aux particules voisines par chocs successifs (et, dans les métaux, grâce aux électrons libres très mobiles).
\end{definition}

La conduction est le mode dominant dans les \textbf{solides}. Elle se manifeste chaque fois qu'un gradient de température existe dans un matériau : barre de fer chauffée à une extrémité, tôle d'un toit exposée au soleil, sol d'une cour en plein après-midi.

\begin{propriete}[Conducteurs et isolants thermiques]
Les matériaux ne conduisent pas tous la chaleur de la même façon :
\begin{itemize}
\item les \textbf{bons conducteurs thermiques} sont les métaux : cuivre, aluminium, fer. Ils transmettent rapidement la chaleur ;
\item les \textbf{mauvais conducteurs} (ou \cle{isolants thermiques}) sont le bois, le plastique, le liège, la laine, la terre et surtout l'\textbf{air immobile} emprisonné dans les matériaux poreux.
\end{itemize}
\end{propriete}

C'est pourquoi les casseroles ont un corps en aluminium (bon conducteur : la chaleur du feu atteint vite les aliments) et un manche en bois ou en plastique (mauvais conducteur : il protège la main). Le choix des matériaux n'est jamais un hasard : c'est de la physique appliquée.

\courssub{La convection thermique}

\textbf{Observation.} Chauffe de l'eau dans une casserole transparente avec quelques grains de riz ou de semoule. On voit les grains \emph{monter} au centre et \emph{redescendre} le long des parois : l'eau elle-même se déplace en courants permanents. Ici, contrairement à la conduction, c'est la \textbf{matière chaude qui se déplace} et transporte son énergie avec elle.

\begin{definition}[Convection thermique]
La \cle{convection} est un mode de transfert de chaleur qui s'effectue \textbf{par déplacement de matière} au sein d'un \textbf{fluide} (liquide ou gaz). Le fluide chauffé se dilate, devient moins dense et s'élève sous l'effet de la poussée d'Archimède ; le fluide froid, plus dense, descend pour prendre sa place. Il s'établit des \textbf{courants de convection}.
\end{definition}

\begin{experience}[Courants de convection dans l'eau]
\textbf{Matériel :} un bécher, de l'eau, quelques cristaux de permanganate de potassium, une source de chaleur (bec électrique ou bougie).\par
\textbf{Protocole :} déposer les cristaux au fond du bécher, près d'un bord, et chauffer doucement juste en dessous.\par
\textbf{Observation :} le filet coloré monte au-dessus de la source chaude, traverse la surface, puis redescend de l'autre côté : on visualise une boucle de courant.\par
\textbf{Interprétation :} l'eau chaude, moins dense, monte ; elle cède sa chaleur en surface, se refroidit, devient plus dense et redescend. La chaleur est transportée \emph{par} le mouvement du fluide.
\end{experience}

La convection explique de nombreux phénomènes quotidiens : l'air chaud s'accumule au plafond d'une pièce (d'où l'emplacement des ventilateurs de plafond), le radiateur placé près du sol chauffe toute la chambre, et la \textbf{brise de mer} qui rafraîchit Kribi l'après-midi : le sable chauffé par le soleil réchauffe l'air qui s'élève, et l'air plus frais au-dessus de l'océan s'engouffre vers la côte pour le remplacer.

\begin{attention}[La convection est impossible dans les solides et dans le vide]
La convection exige un \textbf{fluide} capable de se déplacer. Il n'y a donc \textbf{jamais de convection} dans un solide, ni dans le vide. Dire « la chaleur du Soleil arrive par convection » est une erreur classique au Probatoire : entre le Soleil et la Terre règne le vide spatial.
\end{attention}

\courssub{Le rayonnement thermique}

\textbf{Observation.} Approche ta main (sans toucher !) d'un feu de camp ou d'une braise : tu sens la chaleur \emph{instantanément}, même si l'air entre ta main et le feu est froid, et même si le vent souffle dans l'autre sens. Aucun contact, aucun courant d'air ne peut expliquer cette sensation.

\begin{definition}[Rayonnement thermique]
Le \cle{rayonnement} est un mode de transfert de chaleur qui s'effectue par \textbf{ondes électromagnétiques} (essentiellement infrarouges pour les températures usuelles). Il ne nécessite \textbf{aucun support matériel} : il se propage même dans le \textbf{vide}, à la vitesse de la lumière $c = \SI{3,0e8}{m.s^{-1}}$. Tout corps émet un rayonnement d'autant plus intense qu'il est chaud.
\end{definition}

C'est par rayonnement que la chaleur du Soleil traverse environ 150 millions de kilomètres de vide spatial avant d'atteindre le Cameroun. C'est aussi par rayonnement qu'un four chauffe un gâteau, qu'une lampe à incandescence brûle la main qui s'en approche, et que la Terre, la nuit, se refroidit en émettant vers l'espace.

\begin{attention}[Erreur fréquente : « le rayonnement a besoin d'air »]
Beaucoup d'élèves croient que la chaleur du Soleil nous parvient « par l'air ». C'est faux : l'atmosphère ne représente qu'une fine pellicule ; l'essentiel du trajet se fait dans le \textbf{vide}. Le rayonnement est le \textbf{seul} des trois modes capable de traverser le vide. Retiens : \emph{contact} $\rightarrow$ conduction ; \emph{fluide en mouvement} $\rightarrow$ convection ; \emph{vide ou simple distance} $\rightarrow$ rayonnement.
\end{attention}

\courssub{Schéma de synthèse des trois modes}

\begin{popfigure}
\begin{tikzpicture}[scale=0.9, font=\small]
% ---- Panneau 1 : conduction ----
\fill[poporange!70] (0,0) rectangle (4.5,0.8);
\draw[thick, popdark] (0,0) rectangle (4.5,0.8);
\draw[fill=popink] (0.2,-0.35) circle (0.28);
\node[popink, below] at (0.2,-0.7) {\footnotesize flamme};
\foreach \x in {0.7,1.4,2.1,2.8,3.5,4.2} {
  \draw[popdark] (\x,0.4) circle (0.09);
}
\draw[-{Stealth}, thick, poppink] (0.6,1.05) -- (4.3,1.05);
\node[above, poppink] at (2.45,1.05) {\footnotesize chaleur de proche en proche};
\node[below] at (2.25,-0.95) {\textbf{Conduction} (solide)};
% ---- Panneau 2 : convection ----
\begin{scope}[xshift=6.2cm]
\draw[thick, popdark] (0,0) -- (0,1.8) -- (3.2,1.8) -- (3.2,0);
\fill[popblue!25] (0.05,0.05) rectangle (3.15,1.5);
\draw[-{Stealth}, thick, poporange] (1.6,0.25) -- (1.6,1.35);
\draw[-{Stealth}, thick, popblue] (0.45,1.35) .. controls (0.3,0.9) .. (0.45,0.3);
\draw[-{Stealth}, thick, popblue] (2.75,1.35) .. controls (2.9,0.9) .. (2.75,0.3);
\draw[-{Stealth}, thick, poporange] (0.6,1.42) -- (1.45,1.42);
\draw[-{Stealth}, thick, poporange] (1.75,1.42) -- (2.6,1.42);
\draw[fill=popink] (1.6,-0.3) circle (0.22);
\node[below] at (1.6,-0.75) {\textbf{Convection} (fluide)};
\end{scope}
% ---- Panneau 3 : rayonnement ----
\begin{scope}[xshift=11.6cm]
\draw[fill=popgold, draw=poporange, thick] (0,1.4) circle (0.5);
\foreach \a in {20,60,100,140,160} {
  \draw[-{Stealth}, thick, poporange] (\a:0.65) -- ++(\a:0.5);
}
\draw[fill=popblue!50, draw=popdark] (3.4,0.2) circle (0.4);
\foreach \y in {1.15,0.85,0.55} {
  \draw[-{Stealth}, thick, poppink, decorate, decoration={snake, amplitude=0.8mm}] (0.55,\y) -- (2.95,\y);
}
\node[popmuted] at (1.75,1.6) {\footnotesize vide};
\node[below] at (1.75,-0.55) {\textbf{Rayonnement} (vide possible)};
\end{scope}
\end{tikzpicture}
\legende{Les trois modes de transfert de chaleur : conduction dans une barre métallique (sans déplacement de matière), convection dans l'eau d'une casserole (courants de fluide), rayonnement du Soleil vers la Terre (à travers le vide).}
\end{popfigure}

\courssub{Comparaison des trois modes de transfert}

\begin{center}
\rowcolors{1}{popcream}{white}
\begin{tabularx}{\linewidth}{|l|X|X|X|}
\hline
\rowcolor{popblueL} \textbf{Critère} & \textbf{Conduction} & \textbf{Convection} & \textbf{Rayonnement} \\
\hline
Support matériel & Oui (contact direct) & Oui (fluide) & \textbf{Non} (vide possible) \\
\hline
Déplacement de matière & Non & \textbf{Oui} (courants) & Non \\
\hline
Milieux concernés & Surtout les solides & Liquides et gaz & Tous, y compris le vide \\
\hline
Exemple & Manche métallique chaud & Eau qui bout, brise de mer & Chaleur du Soleil, feu de camp \\
\hline
\end{tabularx}
\end{center}

Dans la réalité, les trois modes coexistent souvent : devant un feu de camp, la braise te chauffe par rayonnement, l'air chaud qui monte te chauffe par convection, et le banc sur lequel tu es assis te chauffe par conduction. L'art du physicien est de repérer le mode \textbf{dominant}.

\begin{methode}[Identifier le mode de transfert dominant]
Face à une situation concrète, pose-toi trois questions, dans l'ordre :
\begin{enumerate}
\item \textbf{Y a-t-il contact direct entre le corps chaud et le corps froid, à travers un solide ?} Si oui $\rightarrow$ \cle{conduction}.
\item \textbf{Y a-t-il un fluide (air, eau) qui peut se déplacer entre les deux ?} Si oui $\rightarrow$ \cle{convection}.
\item \textbf{Les deux corps sont-ils séparés par du vide ou une grande distance, sans contact ni courant de fluide ?} Si oui $\rightarrow$ \cle{rayonnement}.
\end{enumerate}
Si plusieurs réponses sont positives, cite tous les modes présents et précise lequel domine.
\end{methode}

\courssub{Applications : isoler, c'est bloquer les trois modes}

Comprendre les trois modes permet de les \emph{combattre} quand on veut conserver la chaleur (ou la fraîcheur) : glacières pour le transport du poisson fumé, thermos pour le café, isolation des maisons.

\begin{exemplebox}[La bouteille thermos : le triple blocage]
Une bouteille thermos garde l'eau chaude (ou l'eau fraîche) pendant des heures grâce à trois dispositifs, chacun neutralisant un mode de transfert :
\begin{itemize}
\item le \textbf{vide} entre les deux parois bloque la \textbf{conduction} et la \textbf{convection} (pas de matière, pas de transfert par contact ni par courant) ;
\item les parois \textbf{argentées} (réfléchissantes) renvoient le \textbf{rayonnement} vers l'intérieur ;
\item le \textbf{bouchon} en liège ou en plastique, mauvais conducteur, bloque la conduction et empêche la convection de l'air au niveau de l'ouverture.
\end{itemize}
La glacière fonctionne sur le même principe : parois épaisses en mousse (air emprisonné, isolant) et surface intérieure claire et lisse.
\end{exemplebox}

\begin{savaistu}[Les cases en banco : la climatisation des ancêtres]
Les cases traditionnelles en \emph{banco} (terre crue) restent remarquablement fraîches en pleine chaleur. La terre est un \textbf{mauvais conducteur thermique} : les murs épais ralentissent tellement la conduction que la chaleur du jour n'a pas le temps de traverser avant la nuit. Le toit de paille, plein d'air emprisonné, ajoute une isolation supplémentaire. Nos anciens maîtrisaient la physique thermique bien avant les manuels ! À l'inverse, les tôles ondulées, excellentes conductrices, transforment les maisons en fours dès midi — d'où l'intérêt des faux plafonds isolants.
\end{savaistu}

\begin{exoresolu}[Identifier les modes de transfert]
\textbf{Énoncé.} Une marmite d'eau chauffe sur un feu de bois dans une cuisine de Yaoundé. Pour chacun des transferts suivants, indiquer le mode dominant et justifier :
\begin{enumerate}
\item la chaleur traverse le fond en aluminium de la marmite ;
\item la chaleur se répartit dans toute l'eau de la marmite ;
\item le cuisinier, debout à un mètre du feu, sent la chaleur sur son visage ;
\item la chaleur monte du feu vers la marmite à travers l'air.
\end{enumerate}
\tcblower
\textbf{Solution.}
\begin{enumerate}
\item \textbf{Conduction} : l'aluminium est un solide (métal, bon conducteur) traversé par contact direct, sans déplacement de matière.
\item \textbf{Convection} : l'eau est un liquide ; l'eau chaude du fond monte, l'eau froide descend : des courants de convection brassent tout le volume.
\item \textbf{Rayonnement} : le cuisinier n'est ni en contact avec le feu, ni dans le courant d'air chaud qui monte verticalement ; la chaleur lui parvient par ondes électromagnétiques (infrarouges), qui traversent l'air sans le chauffer sensiblement.
\item \textbf{Convection} : l'air chauffé par la flamme se dilate et s'élève ; c'est un déplacement de matière fluide qui transporte la chaleur vers le fond de la marmite.
\end{enumerate}
Remarque : ce simple feu de bois met en jeu les trois modes simultanément — c'est presque toujours le cas dans les situations réelles.
\end{exoresolu}

\begin{aretenir}[L'essentiel de la section]
\begin{itemize}
\item \textbf{Conduction} : contact direct, de proche en proche, sans déplacement de matière ; mode des solides. Métaux = bons conducteurs ; bois, liège, air immobile = isolants.
\item \textbf{Convection} : déplacement de matière dans un fluide (courants dus aux variations de densité) ; impossible dans les solides et dans le vide.
\item \textbf{Rayonnement} : ondes électromagnétiques ; seul mode possible dans le vide ; émis par tout corps chaud.
\item Pour identifier un mode : contact ? $\rightarrow$ conduction ; fluide en mouvement ? $\rightarrow$ convection ; vide ou distance ? $\rightarrow$ rayonnement.
\end{itemize}
\end{aretenir}

Maintenant que nous savons \emph{comment} la chaleur se déplace, il reste à apprendre à la \emph{mesurer} : quelle quantité de chaleur faut-il fournir pour chauffer un litre d'eau, ou pour faire fondre un glaçon ? C'est l'objet des sections suivantes, où nous établirons les relations quantitatives de la calorimétrie.

\coursec{Quantité de chaleur sans changement d'état}

Dans la section précédente, nous avons décrit \emph{comment} la chaleur voyage d'un corps à un autre : par conduction dans les solides, par convection dans les fluides, par rayonnement même à travers le vide. Mais une question fondamentale reste ouverte : \emph{combien} de chaleur faut-il fournir à un corps pour élever sa température d'une valeur donnée ? C'est toute la problématique de la cuisine familiale : combien de temps la marmite d'eau du \emph{koki} doit-elle rester sur le feu ? Combien de bois faut-il pour porter à ébullition l'eau du bain ? Répondre à ces questions, c'est entrer dans le cœur de la calorimétrie. Dans toute cette section, on se limite aux situations où le corps étudié \cle{ne change pas d'état} : un solide reste solide, un liquide reste liquide, un gaz reste gaz. Seule sa température varie.

\courssub{Observation expérimentale : de quoi dépend la chaleur reçue ?}

\begin{experience}[Trois expériences de pensée très concrètes]
Considérons les situations suivantes, que l'on peut réaliser dans n'importe quelle cuisine ou au laboratoire du lycée.
\begin{enumerate}
\item \textbf{Influence de la masse.} Sur deux réchauds identiques, on chauffe d'une part $0,5$ L d'eau, d'autre part $2$ L d'eau, tous deux initialement à $25\;^\circ\mathrm{C}$. On constate qu'il faut environ quatre fois plus de temps (donc quatre fois plus de chaleur) pour porter les $2$ L à ébullition. La chaleur requise \cle{croît avec la masse}.
\item \textbf{Influence de la variation de température.} Porter $1$ L d'eau de $25\;^\circ\mathrm{C}$ à $50\;^\circ\mathrm{C}$ demande deux fois moins de chaleur que de le porter à $75\;^\circ\mathrm{C}$. La chaleur requise \cle{croît avec l'écart de température} à réaliser.
\item \textbf{Influence de la nature du corps.} On dépose sur une plaque chaude une casserole en aluminium de $1$ kg et on verse dans une autre $1$ kg d'eau. À puissance de chauffe égale, la température de l'aluminium s'élève beaucoup plus vite que celle de l'eau. La chaleur requise \cle{dépend de la nature} du matériau.
\end{enumerate}
\textbf{Interprétation.} La quantité de chaleur $Q$ nécessaire pour modifier la température d'un corps dépend de trois facteurs : sa masse $m$, la variation de température $\Delta\theta$ souhaitée, et une caractéristique propre au matériau que nous allons définir : la \cle{chaleur massique}.
\end{experience}

\courssub{La chaleur massique (capacité thermique massique)}

\begin{definition}[Chaleur massique]
La \cle{chaleur massique} (ou \cle{capacité thermique massique}) d'un matériau, notée $c$, est la quantité de chaleur qu'il faut fournir à \textbf{un kilogramme} de ce matériau pour élever sa température de \textbf{un kelvin} (ou d'un degré Celsius), sans changement d'état.
\[
\text{Unité SI : } \mathrm{J\cdot kg^{-1}\cdot K^{-1}}
\]
Plus $c$ est grande, plus le matériau est « difficile à chauffer » : il emmagasine beaucoup d'énergie pour une faible élévation de température.
\end{definition}

\begin{aretenir}[Quelques valeurs remarquables de chaleurs massiques]
\begin{center}
\begin{tabularx}{\linewidth}{|l|c|X|}
\hline
\rowcolor{popblueL} \textbf{Matériau} & $c$ (en $\mathrm{J\cdot kg^{-1}\cdot K^{-1}}$) & \textbf{Conséquence pratique} \\
\hline
Eau liquide & $4180$ & Très grande : l'eau se chauffe et se refroidit lentement \\
\hline
Aluminium & $920$ & Chauffe vite : idéal pour les casseroles \\
\hline
Fer & $460$ & Chauffe vite, stocke bien la chaleur (cocotte) \\
\hline
Cuivre & $385$ & Excellent conducteur, chauffe très vite \\
\hline
Air (à pression constante) & $1000$ & Intermédiaire \\
\hline
\end{tabularx}
\end{center}
La valeur de l'eau est exceptionnellement élevée : il faut plus de quatre fois plus d'énergie pour élever de $1$ K la température de $1$ kg d'eau que celle de $1$ kg d'aluminium, et environ neuf fois plus que pour $1$ kg de fer.
\end{aretenir}

\begin{popfigure}
\begin{center}
\begin{tikzpicture}
% Barres horizontales
\fill[popblue]   (0,3.6) rectangle (4.18,4.3);
\fill[poporange] (0,2.4) rectangle (0.92,3.1);
\fill[popgreen]  (0,1.2) rectangle (0.46,1.9);
\fill[poppurple] (0,0.0) rectangle (0.385,0.7);
% Étiquettes
\node[left] at (0,3.95) {\textbf{Eau}};
\node[left] at (0,2.75) {\textbf{Aluminium}};
\node[left] at (0,1.55) {\textbf{Fer}};
\node[left] at (0,0.35) {\textbf{Cuivre}};
\node[right] at (4.18,3.95) {$4180$};
\node[right] at (0.92,2.75) {$920$};
\node[right] at (0.46,1.55) {$460$};
\node[right] at (0.385,0.35) {$385$};
% Axe
\draw[->,thick] (0,-0.4) -- (4.8,-0.4) node[below left] {$c$ ($\mathrm{J\cdot kg^{-1}\cdot K^{-1}}$)};
\foreach \x/\lab in {1/1000,2/2000,3/3000,4/4000}{
  \draw (\x,-0.45) -- (\x,-0.35) node[below] {\scriptsize \lab};
}
\end{tikzpicture}
\end{center}
\legende{Comparaison des chaleurs massiques de quelques matériaux usuels (échelle : 1 cm pour $1000\ \mathrm{J\cdot kg^{-1}\cdot K^{-1}}$). La barre de l'eau écrase toutes les autres.}
\end{popfigure}

\begin{savaistu}[L'océan, régulateur du climat camerounais]
C'est précisément grâce à la chaleur massique exceptionnelle de l'eau que le climat du littoral camerounais (Douala, Kribi, Limbé) est si doux comparé à celui de l'intérieur des terres. Le jour, l'océan Atlantique absorbe d'énormes quantités de chaleur solaire sans que sa température ne s'élève beaucoup ; la nuit, il restitue lentement cette chaleur. Résultat : les écarts de température entre le jour et la nuit, et entre les saisons, sont fortement atténués près des côtes. À l'inverse, dans le Grand Nord (Maroua), loin de toute grande masse d'eau, le sol sablonneux (faible $c$) chauffe très vite le jour et se refroidit vite la nuit : les amplitudes thermiques y sont considérables. Le même principe explique pourquoi l'eau est utilisée comme fluide de refroidissement dans les moteurs des mototaxis et dans les radiateurs automobiles.
\end{savaistu}

\courssub{La relation fondamentale $Q = m\,c\,(\theta_f - \theta_i)$}

\begin{propriete}[Expression de la quantité de chaleur (admise)]
Lorsqu'un corps de masse $m$, de chaleur massique $c$, voit sa température passer de $\theta_i$ à $\theta_f$ \textbf{sans changer d'état}, la quantité de chaleur qu'il échange avec l'extérieur est :
\[
\boxed{\,Q = m \cdot c \cdot (\theta_f - \theta_i)\,}
\]
avec $Q$ en joules ($\mathrm{J}$), $m$ en kilogrammes ($\mathrm{kg}$), $c$ en $\mathrm{J\cdot kg^{-1}\cdot K^{-1}}$, et $\theta_f - \theta_i$ en kelvins ($\mathrm{K}$) ou en degrés Celsius ($^\circ\mathrm{C}$) — un \emph{écart} de température a la même valeur dans les deux échelles.
\par\smallskip
Cette relation n'est pas démontrable au niveau Première : elle est \textbf{admise} et \textbf{justifiée expérimentalement}. En TP, on vérifie la proportionnalité : à masse et matériau fixés, la chaleur fournie par un thermoplongeur de puissance connue ($Q = \mathcal{P}\,\Delta t$) est proportionnelle à l'élévation de température mesurée ; le rapport $Q/\Delta\theta$ est constant et vaut $m\,c$.
\end{propriete}

\textbf{Vérification par analyse dimensionnelle.} Le produit $m\,c\,\Delta\theta$ s'exprime en
\[
\mathrm{kg} \times \mathrm{J\cdot kg^{-1}\cdot K^{-1}} \times \mathrm{K} = \mathrm{J},
\]
ce qui est bien l'unité d'une quantité de chaleur. La formule est donc homogène.

\begin{attention}[Convention de signe : qui reçoit, qui cède ?]
En thermodynamique, on adopte le point de vue du \textbf{corps étudié} :
\begin{itemize}
\item si le corps \textbf{reçoit} de la chaleur ($\theta_f > \theta_i$), alors $Q > 0$ ;
\item si le corps \textbf{cède} de la chaleur ($\theta_f < \theta_i$), alors $Q < 0$.
\end{itemize}
Le signe de $Q$ est donc automatiquement donné par celui de $(\theta_f - \theta_i)$. Ne prends \emph{jamais} la valeur absolue de l'écart de température dans la formule : tu perdrais l'information essentielle du sens de l'échange, indispensable pour appliquer le principe des échanges de chaleur (section 5).
\end{attention}

\begin{popfigure}
\begin{center}
\begin{tikzpicture}
\begin{axis}[
  width=0.85\linewidth, height=6cm,
  xlabel={Chaleur reçue $Q$ (kJ)},
  ylabel={Température $\theta$ ($^\circ$C)},
  xmin=0, xmax=340, ymin=15, ymax=95,
  grid=both, grid style={popline!40},
  xtick={0,41.8,83.6,125.4,167.2,209,250.8,292.6,334.4},
  x tick label style={/pgf/number format/fixed, /pgf/number format/precision=0, font=\scriptsize},
  xticklabels={0,41,83,125,167,209,250,292,334},
  legend style={at={(0.03,0.97)}, anchor=north west, font=\small},
]
\addplot[popcurve, domain=0:334.4, samples=50] {20 + x/4.18};
\addlegendentry{1 kg d'eau : $\theta = 20 + Q/(m c)$}
\addplot[mark=*, poporange, only marks] coordinates {(125.4,50)};
\node[poporange, font=\scriptsize, anchor=west] at (axis cs:130,48) {$Q = 125{,}4$ kJ $\Rightarrow \theta = 50\,^\circ$C};
\end{axis}
\end{tikzpicture}
\end{center}
\legende{Température d'un kilogramme d'eau initialement à $20\;^\circ$C en fonction de la chaleur reçue. La courbe est une droite de pente $\dfrac{1}{m\,c} = \dfrac{1}{4180}\ \mathrm{K\cdot J^{-1}}$ : la proportionnalité entre $Q$ et $\Delta\theta$ est la signature de la loi $Q = m\,c\,\Delta\theta$.}
\end{popfigure}

\courssub{La capacité calorifique d'un corps donné}

Pour un \emph{objet particulier} (une casserole précise, un radiateur, l'eau d'un ballon d'eau chaude), la masse est fixée une fois pour toutes. Il est alors commode de regrouper $m$ et $c$ en une seule constante.

\begin{definition}[Capacité calorifique]
La \cle{capacité calorifique} (ou capacité thermique) d'un corps donné est le produit de sa masse par sa chaleur massique :
\[
C = m \cdot c
\]
Elle s'exprime en $\mathrm{J\cdot K^{-1}}$ et représente la chaleur nécessaire pour élever de $1$ K la température de \textbf{ce corps entier}. La quantité de chaleur échangée s'écrit alors simplement :
\[
Q = C \cdot (\theta_f - \theta_i)
\]
\end{definition}

\begin{exemplebox}[Capacité calorifique d'une casserole d'eau]
Une casserole contient $2$ L d'eau, soit $m = 2$ kg. Sa capacité calorifique (en négligeant celle de la casserole elle-même) vaut :
\[
C = m\,c = 2 \times 4180 = 8360\ \mathrm{J\cdot K^{-1}}.
\]
Chaque degré supplémentaire « coûte » donc $8360$ J. Pour élever cette eau de $25\;^\circ\mathrm{C}$ à $100\;^\circ\mathrm{C}$ :
\[
Q = C\,\Delta\theta = 8360 \times 75 = 627\,000\ \mathrm{J} = 627\ \mathrm{kJ}.
\]
À titre de comparaison, un réchaud à bois domestique délivre une puissance utile de l'ordre de $1$ kW : il faudrait donc environ $627$ s, soit plus de dix minutes, en supposant aucune perte.
\end{exemplebox}

\courssub{Méthode de calcul et exemples résolus}

\begin{methode}[Calculer une quantité de chaleur sans changement d'état]
\begin{enumerate}
\item \textbf{Identifier} le corps qui échange la chaleur, sa masse $m$, sa nature (donc $c$), ses températures initiale $\theta_i$ et finale $\theta_f$.
\item \textbf{Convertir} toutes les données en unités SI : la masse en kilogrammes (diviser par $1000$ si elle est en grammes), $c$ en $\mathrm{J\cdot kg^{-1}\cdot K^{-1}}$.
\item \textbf{Calculer} $Q = m\,c\,(\theta_f - \theta_i)$ en respectant l'ordre $\theta_f - \theta_i$ (jamais l'inverse).
\item \textbf{Conclure} en interprétant le signe : $Q > 0$ signifie que le corps reçoit de la chaleur ; $Q < 0$ qu'il en cède. Donner le résultat avec son unité et un nombre raisonnable de chiffres significatifs.
\end{enumerate}
\end{methode}

\begin{exemplebox}[Chauffer 500 g d'eau de 20 °C à 80 °C]
On veut chauffer $m = 500$ g d'eau de $\theta_i = 20\;^\circ\mathrm{C}$ à $\theta_f = 80\;^\circ\mathrm{C}$.
\begin{itemize}
\item Conversion : $m = 500\ \mathrm{g} = 0{,}5\ \mathrm{kg}$ ;
\item Écart de température : $\theta_f - \theta_i = 80 - 20 = 60\ \mathrm{K}$ ;
\item Application numérique :
\[
Q = m\,c\,(\theta_f - \theta_i) = 0{,}5 \times 4180 \times 60 = 125\,400\ \mathrm{J} \approx 1,25 \times 10^{5}\ \mathrm{J}.
\]
\end{itemize}
$Q > 0$ : l'eau \textbf{reçoit} environ $125$ kJ de la source de chaleur.
\end{exemplebox}

\begin{exoresolu}[Refroidissement d'un bloc de fer]
\textbf{Énoncé.} Un forgeron de Foumban sort de la forge un bloc de fer de masse $m = 3$ kg porté à $\theta_i = 600\;^\circ\mathrm{C}$ et le laisse refroidir à l'air libre jusqu'à la température ambiante $\theta_f = 30\;^\circ\mathrm{C}$. On donne $c_{\text{fer}} = 460\ \mathrm{J\cdot kg^{-1}\cdot K^{-1}}$.
\begin{enumerate}
\item Calculer la quantité de chaleur échangée par le bloc de fer avec l'extérieur.
\item Interpréter le signe du résultat.
\item Quelle masse d'eau initialement à $30\;^\circ\mathrm{C}$ pourrait-on porter à $80\;^\circ\mathrm{C}$ en récupérant intégralement cette chaleur ?
\end{enumerate}
\tcblower
\textbf{Solution détaillée.}
\begin{enumerate}
\item Le bloc de fer est le corps étudié. Toutes les données sont déjà en unités SI :
\begin{align*}
Q &= m \cdot c_{\text{fer}} \cdot (\theta_f - \theta_i) \\
  &= 3 \times 460 \times (30 - 600) \\
  &= 3 \times 460 \times (-570) \\
  &= -786\,600\ \mathrm{J} \approx -7,87 \times 10^{5}\ \mathrm{J}.
\end{align*}
\item $Q < 0$ : le bloc de fer \textbf{cède} environ $787$ kJ à l'air ambiant. C'est cohérent avec le sens physique du refroidissement.
\item L'eau récupère $Q' = +786\,600$ J. Pour elle, $\theta_f' - \theta_i' = 80 - 30 = 50$ K. D'où :
\[
m_{\text{eau}} = \frac{Q'}{c_{\text{eau}}\,\Delta\theta'} = \frac{786\,600}{4180 \times 50} \approx 3{,}76\ \mathrm{kg}.
\]
On peut donc chauffer environ $3,8$ L d'eau de $30\;^\circ\mathrm{C}$ à $80\;^\circ\mathrm{C}$ grâce au seul refroidissement du bloc de fer — alors que le fer est trois fois plus lourd que cette eau... non : remarque que les masses sont comparables, car la chaleur massique de l'eau est environ neuf fois celle du fer, mais l'écart de température du fer ($570$ K) est bien plus grand que celui de l'eau ($50$ K).
\end{enumerate}
\end{exoresolu}

\begin{attention}[Les deux erreurs classiques du calcul de $Q$]
\begin{itemize}
\item \textbf{Oublier de convertir la masse en kilogrammes.} Écrire $Q = 500 \times 4180 \times 60$ pour $500$ g d'eau donne un résultat \emph{mille fois trop grand} ($1,25 \times 10^{8}$ J !). La masse doit \emph{toujours} être exprimée en kg, car $c$ est donnée par kilogramme.
\item \textbf{Prendre $\Delta\theta$ « positif » sans réfléchir au sens de l'échange.} Pour un corps qui se refroidit, écrire $Q = m\,c\,\Delta\theta$ avec $\Delta\theta = +570$ K donne un $Q$ positif, ce qui signifierait que le fer \emph{reçoit} de la chaleur en refroidissant : absurde ! Écris systématiquement $\theta_f - \theta_i$, dans cet ordre, et laisse le signe parler.
\item \textbf{Confondre température et écart de température.} On ne convertit jamais un \emph{écart} de température en kelvins en ajoutant 273 : un écart de $60\;^\circ$C est un écart de $60$ K.
\end{itemize}
\end{attention}

\begin{exercice}[Pour t'entraîner]
Un chauffe-eau électrique porte $50$ L d'eau de $20\;^\circ\mathrm{C}$ à $65\;^\circ\mathrm{C}$. On donne $c_{\text{eau}} = 4180\ \mathrm{J\cdot kg^{-1}\cdot K^{-1}}$.
\begin{enumerate}
\item Calculer la quantité de chaleur reçue par l'eau.
\item La résistance du chauffe-eau a une puissance de $2000\ \mathrm{W}$. En supposant que toute l'énergie électrique est transmise à l'eau, calculer la durée du chauffage en minutes.
\end{enumerate}
\end{exercice}

\begin{corrige}[Exercice : chauffe-eau électrique]
\begin{enumerate}
\item Masse d'eau : $m = 50\ \mathrm{kg}$ (1 L d'eau a une masse de 1 kg). Écart : $\theta_f - \theta_i = 65 - 20 = 45\ \mathrm{K}$.
\[
Q = m\,c\,(\theta_f - \theta_i) = 50 \times 4180 \times 45 = 9\,405\,000\ \mathrm{J} \approx 9,4 \times 10^{6}\ \mathrm{J}.
\]
$Q > 0$ : l'eau reçoit environ $9,4$ MJ.
\item L'énergie fournie par la résistance pendant la durée $\Delta t$ est $Q = \mathcal{P}\,\Delta t$, d'où :
\[
\Delta t = \frac{Q}{\mathcal{P}} = \frac{9\,405\,000}{2000} \approx 4703\ \mathrm{s} \approx 78\ \mathrm{min}.
\]
Il faut donc environ 1 h 18 min de chauffe, ce qui explique pourquoi on allume le chauffe-eau bien avant la douche du matin.
\end{enumerate}
\end{corrige}

En résumé, dès qu'un corps change de température sans changer d'état, la chaleur échangée se calcule par la relation $Q = m\,c\,(\theta_f - \theta_i)$, dont le signe indique le sens de l'échange. Mais que se passe-t-il lorsque le corps atteint la température d'un changement d'état — fusion de la glace, ébullition de l'eau ? La température cesse alors de varier pendant que la chaleur continue d'être échangée : c'est la notion de \cle{chaleur latente}, objet de la section suivante.

\coursec{Chaleur latente et changements d'état}

Dans la section précédente, nous avons appris à calculer la quantité de chaleur $Q = m\,c\,\Delta\theta$ qu'il faut fournir à un corps pour élever sa température de $\Delta\theta$, \emph{sans qu'il change d'état}. Mais que se passe-t-il lorsqu'on chauffe de la glace jusqu'à la faire fondre, ou de l'eau jusqu'à la faire bouillir ? L'expérience montre un phénomène surprenant : à certains moments, la température refuse de monter, alors même qu'on continue de chauffer. Cette section est consacrée à ces \cle{paliers de température} et à la grandeur qui les caractérise : la \cle{chaleur latente de changement d'état}.

\courssub{Observation expérimentale : le palier de température}

\begin{experience}[Fusion de la glace : la température qui « attend »]
\textbf{Matériel :} un bécher contenant de la glace pilée (environ \SI{100}{g}) sortie d'un congélateur à $\theta_i = -10\ \mathrm{^\circ C}$, un thermomètre, un agitateur, une plaque chauffante de puissance constante, un chronomètre.

\textbf{Protocole :} on chauffe la glace à puissance constante et on relève la température toutes les minutes, en agitant régulièrement pour homogénéiser.

\textbf{Observations :}
\begin{itemize}
\item de $-10\ \mathrm{^\circ C}$ à $0\ \mathrm{^\circ C}$, la température de la glace monte régulièrement ;
\item arrivée à $0\ \mathrm{^\circ C}$, la glace commence à fondre : \textbf{la température reste bloquée à $0\ \mathrm{^\circ C}$} tant qu'il reste de la glace, même si l'on chauffe toujours ;
\item une fois toute la glace fondue, la température de l'eau remonte jusqu'à $100\ \mathrm{^\circ C}$ ;
\item à $100\ \mathrm{^\circ C}$, l'eau bout : \textbf{nouveau palier de température} pendant toute la vaporisation.
\end{itemize}

\textbf{Interprétation :} pendant un changement d'état, la chaleur fournie ne sert plus à augmenter l'agitation thermique des particules (donc la température), mais à \emph{briser les liaisons} entre les particules pour les éloigner les unes des autres : c'est un travail de « désorganisation » de la matière. L'énergie est absorbée sans élévation de température : on dit qu'elle est \emph{latente} (cachée).
\end{experience}

\begin{propriete}[Température constante pendant un changement d'état — admise]
Pendant tout changement d'état d'un \cle{corps pur} effectué à pression constante, la \textbf{température reste constante}, égale à la température de changement d'état ($\theta_{\text{fus}}$, $\theta_{\text{éb}}$...), même si le système échange de la chaleur avec l'extérieur. La quantité de chaleur échangée pendant le changement d'état complet d'une masse $m$ est :
\[ Q = m \times L \]
où $L$ est la chaleur latente de changement d'état du corps. Pour les transformations inverses (solidification, liquéfaction), le corps \emph{cède} cette chaleur : $Q = -m\,L$.
\end{propriete}

\begin{popfigure}
\begin{center}
\begin{tikzpicture}
\begin{axis}[
  width=0.95\linewidth, height=7.5cm,
  xlabel={Chaleur fournie $Q$ (kJ)}, ylabel={Température $\theta$ ($^\circ$C)},
  xmin=0, xmax=100, ymin=-20, ymax=120,
  xtick={0,5,15,50,95}, ytick={-10,0,100},
  grid=both, grid style={popline!40},
  axis lines=left, thick
]
% palier fusion : de Q=5 à Q=15 à 0°C ; chauffage eau : 15->50 à 100°C ; palier vap : 50->95
\addplot[popcurve,domain=0:5,samples=10] {2*x - 10};
\addplot[popcurve,domain=5:15,samples=2] {0};
\addplot[popcurve,domain=15:50,samples=10] {(100/35)*(x-15)};
\addplot[popcurve,domain=50:95,samples=2] {100};
% annotations
\draw[dashed, popmuted] (axis cs:0,0) -- (axis cs:15,0);
\draw[dashed, popmuted] (axis cs:0,100) -- (axis cs:95,100);
\node[popdark, font=\small, anchor=south] at (axis cs:10,1) {palier de fusion ($0\ ^\circ$C)};
\node[popdark, font=\small, anchor=south] at (axis cs:72,101) {palier de vaporisation ($100\ ^\circ$C)};
\node[popblue, font=\small, rotate=55, anchor=south] at (axis cs:2.5,-6) {glace};
\node[popblue, font=\small, rotate=70, anchor=south west] at (axis cs:20,10) {eau liquide};
\node[popblue, font=\small, anchor=south] at (axis cs:97,103) {vapeur};
\end{axis}
\end{tikzpicture}
\end{center}
\legende{Courbe de chauffage de l'eau (pression atmosphérique) : température en fonction de la chaleur fournie. Les deux paliers horizontaux correspondent à la fusion ($0\ \mathrm{^\circ C}$) et à la vaporisation ($100\ \mathrm{^\circ C}$) : pendant ces étapes, la chaleur fournie ne modifie pas la température.}
\end{popfigure}

\courssub{Chaleur latente de changement d'état}

\begin{definition}[Chaleur latente de changement d'état]
La \cle{chaleur latente de changement d'état} $L$ d'un corps pur est la quantité de chaleur qu'il faut fournir à \textbf{1 kg} de ce corps pour le faire changer d'état, \textbf{à température constante} (et à pression constante).
\begin{itemize}
\item \cle{Chaleur latente de fusion} $L_f$ : chaleur nécessaire pour faire passer 1 kg de corps de l'état solide à l'état liquide, à sa température de fusion ;
\item \cle{Chaleur latente de vaporisation} $L_v$ : chaleur nécessaire pour faire passer 1 kg de corps de l'état liquide à l'état gazeux, à sa température d'ébullition.
\end{itemize}
Unité SI : le \textbf{joule par kilogramme}, de symbole $\mathrm{J \cdot kg^{-1}}$.
\end{definition}

\begin{aretenir}[Valeurs remarquables pour l'eau]
\begin{center}
\begin{tabular}{|l|c|c|}
\hline
\rowcolor{popblueL} Changement d'état & Température & Chaleur latente \\
\hline
Fusion de la glace & $0\ \mathrm{^\circ C}$ & $L_f = \num{3,3e5}\ \mathrm{J \cdot kg^{-1}}$ \\
\hline
Vaporisation de l'eau & $100\ \mathrm{^\circ C}$ & $L_v = \num{2,26e6}\ \mathrm{J \cdot kg^{-1}}$ \\
\hline
\end{tabular}
\end{center}
Remarque : $L_v \approx 7\,L_f$ : il faut beaucoup plus d'énergie pour séparer complètement les molécules (vaporisation) que pour simplement désordonner le solide (fusion). À titre de comparaison, vaporiser 1 kg d'eau à $100\ \mathrm{^\circ C}$ demande autant d'énergie que chauffer 1 kg d'eau de $0\ \mathrm{^\circ C}$ à environ $540\ \mathrm{^\circ C}$ !
\end{aretenir}

\begin{propriete}[Quantité de chaleur échangée lors d'un changement d'état]
Pour une masse $m$ d'un corps pur changeant d'état à température constante :
\begin{itemize}
\item \textbf{fusion, vaporisation} (le corps \emph{reçoit} de la chaleur) : \[ Q = +\,m \times L \]
\item \textbf{solidification, liquéfaction} (le corps \emph{cède} de la chaleur) : \[ Q = -\,m \times L \]
\end{itemize}
Vérification par analyse dimensionnelle : $[m \times L] = \mathrm{kg} \times \mathrm{J \cdot kg^{-1}} = \mathrm{J}$ : le résultat est bien homogène à une énergie.
\end{propriete}

\begin{exemplebox}[Fonte de 200 g de glace à $0\ \mathrm{^\circ C}$]
Quelle quantité de chaleur faut-il fournir pour faire fondre complètement $m = 200\ \mathrm{g}$ de glace déjà à $0\ \mathrm{^\circ C}$ ?

La glace est \emph{déjà} à sa température de fusion : il n'y a pas d'étape de chauffage préalable. On applique directement :
\[ Q = m \times L_f = \num{0,200} \times \num{3,3e5} = \num{6,6e4}\ \mathrm{J} \]
Il faut donc fournir $Q = \SI{66}{kJ}$, soit l'énergie dégagée par une lampe de \SI{60}{W} allumée pendant plus de 18 minutes. On comprend pourquoi la glace d'un glaçon ne fond pas instantanément dans un jus de fruit !
\end{exemplebox}

\begin{attention}[Piège classique : ne pas appliquer $Q = m\,c\,\Delta\theta$ pendant un palier]
Pendant un changement d'état, $\Delta\theta = 0$. Appliquer la formule $Q = m\,c\,\Delta\theta$ donnerait $Q = 0$, ce qui est absurde : on fournit bien de la chaleur !

\textbf{Règle de sécurité :}
\begin{itemize}
\item la température \emph{change} $\Rightarrow$ on utilise $Q = m\,c\,\Delta\theta$ ;
\item l'\emph{état} change (à température constante) $\Rightarrow$ on utilise $Q = m\,L$.
\end{itemize}
Autre erreur fréquente : oublier de convertir la masse en kilogrammes ($200\ \mathrm{g} = \num{0,200}\ \mathrm{kg}$) alors que $L$ est donnée en $\mathrm{J \cdot kg^{-1}}$.
\end{attention}

\courssub{Calculs combinés : découpage en étapes}

Dans les problèmes réels, le corps subit souvent \emph{plusieurs} étapes successives : chauffage, changement d'état, nouveau chauffage... La chaleur totale est la somme des chaleurs de chaque étape.

\begin{methode}[Calculer une chaleur totale en plusieurs étapes]
\begin{enumerate}
\item \textbf{Repérer les étapes} sur le « chemin thermique » : chaque portion où la température varie (formule $m\,c\,\Delta\theta$) et chaque palier de changement d'état (formule $m\,L$). Attention : la capacité thermique massique $c$ dépend de l'état ($c_{\text{glace}} \neq c_{\text{eau}}$).
\item \textbf{Calculer séparément} chaque quantité de chaleur $Q_1, Q_2, Q_3, \ldots$ en veillant aux unités (masses en kg, températures en $^\circ$C ou K pour les écarts).
\item \textbf{Additionner algébriquement} : $Q_{\text{totale}} = Q_1 + Q_2 + Q_3 + \cdots$
\item \textbf{Vérifier} l'ordre de grandeur et le signe du résultat (positif si le corps reçoit de la chaleur).
\end{enumerate}
\end{methode}

\begin{exoresolu}[De la glace à $-5\ \mathrm{^\circ C}$ à la vapeur à $100\ \mathrm{^\circ C}$]
\textbf{Énoncé.} On veut transformer $m = 500\ \mathrm{g}$ de glace prise à $\theta_i = -5\ \mathrm{^\circ C}$ en vapeur d'eau à $100\ \mathrm{^\circ C}$, sous pression atmosphérique. Calculer la quantité de chaleur totale à fournir.

Données : $c_{\text{glace}} = \num{2,1e3}\ \mathrm{J \cdot kg^{-1} \cdot ^\circ C^{-1}}$ ; $c_{\text{eau}} = \num{4,18e3}\ \mathrm{J \cdot kg^{-1} \cdot ^\circ C^{-1}}$ ; $L_f = \num{3,3e5}\ \mathrm{J \cdot kg^{-1}}$ ; $L_v = \num{2,26e6}\ \mathrm{J \cdot kg^{-1}}$.

\tcblower
\textbf{Solution détaillée.}

\textbf{Étape 1 — Chauffage de la glace de $-5\ \mathrm{^\circ C}$ à $0\ \mathrm{^\circ C}$ :}
\[ Q_1 = m\,c_{\text{glace}}\,\Delta\theta_1 = \num{0,500} \times \num{2,1e3} \times 5 = \num{5,25e3}\ \mathrm{J} \]

\textbf{Étape 2 — Fusion complète de la glace à $0\ \mathrm{^\circ C}$ :}
\[ Q_2 = m \times L_f = \num{0,500} \times \num{3,3e5} = \num{1,65e5}\ \mathrm{J} \]

\textbf{Étape 3 — Chauffage de l'eau de $0\ \mathrm{^\circ C}$ à $100\ \mathrm{^\circ C}$ :}
\[ Q_3 = m\,c_{\text{eau}}\,\Delta\theta_3 = \num{0,500} \times \num{4,18e3} \times 100 = \num{2,09e5}\ \mathrm{J} \]

\textbf{Étape 4 — Vaporisation complète de l'eau à $100\ \mathrm{^\circ C}$ :}
\[ Q_4 = m \times L_v = \num{0,500} \times \num{2,26e6} = \num{1,13e6}\ \mathrm{J} \]

\textbf{Bilan :}
\begin{align*}
Q_{\text{totale}} &= Q_1 + Q_2 + Q_3 + Q_4 \\
&= \num{5,25e3} + \num{1,65e5} + \num{2,09e5} + \num{1,13e6} \\
&\approx \num{1,51e6}\ \mathrm{J}
\end{align*}

Il faut fournir environ $Q_{\text{totale}} \approx \SI{1,51}{MJ}$. Remarque : l'étape de vaporisation ($Q_4$) représente à elle seule environ $75\ \%$ de l'énergie totale — c'est l'étape de loin la plus coûteuse, ce qui confirme que $L_v \gg L_f$.
\end{exoresolu}

\courssub{Applications dans la vie quotidienne et la technologie}

\begin{itemize}
\item \textbf{Fonte des glaces et climat.} La fonte des glaciers exige d'énormes quantités de chaleur ($L_f$ élevée) : elle se fait à $0\ \mathrm{^\circ C}$ sans élévation de température, ce qui « amortit » le réchauffement local mais traduit un apport énergétique considérable.
\item \textbf{Cuisson à la vapeur.} La vapeur d'eau à $100\ \mathrm{^\circ C}$ transporte une grande énergie ($L_v$) qu'elle restitue en se liquéfiant au contact des aliments : la cuisson à la vapeur est rapide et homogène. C'est aussi pourquoi une brûlure par la vapeur est bien plus grave qu'une brûlure par l'eau bouillante : la vapeur cède $m\,L_v$ \emph{en plus} du refroidissement de l'eau formée.
\item \textbf{Réfrigération et climatisation.} Un fluide frigorigène s'évapore dans le circuit intérieur du réfrigérateur en \emph{prélevant} sa chaleur latente de vaporisation aux denrées (il les refroidit), puis se liquéfie à l'extérieur en rejetant cette chaleur. Le même principe équipe les climatiseurs des hôpitaux de Yaoundé et de Douala.
\item \textbf{Le glaçon dans la boisson.} Tant qu'il reste de la glace, le mélange reste à $0\ \mathrm{^\circ C}$ : la boisson reste fraîche durablement.
\end{itemize}

\begin{savaistu}[Pourquoi la transpiration rafraîchit-elle ?]
Quand il fait très chaud à Douala ou à Maroua, ton corps sécrète de la sueur. Pour s'évaporer, chaque kilogramme de sueur doit prélever environ $L_v \approx \num{2,4e6}\ \mathrm{J \cdot kg^{-1}}$ (à la température de la peau) : cette chaleur est prise \emph{à ton corps}, dont la température baisse. La vaporisation de seulement \SI{30}{g} de sueur évacue environ \SI{72}{kJ} ! C'est pourquoi on se sent rafraîchi quand un ventilateur souffle sur la peau moite : l'air en mouvement accélère l'évaporation, donc le prélèvement de chaleur. C'est aussi pour cela que la chaleur humide est plus éprouvante que la chaleur sèche : l'air déjà saturé de vapeur d'eau évapore mal la sueur, et le corps peine à se refroidir.
\end{savaistu}

\begin{exercice}[Chaleur dégagée par la solidification]
De l'eau liquide de masse $m = 250\ \mathrm{g}$ à $0\ \mathrm{^\circ C}$ se solidifie complètement à $0\ \mathrm{^\circ C}$ dans un congélateur. Calculer la quantité de chaleur échangée par l'eau, en précisant son signe et sa signification physique. On donne $L_f = \num{3,3e5}\ \mathrm{J \cdot kg^{-1}}$.
\end{exercice}

\begin{corrige}[Exercice : solidification de l'eau]
La solidification est la transformation inverse de la fusion : l'eau \emph{cède} de la chaleur au congélateur. On applique :
\[ Q = -m \times L_f = -\num{0,250} \times \num{3,3e5} = -\num{8,25e4}\ \mathrm{J} \]
Le signe négatif signifie que l'eau perd de l'énergie : elle cède $8,25 \times 10^{4}\ \mathrm{J}$ au milieu extérieur. C'est cette chaleur que le circuit du congélateur doit évacuer.
\end{corrige}

\begin{attention}[Récapitulatif des pièges de la section]
\begin{itemize}
\item Ne jamais écrire $Q = m\,c\,\Delta\theta$ sur un palier de changement d'état ($\Delta\theta = 0$ mais $Q \neq 0$).
\item Ne pas confondre $L_f$ et $L_v$ : la vaporisation exige environ 7 fois plus d'énergie que la fusion pour l'eau.
\item Toujours convertir la masse en kilogrammes.
\item Dans un calcul multi-étapes, utiliser la \emph{bonne} capacité thermique massique à chaque étape : $c_{\text{glace}}$ pour la glace, $c_{\text{eau}}$ pour l'eau liquide.
\item Pour les transformations inverses (solidification, liquéfaction), ne pas oublier le signe moins : le corps cède la chaleur.
\end{itemize}
\end{attention}

Tu maîtrises maintenant les deux briques de base du calcul thermique : $Q = m\,c\,\Delta\theta$ quand la température varie, et $Q = m\,L$ quand l'état change. Il ne reste plus qu'à réunir ces outils dans un dispositif expérimental fermé, le calorimètre, et à exploiter le principe des échanges de chaleur : ce sera l'objet de la section suivante.

\coursec{Le calorimètre et le principe des échanges de chaleur}

Dans la section précédente, nous avons appris à calculer la quantité de chaleur échangée par un corps, que ce soit lors d'une simple variation de température ($Q = m\,c\,\Delta\theta$) ou lors d'un changement d'état ($Q = m\,L$). Mais une question pratique demeure : comment mesurer expérimentalement ces quantités de chaleur ? Comment déterminer, au laboratoire, la chaleur massique du cuivre ou la chaleur latente de fusion de la glace ? Il faut pour cela un instrument capable de « piéger » la chaleur échangée entre les corps étudiés, sans en laisser fuir vers l'extérieur. Cet instrument existe : c'est le \cle{calorimètre}. Cette section lui est consacrée, ainsi qu'au principe fondamental qui gouverne tous les échanges de chaleur en enceinte isolée.

\courssub{Fonction et description du calorimètre}

\begin{definition}[Calorimètre]
Un \cle{calorimètre} est une enceinte \textbf{thermiquement isolée} (adiabatique) dans laquelle on réalise des échanges de chaleur entre plusieurs corps, de manière à pouvoir mesurer les quantités de chaleur échangées. L'isolation thermique garantit que les échanges avec le milieu extérieur sont négligeables pendant la durée de l'expérience.
\end{definition}

L'idée est simple : si la chaleur ne peut ni entrer ni sortir, alors toute la chaleur cédée par un corps chaud est obligatoirement récupérée par les corps froids placés avec lui. Il suffit alors de connaître les propriétés thermiques des uns pour déduire celles des autres.

Le calorimètre de laboratoire le plus répandu est le \textbf{vase Dewar} (du nom du physicien écossais James Dewar, 1892). Il comprend :

\begin{itemize}
\item un \cle{vase intérieur} (récipient en verre ou en métal) qui contient le liquide, généralement de l'eau ;
\item des \cle{parois isolantes} : double paroi avec vide d'air ou matériau isolant (liège, polystyrène), qui supprime la conduction et la convection à travers les parois ; les surfaces argentées réduisent le rayonnement ;
\item un \cle{couvercle} isolant, percé d'orifices, qui limite les échanges par convection avec l'air extérieur et l'évaporation ;
\item un \cle{thermomètre} pour suivre l'évolution de la température ;
\item un \cle{agitateur} (tige métallique en anneau) qui homogénéise la température du mélange.
\end{itemize}

\begin{popfigure}
\begin{center}
\begin{tikzpicture}[scale=1.0, every node/.style={font=\small}]
% Enveloppe extérieure
\draw[thick, popdark, fill=popcream] (-2.6,-3.2) rectangle (2.6,0.4);
% Isolant (hachuré)
\fill[pattern=north east lines, pattern color=popmuted] (-2.3,-2.9) rectangle (2.3,0.1);
% Vase intérieur
\draw[thick, popblue, fill=popblueL] (-1.7,-2.6) rectangle (1.7,-0.2);
% Eau
\fill[popblue!40] (-1.65,-2.55) rectangle (1.65,-1.1);
\draw[popblue, thick] (-1.65,-1.1) -- (1.65,-1.1);
% Couvercle
\draw[thick, popdark, fill=popgoldL] (-2.6,0.4) rectangle (2.6,0.75);
% Thermomètre
\draw[thick, popink, fill=white] (0.7,0.75) rectangle (0.95,2.2);
\fill[popink] (0.7,-1.6) rectangle (0.95,0.75);
\draw[popink, fill=popink] (0.825,-1.75) circle (0.16);
% Agitateur
\draw[thick, poporange] (-0.9,0.75) -- (-0.9,1.9);
\draw[thick, poporange] (-1.15,1.9) -- (-0.65,1.9);
\draw[thick, poporange] (-0.9,-1.9) circle (0.28);
\draw[thick, poporange] (-0.9,0.75) -- (-0.9,-1.62);
% Bloc métallique dans l'eau
\draw[thick, poppurple, fill=poppurpleL] (0.1,-2.4) rectangle (0.6,-1.7);
% Annotations
\draw[->, popdark] (2.6,-0.6) -- (3.4,-0.6) node[right, align=left]{paroi\\isolante};
\draw[->, popblue] (-1.7,-2.2) -- (-3.1,-2.2) node[left, align=right]{vase\\intérieur};
\draw[->, popblue] (1.65,-1.5) -- (3.4,-1.9) node[right]{eau};
\draw[->, popink] (0.95,1.6) -- (3.4,1.6) node[right]{thermomètre};
\draw[->, poporange] (-1.15,1.9) -- (-3.1,1.9) node[left]{agitateur};
\draw[->, popdark] (2.6,0.57) -- (3.4,0.9) node[right]{couvercle};
\draw[->, poppurple] (0.35,-2.05) -- (-3.1,-0.9) node[left, align=right]{corps\\étudié};
\end{tikzpicture}
\end{center}
\legende{Schéma annoté en coupe d'un calorimètre et de ses accessoires (schéma exigible au Probatoire).}
\end{popfigure}

\begin{attention}[Schéma exigible]
Le schéma annoté du calorimètre est explicitement demandé par le programme. Entraîne-toi à le reproduire rapidement avec les cinq légendes : vase intérieur, paroi isolante, couvercle, thermomètre, agitateur. N'oublie pas que l'agitateur et le thermomètre traversent le couvercle et plongent dans le liquide.
\end{attention}

\courssub{Capacité calorifique du calorimètre et valeur en eau}

Le vase, le thermomètre et l'agitateur participent eux aussi aux échanges thermiques : quand on verse de l'eau chaude dans le calorimètre, une partie de la chaleur sert à réchauffer l'appareil lui-même. On caractérise cette participation par sa \textbf{capacité calorifique} $C_{\text{cal}}$ (en $\mathrm{J\,K^{-1}}$). Plutôt que de donner $C_{\text{cal}}$ directement, on utilise une grandeur équivalente très pratique.

\begin{definition}[Valeur en eau d'un calorimètre]
La \cle{valeur en eau} $\mu$ d'un calorimètre est la masse d'eau qui aurait la même capacité calorifique que le calorimètre et ses accessoires :
\[ C_{\text{cal}} = \mu \cdot c_{\text{eau}} \]
avec $\mu$ en $\mathrm{kg}$ (SI) et $c_{\text{eau}} = \SI{4180}{J\,kg^{-1}\,K^{-1}}$. Si $\mu$ est donnée en $\mathrm{g}$, on utilise $c_{\text{eau}} = \SI{4,18}{J\,g^{-1}\,K^{-1}}$.
\end{definition}

Dire qu'un calorimètre a une valeur en eau $\mu = \SI{40}{g}$ signifie que, thermiquement, tout se passe comme si le calorimètre contenait $\SI{40}{g}$ d'eau supplémentaire. Dans les bilans, on traitera donc le calorimètre comme une masse $\mu$ d'eau à la température initiale du calorimètre.

\courssub{Le principe des échanges de chaleur}

Voici le principe qui commande tous les problèmes de calorimétrie.

\begin{propriete}[Principe des échanges de chaleur]
Lorsque des corps échangent de la chaleur \textbf{uniquement entre eux} (système isolé, par exemple à l'intérieur d'un calorimètre parfait), la somme algébrique des quantités de chaleur échangées est nulle :
\[ Q_1 + Q_2 + \cdots + Q_n = 0 \qquad \text{soit} \qquad \sum_i Q_i = 0 \]
Convention de signe : $Q > 0$ pour une chaleur \textbf{reçue} par un corps, $Q < 0$ pour une chaleur \textbf{cédée}. Autrement dit : la chaleur cédée par les corps chauds est intégralement reçue par les corps froids.
\end{propriete}

\begin{experience}[Démonstration guidée du principe]
Considérons le système $\Sigma$ formé de tous les corps placés dans le calorimètre (eau, corps chaud, accessoires).
\begin{enumerate}
\item Le calorimètre est thermiquement isolé : le système $\Sigma$ n'échange \textbf{aucune chaleur} avec l'extérieur, et aucun travail mécanique n'est mis en jeu ($W=0$).
\item La \textbf{conservation de l'énergie} (premier principe de la thermodynamique) impose alors que l'énergie interne totale du système reste constante : $\Delta U_{\Sigma} = Q + W = 0$.
\item Or l'énergie interne totale est la somme des énergies internes de chaque corps : $\Delta U_{\Sigma} = \Delta U_1 + \Delta U_2 + \cdots + \Delta U_n$.
\item Pour chaque corps, la variation d'énergie interne est précisément la chaleur qu'il a échangée (puisque $W=0$) : $\Delta U_i = Q_i$.
\item On en déduit : $Q_1 + Q_2 + \cdots + Q_n = 0$. \textbf{C.Q.F.D.}
\end{enumerate}
Le principe des échanges de chaleur n'est donc rien d'autre que la \emph{conservation de l'énergie} appliquée à un système isolé.
\end{experience}

\begin{popfigure}
\begin{center}
\begin{tikzpicture}[scale=1.0, every node/.style={font=\small}]
% Bloc chaud
\draw[thick, popink, fill=popink!20] (-4.6,0) rectangle (-2.6,1.4);
\node at (-3.6,0.7) {\textbf{Corps chaud}};
\node[below] at (-3.6,0) {$m_1,\ c_1,\ \theta_1$};
% Eau froide
\draw[thick, popblue, fill=popblueL] (-0.9,0) rectangle (1.1,1.4);
\node at (0.1,0.7) {\textbf{Eau froide}};
\node[below] at (0.1,0) {$m_2,\ c_{\text{eau}},\ \theta_2$};
% Calorimètre
\draw[thick, popgold, fill=popgoldL] (2.6,0) rectangle (4.6,1.4);
\node at (3.6,0.7) {\textbf{Calorimètre}};
\node[below] at (3.6,0) {$\mu,\ \theta_2$};
% Flèches de chaleur
\draw[->, very thick, popink] (-3.6,-0.45) -- (-3.6,-1.5) node[midway, left]{$Q_1<0$ (cédée)};
\draw[->, very thick, popgreen] (0.1,-1.5) -- (0.1,-0.45) node[midway, left]{$Q_2>0$ (reçue)};
\draw[->, very thick, popgreen] (3.6,-1.5) -- (3.6,-0.45) node[midway, right]{$Q_3>0$ (reçue)};
% Bilan
\node[popdark, font=\bfseries] at (0,-2.2) {$Q_1 + Q_2 + Q_3 = 0$ \quad (système isolé, équilibre à $\theta_{eq}$)};
\end{tikzpicture}
\end{center}
\legende{Schéma-bilan des échanges thermiques dans un calorimètre : la chaleur cédée par le corps chaud est reçue par l'eau froide et par le calorimètre.}
\end{popfigure}

\courssub{Méthode générale de résolution d'un problème de calorimétrie}

\begin{methode}[Résoudre un problème de calorimétrie en quatre étapes]
\begin{enumerate}
\item \textbf{Inventaire} : liste tous les corps participant aux échanges (y compris le calorimètre via sa valeur en eau $\mu$), avec pour chacun sa masse, sa chaleur massique et sa température initiale. Identifie les corps chauds (qui cèdent) et les corps froids (qui reçoivent).
\item \textbf{Écriture des quantités de chaleur} : pour chaque corps, écris $Q_i$ avec la bonne formule : $Q = m\,c\,(\theta_{eq} - \theta_{\text{init}})$ sans changement d'état ; $Q = \pm m\,L$ si changement d'état. Garde la convention algébrique : la formule $m c \Delta\theta$ donne automatiquement le bon signe si $\Delta\theta = \theta_{\text{final}} - \theta_{\text{initial}}$.
\item \textbf{Application du principe} : écris $\sum_i Q_i = 0$.
\item \textbf{Résolution} : résous l'équation obtenue par rapport à l'inconnue ($\theta_{eq}$, $c$, $\mu$, $L$ ou une masse), puis vérifie la cohérence du résultat (unité, ordre de grandeur, $\theta_{eq}$ comprise entre les températures initiales).
\end{enumerate}
\end{methode}

\begin{exemplebox}[Mélange d'eau chaude et d'eau froide]
On mélange dans un calorimètre de valeur en eau négligeable $m_1 = \SI{200}{g}$ d'eau à $\theta_1 = \SI{80}{\celsius}$ et $m_2 = \SI{300}{g}$ d'eau à $\theta_2 = \SI{20}{\celsius}$. Le principe des échanges donne :
\[ m_1\,c_{\text{eau}}\,(\theta_{eq} - \theta_1) + m_2\,c_{\text{eau}}\,(\theta_{eq} - \theta_2) = 0 \]
La chaleur massique $c_{\text{eau}}$ se simplifie :
\[ \theta_{eq} = \frac{m_1 \theta_1 + m_2 \theta_2}{m_1 + m_2} = \frac{200 \times 80 + 300 \times 20}{500} = \frac{22000}{500} = \SI{44}{\celsius} \]
La température d'équilibre est $\theta_{eq} = \SI{44}{\celsius}$ : c'est la moyenne des températures pondérée par les masses. Elle est bien comprise entre $\SI{20}{\celsius}$ et $\SI{80}{\celsius}$, plus proche de $\SI{20}{\celsius}$ car l'eau froide est plus abondante.
\end{exemplebox}

\begin{exoresolu}[Détermination de la chaleur massique d'un métal]
Un bloc de cuivre de masse $m_1 = \SI{150}{g}$ est chauffé à $\theta_1 = \SI{95}{\celsius}$ puis plongé dans un calorimètre contenant $m_2 = \SI{250}{g}$ d'eau à $\theta_2 = \SI{18}{\celsius}$. La valeur en eau du calorimètre est $\mu = \SI{30}{g}$. La température d'équilibre est $\theta_{eq} = \SI{22,4}{\celsius}$. Déterminer la chaleur massique $c_1$ du cuivre. On donne $c_{\text{eau}} = \SI{4180}{J\,kg^{-1}\,K^{-1}}$.
\tcblower
\textbf{Étape 1 — Inventaire.} Corps chaud : le cuivre ($m_1 = \SI{0,150}{kg}$, $c_1$ inconnue, $\theta_1 = \SI{95}{\celsius}$). Corps froids : l'eau ($m_2 = \SI{0,250}{kg}$, $\theta_2 = \SI{18}{\celsius}$) et le calorimètre, équivalent à $\mu = \SI{0,030}{kg}$ d'eau à $\theta_2$.

\textbf{Étape 2 — Quantités de chaleur.}
\begin{align*}
Q_1 &= m_1\,c_1\,(\theta_{eq} - \theta_1) = 0,150 \times c_1 \times (22,4 - 95) = -10,89\,c_1 \\
Q_2 &= m_2\,c_{\text{eau}}\,(\theta_{eq} - \theta_2) = 0,250 \times 4180 \times 4,4 = \SI{4598}{J} \\
Q_3 &= \mu\,c_{\text{eau}}\,(\theta_{eq} - \theta_2) = 0,030 \times 4180 \times 4,4 = \SI{551,8}{J}
\end{align*}

\textbf{Étape 3 — Principe des échanges.} $Q_1 + Q_2 + Q_3 = 0$ :
\[ -10,89\,c_1 + 4598 + 551,8 = 0 \]

\textbf{Étape 4 — Résolution.}
\[ c_1 = \frac{4598 + 551,8}{10,89} = \frac{5149,8}{10,89} \approx \SI{473}{J\,kg^{-1}\,K^{-1}} \]
Vérification : la valeur tabulée du cuivre est $\SI{385}{J\,kg^{-1}\,K^{-1}}$ ; l'écart s'explique par les pertes thermiques résiduelles du calorimètre réel et les incertitudes de mesure. L'ordre de grandeur (quelques centaines de $\mathrm{J\,kg^{-1}\,K^{-1}}$ pour un métal) est correct.
\end{exoresolu}

\begin{attention}[Erreur fréquente : oublier la valeur en eau]
Lorsque l'énoncé fournit la valeur en eau $\mu$ du calorimètre, tu \textbf{dois} inclure le terme $Q_{\text{cal}} = \mu\,c_{\text{eau}}\,(\theta_{eq} - \theta_2)$ dans le bilan. L'omettre fausse le résultat : dans l'exercice résolu ci-dessus, négliger $\mu$ aurait donné $c_1 \approx \SI{422}{J\,kg^{-1}\,K^{-1}}$, soit une erreur de plus de $10\,\%$. En revanche, si l'énoncé précise « calorimètre de capacité négligeable » ou « calorimètre parfait », alors $\mu = 0$ et ce terme disparaît. Autre piège : la température initiale du calorimètre est celle de l'eau froide qu'il contient, pas celle du corps chaud.
\end{attention}

\courssub{Les analogues du calorimètre dans l'environnement}

Le principe d'isolation thermique du calorimètre se retrouve dans de nombreux objets de la vie quotidienne :

\begin{itemize}
\item la \textbf{glacière} : parois en polystyrène ou en mousse isolante, elle garde les boissons fraîches en ralentissant l'entrée de chaleur ;
\item la \textbf{bouteille thermos} : c'est littéralement un vase Dewar domestique, à double paroi de verre argenté sous vide ;
\item la \textbf{marmite norvégienne} : caisse isolée dans laquelle on enferme une casserole portée à ébullition ; la cuisson se poursuit lentement sans apport d'énergie ;
\item les \textbf{conteneurs isothermes} de transport des vaccins et des produits sanguins, indispensables dans les campagnes de vaccination au Cameroun.
\end{itemize}

\begin{savaistu}[Des glacières de nos marchés au vase Dewar]
Les glacières des vendeuses de boissons fraîches que l'on croise dans les rues de Douala ou de Yaoundé exploitent exactement le même principe physique que le calorimètre de laboratoire : une enceinte aux parois isolantes qui minimise les échanges de chaleur avec l'extérieur. À l'intérieur, la glace fond lentement en absorbant la chaleur qui traverse malgré tout les parois ($Q = m L_f$), ce qui maintient les boissons au froid pendant des heures. James Dewar, qui a inventé le vase isolant en 1892 pour liquéfier les gaz, n'a jamais breveté son invention : ce sont d'autres qui en ont tiré fortune avec la bouteille thermos !
\end{savaistu}

\begin{exercice}[Température d'équilibre avec calorimètre réel]
Un calorimètre de valeur en eau $\mu = \SI{25}{g}$ contient $m_2 = \SI{200}{g}$ d'eau à $\theta_2 = \SI{15}{\celsius}$. On y verse $m_1 = \SI{150}{g}$ d'eau à $\theta_1 = \SI{70}{\celsius}$. Calculer la température d'équilibre $\theta_{eq}$.
\end{exercice}

\begin{corrige}[Exercice 1]
Inventaire : l'eau chaude cède $Q_1 = m_1 c_{\text{eau}} (\theta_{eq} - \theta_1)$ ; l'eau froide reçoit $Q_2 = m_2 c_{\text{eau}} (\theta_{eq} - \theta_2)$ ; le calorimètre reçoit $Q_3 = \mu c_{\text{eau}} (\theta_{eq} - \theta_2)$. Le principe $\sum Q_i = 0$, après simplification par $c_{\text{eau}}$, donne :
\[ m_1(\theta_{eq} - \theta_1) + (m_2 + \mu)(\theta_{eq} - \theta_2) = 0 \]
\[ \theta_{eq} = \frac{m_1 \theta_1 + (m_2 + \mu)\theta_2}{m_1 + m_2 + \mu} = \frac{150 \times 70 + 225 \times 15}{375} = \frac{10500 + 3375}{375} = \SI{37}{\celsius} \]
Vérification : $\SI{15}{\celsius} < \SI{37}{\celsius} < \SI{70}{\celsius}$, et sans le calorimètre on aurait trouvé $\SI{38,6}{\celsius}$ : le calorimètre « absorbe » une partie de la chaleur, ce qui abaisse légèrement l'équilibre. Cohérent.
\end{corrige}

\begin{aretenir}[L'essentiel de la section]
\begin{itemize}
\item Le calorimètre est une enceinte thermiquement isolée (vase, parois isolantes, couvercle, thermomètre, agitateur) permettant de mesurer des quantités de chaleur.
\item Sa valeur en eau $\mu$ est la masse d'eau de même capacité calorifique que l'appareil : $C_{\text{cal}} = \mu\,c_{\text{eau}}$.
\item Principe des échanges de chaleur : dans un système isolé, $\sum_i Q_i = 0$ ; c'est la conservation de l'énergie.
\item Méthode : inventaire, écriture algébrique de chaque $Q_i$, application de $\sum Q_i = 0$, résolution et vérification.
\item Glacière, thermos, marmite norvégienne et conteneurs isothermes sont les « calorimètres » de la vie courante.
\end{itemize}
\end{aretenir}

\coursec{Synthèse et résolution de problèmes de calorimétrie}

Dans la section précédente, tu as découvert le calorimètre, cette enceinte quasi adiabatique dans laquelle le \cle{principe des échanges de chaleur} s'écrit simplement $\sum Q = 0$ : tout ce que les corps chauds perdent, les corps froids le gagnent. Tu sais aussi calculer une quantité de chaleur sans changement d'état ($Q = m\,c\,\Delta\theta$) et avec changement d'état ($Q = m\,L$). Le moment est venu de rassembler tous ces outils et d'apprendre à résoudre \emph{n'importe quel} problème de calorimétrie, comme on te le demandera au Probatoire. C'est l'objet de cette dernière section : d'abord un bilan des formules, puis une méthode générale en cinq étapes, enfin quatre problèmes types entièrement résolus.

\courssub{Bilan des formules de la calorimétrie}

Toute la calorimétrie de Première tient dans cinq relations. Avant de les utiliser, assure-toi de bien comprendre ce que chacune représente physiquement :

\begin{itemize}
\item $Q = m\,c\,\Delta\theta$ : chaleur échangée par un corps de masse $m$ dont la température varie de $\Delta\theta$, \textbf{sans changement d'état} ;
\item $Q = m\,L$ : chaleur échangée lors d'un \textbf{changement d'état}, à température constante ;
\item $C = m\,c$ : capacité calorifique d'un corps (chaleur nécessaire pour élever sa température de $1\ \mathrm{K}$) ;
\item $C_{\mathrm{cal}} = \mu\,c_{\mathrm{eau}}$ : la valeur en eau $\mu$ remplace le calorimètre par une masse fictive d'eau équivalente ;
\item $\sum Q_i = 0$ : principe des échanges de chaleur dans un système isolé.
\end{itemize}

\begin{center}
\rowcolors{1}{popcream}{white}
\begin{tabularx}{\linewidth}{|l|X|l|}
\hline
\rowcolor{popblueL} \textbf{Grandeur / relation} & \textbf{Signification} & \textbf{Unité SI} \\
\hline
$Q$ & Quantité de chaleur échangée & joule (J) \\
$m$ & Masse du corps & kilogramme (kg) \\
$c$ & Chaleur massique du corps & $\mathrm{J\,kg^{-1}\,K^{-1}}$ \\
$\Delta\theta = \theta_f - \theta_i$ & Variation de température & kelvin (K) ou degré ($^\circ$C) pour un \emph{écart} \\
$C = m\,c$ & Capacité calorifique & $\mathrm{J\,K^{-1}}$ \\
$\mu$ & Valeur en eau du calorimètre & kg (souvent exprimée en g) \\
$L_f$, $L_v$, $L_s$ & Chaleurs latentes (fusion, vaporisation, solidification) & $\mathrm{J\,kg^{-1}}$ \\
\hline
\end{tabularx}
\end{center}

\begin{aretenir}[Valeurs usuelles à connaître]
$c_{\mathrm{eau}} = \SI{4180}{J.kg^{-1}.K^{-1}}$ ; $c_{\mathrm{glace}} \approx \SI{2100}{J.kg^{-1}.K^{-1}}$ ; $L_f(\mathrm{glace}) = \SI{3,34e5}{J.kg^{-1}}$ ; $L_v(\mathrm{eau}) = \SI{2,26e6}{J.kg^{-1}}$. Une variation de température a la \textbf{même valeur} en kelvins et en degrés Celsius : $\Delta\theta = \SI{5}{K} = \SI{5}{\degree C}$. En revanche, une température absolue s'exprime en kelvins : $T(\mathrm{K}) = \theta(^\circ\mathrm{C}) + 273{,}15$.
\end{aretenir}

\begin{attention}[Convention de signe : la clé de tout le chapitre]
On compte $Q$ \textbf{algébriquement} : $Q > 0$ si le corps \emph{reçoit} de la chaleur, $Q < 0$ s'il en \emph{cède}. En écrivant $Q = m\,c\,(\theta_f - \theta_i)$ avec $\theta_f$ la température finale (souvent $\theta_{eq}$), le signe sort tout seul : un corps qui se refroidit a $\theta_f < \theta_i$, donc $Q < 0$, automatiquement. Ne mets \textbf{jamais} de valeur absolue sur $\Delta\theta$ puis un signe « à la main » : c'est la première source d'erreur au Probatoire.
\end{attention}

\courssub{Stratégie générale de résolution : la méthode en cinq étapes}

Face à un problème de calorimétrie, l'élève qui fonce sur les calculs sans organiser ses idées se trompe presque toujours. Adopte systématiquement la démarche suivante, qui transforme un problème intimidant en une suite de petites questions faciles.

\begin{methode}[Résoudre un problème de calorimétrie en 5 étapes]
\begin{enumerate}
\item \textbf{Schéma et inventaire.} Dessine la situation (calorimètre, corps introduits) et liste \emph{tous} les acteurs thermiques : les corps, mais aussi le calorimètre et son eau si elle existe. Note pour chacun : masse, nature (donc $c$ ou $L$), température initiale.
\item \textbf{Bilan des échanges.} Pour chaque acteur, écris la chaleur échangée avec la bonne formule : $m\,c\,(\theta_{eq} - \theta_i)$ si sa température change, $m\,L$ s'il change d'état. Remplace le calorimètre par sa valeur en eau $\mu$.
\item \textbf{Équation.} Écris le principe des échanges : $\sum Q_i = 0$. C'est une équation où l'inconnue est $\theta_{eq}$, $c$, $\mu$ ou une masse, selon l'énoncé.
\item \textbf{Résolution numérique.} Isole l'inconnue littéralement \emph{d'abord}, puis fais l'application numérique avec les unités SI (masses en kg).
\item \textbf{Vérification de cohérence.} Contrôle les unités (analyse dimensionnelle), le signe, et surtout le caractère physique : une température d'équilibre doit être comprise entre la plus froide et la plus chaude des températures initiales ; une masse ou une chaleur massique doit être positive.
\end{enumerate}
\end{methode}

\begin{popfigure}
\begin{center}
\begin{tikzpicture}[node distance=0.55cm,
etape/.style={rectangle, rounded corners=3pt, draw=popblue, fill=popblueL, text width=10.5cm, align=left, inner sep=5pt, font=\small},
fleche/.style={-{Stealth[length=2.5mm]}, thick, poporange}]
\node[etape] (e1) {\textbf{Étape 1 — Schéma et inventaire} : identifier tous les corps (et le calorimètre via $\mu$), relever $m$, $c$, $\theta_i$ de chacun.};
\node[etape, below=of e1] (e2) {\textbf{Étape 2 — Bilan des chaleurs} : $Q_i = m_i c_i(\theta_{eq}-\theta_i)$ ou $Q_i = \pm m_i L$ selon le cas.};
\node[etape, below=of e2] (e3) {\textbf{Étape 3 — Équation} : principe des échanges $\sum Q_i = 0$.};
\node[etape, below=of e3] (e4) {\textbf{Étape 4 — Résolution} : isoler l'inconnue littéralement, puis application numérique en unités SI.};
\node[etape, below=of e4, draw=popgreen, fill=popgreenL] (e5) {\textbf{Étape 5 — Vérification} : unités, signe, et $\theta_{\min} < \theta_{eq} < \theta_{\max}$.};
\draw[fleche] (e1) -- (e2);
\draw[fleche] (e2) -- (e3);
\draw[fleche] (e3) -- (e4);
\draw[fleche] (e4) -- (e5);
\end{tikzpicture}
\end{center}
\legende{Organigramme de la méthode générale de résolution d'un problème de calorimétrie.}
\end{popfigure}

\courssub{Problème type 1 : température d'équilibre d'un mélange avec calorimètre non négligeable}

\textbf{Situation.} Dans un calorimètre de valeur en eau $\mu = \SI{25}{g}$ contenant $m_1 = \SI{300}{g}$ d'eau froide à $\theta_1 = \SI{20}{\celsius}$, on verse $m_2 = \SI{200}{g}$ d'eau chaude à $\theta_2 = \SI{70}{\celsius}$. On cherche la température d'équilibre $\theta_{eq}$.

\textbf{Étapes 1 et 2.} Trois acteurs : l'eau froide, l'eau chaude, le calorimètre (remplacé par $\mu$ d'eau à $\theta_1$). L'eau froide et le calorimètre reçoivent de la chaleur ; l'eau chaude en cède :
\begin{align*}
Q_1 &= m_1\,c_{\mathrm{eau}}\,(\theta_{eq} - \theta_1) \\
Q_{\mathrm{cal}} &= \mu\,c_{\mathrm{eau}}\,(\theta_{eq} - \theta_1) \\
Q_2 &= m_2\,c_{\mathrm{eau}}\,(\theta_{eq} - \theta_2)
\end{align*}

\textbf{Étape 3.} $Q_1 + Q_{\mathrm{cal}} + Q_2 = 0$. Le facteur $c_{\mathrm{eau}}$ se simplifie :
\[ (m_1 + \mu)(\theta_{eq} - \theta_1) + m_2(\theta_{eq} - \theta_2) = 0 \]

\textbf{Étape 4.} On isole $\theta_{eq}$ :
\[ \theta_{eq} = \frac{(m_1 + \mu)\,\theta_1 + m_2\,\theta_2}{m_1 + \mu + m_2} = \frac{(0{,}325)(20) + (0{,}200)(70)}{0{,}325 + 0{,}200} = \frac{6{,}5 + 14}{0{,}525} \approx \SI{39,0}{\celsius} \]

\textbf{Étape 5.} Vérification : $\SI{20}{\celsius} < \SI{39,0}{\celsius} < \SI{70}{\celsius}$ : cohérent. Sans calorimètre on aurait trouvé $\SI{40,0}{\celsius}$ : le calorimètre, qui « boit » un peu de chaleur, abaisse légèrement l'équilibre. C'est physiquement attendu.

\begin{attention}[Le calorimètre n'est jamais gratuit]
Oublier la valeur en eau $\mu$ quand elle est donnée est une erreur éliminatoire en calorimétrie. Dès que l'énoncé mentionne « calorimètre de valeur en eau $\mu$ » ou « de capacité calorifique $C_{\mathrm{cal}}$ », ajoute le terme $\mu\,c_{\mathrm{eau}}(\theta_{eq}-\theta_1)$ ou $C_{\mathrm{cal}}(\theta_{eq}-\theta_1)$ dans le bilan. Le calorimètre est toujours initialement à la température de l'eau froide qu'il contient.
\end{attention}

\courssub{Problème type 2 : mesure de la chaleur massique d'un métal (méthode des mélanges)}

C'est \emph{le} grand classique des TP et du Probatoire : on chauffe un bloc de métal, on le plonge dans l'eau froide d'un calorimètre, et la température d'équilibre permet de remonter à $c_{\mathrm{métal}}$.

\begin{exemplebox}[Chaleur massique du cuivre]
Un bloc de cuivre de masse $m_c = \SI{150}{g}$, porté à $\theta_c = \SI{95}{\celsius}$ dans un bain d'eau bouillante, est plongé rapidement dans un calorimètre de valeur en eau $\mu = \SI{20}{g}$ contenant $m_e = \SI{250}{g}$ d'eau à $\theta_e = \SI{18}{\celsius}$. La température d'équilibre est $\theta_{eq} = \SI{21}{\celsius}$. Déterminons $c_{\mathrm{cuivre}}$.

\textbf{Bilan.} Le cuivre cède : $Q_c = m_c\,c_c\,(\theta_{eq} - \theta_c)$ (négatif). L'eau et le calorimètre reçoivent : $Q_e + Q_{\mathrm{cal}} = (m_e + \mu)\,c_{\mathrm{eau}}\,(\theta_{eq} - \theta_e)$ (positif).

\textbf{Équation et résolution.}
\[ m_c\,c_c\,(\theta_{eq} - \theta_c) + (m_e + \mu)\,c_{\mathrm{eau}}\,(\theta_{eq} - \theta_e) = 0 \]
\[ c_c = \frac{(m_e + \mu)\,c_{\mathrm{eau}}\,(\theta_{eq} - \theta_e)}{m_c\,(\theta_c - \theta_{eq})} \]
Application numérique (masses en kg) :
\[ c_c = \frac{(0{,}270)(4180)(21 - 18)}{0{,}150 \times (95 - 21)} = \frac{3385{,}8}{11{,}1} \approx \SI{305}{J.kg^{-1}.K^{-1}} \]

\textbf{Vérification.} Le résultat est positif, de l'ordre de quelques centaines de $\mathrm{J\,kg^{-1}\,K^{-1}}$ comme tous les métaux, et proche de la valeur tabulée du cuivre ($\SI{385}{J.kg^{-1}.K^{-1}}$). L'écart s'explique par les pertes thermiques lors du transfert rapide du bloc et par l'imprécision des thermomètres : c'est une excellente discussion à rédiger en conclusion de TP.
\end{exemplebox}

\courssub{Problème type 3 : glace dans l'eau chaude — toute la glace fond-elle ?}

Voici le problème le plus subtil du chapitre. Quand on jette de la glace à $\SI{0}{\celsius}$ dans de l'eau chaude, deux étapes se succèdent pour la glace : la \textbf{fusion} (à $\SI{0}{\celsius}$, coût $m_g L_f$) puis le \textbf{réchauffement} de l'eau de fusion. Mais rien ne garantit que l'eau chaude puisse fournir assez de chaleur pour tout fondre ! Il faut donc \emph{tester l'hypothèse} avant de conclure.

\begin{attention}[Erreur fréquente : conclure sans vérifier]
Beaucoup d'élèves écrivent directement $\sum Q = 0$ en supposant toute la glace fondue, trouvent $\theta_{eq}$... et s'arrêtent. Or si le calcul donne $\theta_{eq} < \SI{0}{\celsius}$, l'hypothèse est \textbf{absurde} : un mélange eau-glace ne peut pas descendre sous $\SI{0}{\celsius}$ à pression ambiante. La bonne démarche : comparer la chaleur \emph{disponible} (ce que l'eau chaude peut céder au maximum en descendant à $\SI{0}{\celsius}$) à la chaleur \emph{nécessaire} pour fondre toute la glace. Si la disponible est insuffisante, l'état final est un mélange eau + glace à $\theta_{eq} = \SI{0}{\celsius}$, et on calcule la masse de glace réellement fondue.
\end{attention}

\begin{exoresolu}[Glace dans l'eau chaude : hypothèse à tester]
\textbf{Énoncé.} Dans un calorimètre de capacité négligeable contenant $m_1 = \SI{200}{g}$ d'eau à $\theta_1 = \SI{40}{\celsius}$, on introduit $m_g = \SI{50}{g}$ de glace à $\SI{0}{\celsius}$. Déterminer l'état final du système. On donne $c_{\mathrm{eau}} = \SI{4180}{J.kg^{-1}.K^{-1}}$ et $L_f = \SI{3,34e5}{J.kg^{-1}}$.
\tcblower
\textbf{Solution.} \textbf{Étape 1.} Acteurs : l'eau chaude ($m_1$, $\theta_1$) et la glace ($m_g$, $\SI{0}{\celsius}$).

\textbf{Étape 2 (test de l'hypothèse).} Chaleur maximale que l'eau chaude peut céder en se refroidissant jusqu'à $\SI{0}{\celsius}$ :
\[ Q_{\mathrm{dispo}} = m_1\,c_{\mathrm{eau}}\,(\theta_1 - 0) = 0{,}200 \times 4180 \times 40 = \SI{33440}{J} \]
Chaleur nécessaire pour fondre toute la glace :
\[ Q_{\mathrm{néc}} = m_g\,L_f = 0{,}050 \times 3{,}34\times 10^{5} = \SI{16700}{J} \]
Comme $Q_{\mathrm{dispo}} > Q_{\mathrm{néc}}$ ($\SI{33440}{J} > \SI{16700}{J}$), \textbf{toute la glace fond} et l'équilibre se situe au-dessus de $\SI{0}{\celsius}$.

\textbf{Étape 3.} Bilan complet : l'eau chaude cède $Q_1 = m_1 c_{\mathrm{eau}}(\theta_{eq} - \theta_1)$ ; la glace reçoit $Q_g = m_g L_f + m_g c_{\mathrm{eau}}(\theta_{eq} - 0)$ (fusion \emph{puis} réchauffement de l'eau de fusion).

\textbf{Étape 4.} $Q_1 + Q_g = 0$ :
\[ \theta_{eq} = \frac{m_1 c_{\mathrm{eau}} \theta_1 - m_g L_f}{(m_1 + m_g)\,c_{\mathrm{eau}}} = \frac{33440 - 16700}{0{,}250 \times 4180} = \frac{16740}{1045} \approx \SI{16,0}{\celsius} \]

\textbf{Étape 5.} Vérification : $0 < \SI{16,0}{\celsius} < \SI{40}{\celsius}$, positif et supérieur à $\SI{0}{\celsius}$ : l'hypothèse « toute la glace fond » est confirmée a posteriori. Le résultat est cohérent.

\textbf{Variante.} Avec seulement $m_1 = \SI{80}{g}$ d'eau à $\SI{40}{\celsius}$, on aurait $Q_{\mathrm{dispo}} = 0{,}080 \times 4180 \times 40 = \SI{13376}{J} < Q_{\mathrm{néc}} = \SI{16700}{J}$ : la glace ne fond que partiellement. L'état final est un mélange à $\theta_{eq} = \SI{0}{\celsius}$ et la masse fondue vaut $m_f = \dfrac{Q_{\mathrm{dispo}}}{L_f} = \dfrac{13376}{3{,}34\times 10^{5}} \approx \SI{40,0}{g}$. Il reste $\SI{10,0}{g}$ de glace.
\end{exoresolu}

\courssub{Problème type 4 : détermination de la valeur en eau d'un calorimètre}

Avant toute mesure précise, il faut \emph{étalonner} le calorimètre, c'est-à-dire mesurer sa valeur en eau $\mu$. Le principe : on y mélange deux masses d'eau connues à températures connues, et on déduit $\mu$ de la température d'équilibre.

\textbf{Situation.} Un calorimètre contient $m_1 = \SI{250}{g}$ d'eau à $\theta_1 = \SI{17}{\celsius}$. On y verse $m_2 = \SI{150}{g}$ d'eau à $\theta_2 = \SI{60}{\celsius}$ ; l'équilibre s'établit à $\theta_{eq} = \SI{32}{\celsius}$. Cherchons $\mu$.

\textbf{Bilan et équation.} L'eau chaude cède, l'eau froide et le calorimètre reçoivent :
\[ m_2\,c_{\mathrm{eau}}\,(\theta_{eq} - \theta_2) + (m_1 + \mu)\,c_{\mathrm{eau}}\,(\theta_{eq} - \theta_1) = 0 \]
Le facteur $c_{\mathrm{eau}}$ se simplifie ; on isole $\mu$ :
\[ \mu = \frac{m_2\,(\theta_2 - \theta_{eq})}{\theta_{eq} - \theta_1} - m_1 = \frac{150 \times (60 - 32)}{32 - 17} - 250 = \frac{4200}{15} - 250 = 280 - 250 = \SI{30}{g} \]

\textbf{Vérification.} $\mu > 0$, de l'ordre de quelques dizaines de grammes, ce qui est typique d'un vase Dewar de laboratoire scolaire. Remarque : si le calorimètre était parfait ($\mu = 0$), on aurait $\theta_{eq} = \dfrac{m_1\theta_1 + m_2\theta_2}{m_1+m_2} = \SI{33,1}{\celsius}$ ; l'écart mesuré trahit la participation thermique du calorimètre.

\courssub{Vérification systématique : le réflexe du physicien}

Un résultat non vérifié est un résultat suspect. Voici la check-list à passer mentalement avant de rendre ta copie :

\begin{enumerate}
\item \textbf{Unités.} Les masses sont-elles en kg ? Les chaleurs en J ? Fais l'analyse dimensionnelle de la formule littérale : dans $c_c = \dfrac{(m_e+\mu)c_{\mathrm{eau}}\Delta\theta}{m_c \Delta\theta'}$, les masses et les écarts de température se simplifient partiellement et il reste bien des $\mathrm{J\,kg^{-1}\,K^{-1}}$.
\item \textbf{Signe.} Une chaleur massique, une valeur en eau, une chaleur latente, une masse : toujours positives. Un résultat négatif signale une erreur de signe dans le bilan.
\item \textbf{Encadrement.} Sans changement d'état, $\theta_{\min} < \theta_{eq} < \theta_{\max}$. Avec de la glace à $\SI{0}{\celsius}$, $\theta_{eq} \geq \SI{0}{\celsius}$ ; avec de la vapeur à $\SI{100}{\celsius}$, $\theta_{eq} \leq \SI{100}{\celsius}$.
\item \textbf{Ordre de grandeur.} Fondre $\SI{1}{g}$ de glace coûte environ $\SI{334}{J}$, soit autant que chauffer $\SI{1}{g}$ d'eau de $\SI{80}{K}$ : la fusion « coûte très cher ». Si ton calcul fait fondre un kilo de glace avec une cuillère d'eau tiède, méfiance.
\end{enumerate}

\begin{savaistu}[La calorimétrie mesure aussi l'énergie des aliments]
Les « calories » affichées sur les emballages (en réalité des kilocalories, $\SI{1}{kcal} = \SI{4180}{J}$) sont mesurées avec une \textbf{bombe calorimétrique} : on brûle complètement un échantillon d'aliment dans une enceinte remplie d'oxygène, placée dans un calorimètre d'eau, et l'élévation de température donne l'énergie libérée. Ainsi, $\SI{100}{g}$ d'arachides grillées — très appréciées au Cameroun — libèrent environ $\SI{570}{kcal} \approx \SI{2,4e6}{J}$ : de quoi porter près de $7$ litres d'eau de $\SI{20}{\celsius}$ à ébullition ! La calorimétrie est aussi utilisée en médecine (métabolisme de base), dans l'industrie (pouvoir calorifique des carburants) et au laboratoire pour mesurer des chaleurs de réaction chimique.
\end{savaistu}

\begin{exercice}[Mélange dans un calorimètre]
Un calorimètre de valeur en eau $\mu = \SI{15}{g}$ contient $m_1 = \SI{400}{g}$ d'eau à $\theta_1 = \SI{15}{\celsius}$. On y plonge un bloc d'aluminium de masse $m_2 = \SI{200}{g}$ sortant d'une étuve à $\theta_2 = \SI{80}{\celsius}$. On donne $c_{\mathrm{al}} = \SI{900}{J.kg^{-1}.K^{-1}}$. Calculer la température d'équilibre.
\end{exercice}

\begin{corrige}[Exercice 1]
Bilan : $m_2 c_{\mathrm{al}}(\theta_{eq}-\theta_2) + (m_1+\mu)c_{\mathrm{eau}}(\theta_{eq}-\theta_1) = 0$, d'où
\[ \theta_{eq} = \frac{m_2 c_{\mathrm{al}}\theta_2 + (m_1+\mu)c_{\mathrm{eau}}\theta_1}{m_2 c_{\mathrm{al}} + (m_1+\mu)c_{\mathrm{eau}}} = \frac{0{,}2\times 900\times 80 + 0{,}415\times 4180\times 15}{0{,}2\times 900 + 0{,}415\times 4180} \]
\[ \theta_{eq} = \frac{14400 + 26020{,}5}{180 + 1734{,}7} = \frac{40420{,}5}{1914{,}7} \approx \SI{21,1}{\celsius} \]
Vérification : $\SI{15}{\celsius} < \SI{21,1}{\celsius} < \SI{80}{\celsius}$ ; l'aluminium, de faible capacité calorifique ($180\ \mathrm{J\,K^{-1}}$ contre $1735\ \mathrm{J\,K^{-1}}$ pour l'eau), ne réchauffe que peu le mélange : cohérent.
\end{corrige}

\begin{exercice}[Glace et eau chaude]
Dans un calorimètre de capacité négligeable contenant $\SI{150}{g}$ d'eau à $\SI{25}{\celsius}$, on dépose $\SI{60}{g}$ de glace à $\SI{0}{\celsius}$. Déterminer l'état final (température d'équilibre et masses éventuelles de glace restante).
\end{exercice}

\begin{corrige}[Exercice 2]
Chaleur disponible : $Q_{\mathrm{dispo}} = 0{,}150 \times 4180 \times 25 = \SI{15675}{J}$. Chaleur nécessaire pour tout fondre : $Q_{\mathrm{néc}} = 0{,}060 \times 3{,}34\times 10^{5} = \SI{20040}{J}$. Comme $Q_{\mathrm{dispo}} < Q_{\mathrm{néc}}$, la glace ne fond que partiellement : $\theta_{eq} = \SI{0}{\celsius}$ et la masse fondue est $m_f = \dfrac{15675}{3{,}34\times 10^{5}} \approx \SI{46,9}{g}$. Il reste environ $\SI{13,1}{g}$ de glace dans $\SI{196,9}{g}$ d'eau, le tout à $\SI{0}{\celsius}$.
\end{corrige}

\begin{aretenir}[L'essentiel de la leçon]
La calorimétrie repose sur un unique principe — dans un système isolé, $\sum Q = 0$ — et deux formules : $Q = m\,c\,\Delta\theta$ (température qui change) et $Q = m\,L$ (état qui change). La méthode en cinq étapes (schéma, bilan, équation, résolution, vérification) résout tous les problèmes. Face à un changement d'état, \textbf{toujours tester l'hypothèse} en comparant chaleur disponible et chaleur nécessaire avant de conclure. Et n'oublie jamais le calorimètre : sa valeur en eau $\mu$ participe aux échanges comme une masse d'eau supplémentaire.
\end{aretenir}

\newpage

\coursec{Activité d'intégration}

\coursec{Chaleur, température et échanges thermiques}

\courssub{Notions fondamentales : température et chaleur}
\begin{experience}[Température et dilatation]
On plonge un ballon muni d'un tube capillaire dans l'eau chaude : le liquide monte. Dans l'eau froide : il descend. La \cle{température} $\theta$ (en \si{\celsius} ou K) est une mesure de l'agitation thermique microscopique des particules. Elle se mesure avec un thermomètre (dilatation liquide, thermocouple, résistance Pt100, pyromètre).
\end{experience}

\begin{definition}[Chaleur]
La \cle{chaleur} $Q$ (en joules, J) est un \emph{transfert d'énergie} entre deux systèmes due à une différence de température. Elle n'est pas une propriété d'état : on ne dit pas « un corps possède de la chaleur », mais « un corps a une énergie interne $U$ » ; la chaleur est l'énergie en transit.
\end{definition}

\begin{propriete}[Équilibre thermique]
Deux corps mis en contact évoluent vers une température commune $\theta_f$ : c'est l'\cle{équilibre thermique}. Le principe zéro de la thermodynamique assure la transitivité : si A est en équilibre avec B et B avec C, alors A est en équilibre avec C.
\end{propriete}

\courssub{Énergie interne et premier principe (rappel)}
Pour un système fermé au repos, la variation d'énergie interne $\Delta U$ est la somme du travail $W$ reçu et de la chaleur $Q$ reçue :
\[ \Delta U = Q + W \qquad \text{(signe convention chimie/physique : reçu $>0$)} \]
En calorimétrie, on néglige souvent le travail ($W=0$), d'où $\Delta U = Q$.

\begin{aretenir}[Vocabulaire du Probatoire]
\begin{itemize}
\item \textbf{Température $\theta$} : grandeur intensive, unité SI : kelvin (K) ; en pratique \si{\celsius} ($\theta/\si{\celsius} = T/\si{K} - 273,15$).
\item \textbf{Chaleur $Q$} : énergie transférée, unité SI : joule (J). $Q>0$ : chaleur reçue (réchauffement) ; $Q<0$ : chaleur cédée (refroidissement).
\item \textbf{Système calorifugé} : système qui n'échange pas de chaleur avec l'extérieur ($Q_{\text{ext}}=0$).
\end{itemize}
\end{aretenir}

\coursec{Les trois modes de transfert de chaleur}

\begin{propriete}[Conduction]
Transfert par contact direct, sans déplacement macroscopique de la matière. Dominant dans les \cle{solides} (phonons, électrons libres). Loi de Fourier (1D, régime permanent) : $\Phi = -k S \dfrac{\Delta\theta}{\Delta x}$ où $\Phi$ (W) est le flux thermique, $k$ (W$\cdot$m$^{-1}$$\cdot$K$^{-1}$) la conductivité thermique.
\begin{itemize}
\item Exemple camerounais : la poignée en métal d'une casserole sur le feu (conduction rapide) vs manche en bois/plastique (isolant, $k$ faible).
\item Parois de la glacière : mousse polyuréthane ($k \approx 0,025$ W$\cdot$m$^{-1}$$\cdot$K$^{-1}$) ou liège pour freiner la conduction.
\end{itemize}
\end{propriete}

\begin{propriete}[Convection]
Transfert dans un \cle{fluide} (liquide ou gaz) par déplacement de matière. \textbf{Convection naturelle} : différences de densité (Archimède) ; \textbf{convection forcée} : ventilateur, pompe.
\begin{itemize}
\item Exemple : air chaud qui monte au-dessus de l'étal de Mama Élise ; circulation d'eau dans le radiateur d'un groupe électrogène ENEO.
\item Dans la glacière : ouverture du couvercle $\rightarrow$ entrée d'air chaud (convection forcée) ; joints défectueux $\rightarrow$ fuites convectives.
\end{itemize}
\end{propriete}

\begin{propriete}[Rayonnement thermique]
Transfert par \cle{ondes électromagnétiques} (infrarouge principalement), sans support matériel. Tout corps à $T>0$ K émet. Loi de Stefan-Boltzmann : $\Phi = \varepsilon \sigma S T^4$ ($\sigma = 5,67\times10^{-8}$ W$\cdot$m$^{-2}$$\cdot$K$^{-4}$, $\varepsilon$ émissivité).
\begin{itemize}
\item Exemple : Soleil $\rightarrow$ couvercle de la glacière (rayonnement solaire $\approx 1000$ W/m$^2$ au zénith) ; four à micro-ondes (rayonnement $\approx 12$ cm) ; imagerie thermique médicale.
\item Surface claire/blanche : $\varepsilon$ faible (réfléchit) ; surface noire/mate : $\varepsilon \approx 1$ (absorbe/émet).
\end{itemize}
\end{propriete}

\begin{savaistu}[Le pot en terre cuite « jarre » : climatiseur naturel]
Au Nord-Cameroun, l'eau stockée dans une jarre en terre cuite poreuse reste fraîche. L'eau s'infiltre dans les pores et s'évapore à la surface externe : l'évaporation puise l'énergie latente de vaporisation ($L_v \approx 2,26\times10^6$ J/kg) dans la jarre et l'eau qu'elle contient $\rightarrow$ refroidissement par \emph{changement d'état} couplé à la convection/rayonnement. C'est un ancêtre physique du réfrigérateur sans électricité !
\end{savaistu}

\begin{exemplebox}[Identification des modes pour la glacière Mokolo]
\begin{center}
\begin{tabularx}{\linewidth}{|l|X|X|}\hline
\rowcolor{popblueL} Mode & Trajet physique & Action sur la glacière \\ \hline
\textbf{Rayonnement} & Soleil $\rightarrow$ couvercle/parois exposées & Échauffement direct des parois externes, même sans contact \\ \hline
\textbf{Conduction} & Parois externes $\rightarrow$ isolant $\rightarrow$ vase intérieur ; Étal bois $\rightarrow$ fond glacière & Traversée de la matière solide ; freinée par l'isolant (mousse) \\ \hline
\textbf{Convection} & Air ambiant chaud $\rightarrow$ parois externes ; Ouverture couvercle $\rightarrow$ air chaud entrant & Échange fluide/paroi ; limité par joints étanches et ouvertures rares \\ \hline
\end{tabularx}
\end{center}
\end{exemplebox}

\coursec{Quantité de chaleur sans changement d'état}

\begin{experience}[Mesure de la capacité thermique massique]
On mélange une masse $m_1$ d'eau chaude ($\theta_1$) et $m_2$ d'eau froide ($\theta_2$) dans un calorimètre. À l'équilibre $\theta_f$ : $m_1 c_e (\theta_f-\theta_1) + m_2 c_e (\theta_f-\theta_2) + \mu c_e (\theta_f-\theta_2) = 0$. On en déduit $c_e \approx 4185$ J$\cdot$kg$^{-1}$$\cdot$K$^{-1}$.
\end{experience}

\begin{definition}[Capacité thermique massique (chaleur massique) $c$]
Quantité de chaleur nécessaire pour élever de 1 K la température de 1 kg de matière, \textbf{sans changement d'état} :
\[ c = \frac{1}{m} \frac{\delta Q}{\mathrm{d}\theta} \quad \Rightarrow \quad Q = m \int_{\theta_i}^{\theta_f} c(\theta)\, \mathrm{d}\theta \]
Si $c$ constant sur l'intervalle : $\boxed{Q = m\,c\,\Delta\theta}$ avec $\Delta\theta = \theta_f - \theta_i$ (en K ou \si{\celsius}, même valeur pour un écart).
\end{definition}

\begin{definition}[Capacité calorifique $C$ et valeur en eau $\mu$]
$C = m\,c$ (J/K) : chaleur pour élever de 1 K le corps entier. \cle{Valeur en eau} $\mu = \dfrac{C}{c_e}$ (kg) : masse d'eau ayant même capacité thermique que le corps (calorimètre + accessoires). Dans les bilans, on remplace le calorimètre par une masse $\mu$ d'eau.
\end{definition}

\begin{aretenir}[Tableau de quelques $c$ (J$\cdot$kg$^{-1}$$\cdot$K$^{-1}$) à 20°C]
\begin{center}
\begin{tabular}{|l|c|c|c|c|c|}\hline
Eau & Alu & Fer & Cuivre & Verre & Air (sec) \\ \hline
4185 & 900 & 460 & 385 & 750 & 1005 \\ \hline
\end{tabular}
\end{center}
\textbf{Analyse dimensionnelle} : $[c] = \frac{[Q]}{[m][\Delta\theta]} = \frac{\si{J}}{\si{kg}\cdot\si{K}} = \si{J.kg^{-1}.K^{-1}}$. Vérifié.
\end{aretenir}

\coursec{Chaleur latente et changements d'état}

\begin{experience}[Courbe de chauffage de la glace (TP classique)]
On chauffe de la glace à $-10^\circ$C à puissance constante $P$ et on relève $\theta(t)$. On observe deux paliers :
\begin{enumerate}
\item $\theta = 0^\circ$C : \textbf{fusion} (solide $\rightarrow$ liquide). Température constante, énergie fournie $=$ chaleur latente de fusion $L_f$.
\item $\theta = 100^\circ$C : \textbf{vaporisation} (liquide $\rightarrow$ gaz). Température constante, énergie fournie $=$ chaleur latente de vaporisation $L_v$.
\end{enumerate}
\end{experience}

\begin{definition}[Chaleur latente $L$]
Quantité de chaleur absorbée (ou libérée) lors d'un changement d'état \textbf{à température constante} pour une masse $m$ :
\[ \boxed{Q = m\,L} \qquad L_f \text{ (fusion)},\; L_v \text{ (vaporisation)},\; L_s \text{ (sublimation)} \]
Unités : J/kg. Pour l'eau : $L_f = 3,34\times10^5$ J/kg ; $L_v = 2,26\times10^6$ J/kg (à 100°C, 1 atm).
\end{definition}

\begin{propriete}[Interprétation microscopique]
Lors de la fusion, l'énergie fournie brise les liaisons hydrogène du réseau cristallin (augmentation de l'énergie potentielle interne), sans augmenter l'agitation cinétique (température). $L_v \gg L_f$ car la vaporisation sépare complètement les molécules (gaz parfait).
\end{propriete}

\begin{attention}[Pièges classiques sur les changements d'état]
\begin{itemize}
\item Confondre $c$ (chaleur massique) et $L$ (chaleur latente) : $c$ pour $\Delta\theta \neq 0$, $L$ pour $\Delta\theta = 0$ (palier).
\item Oublier que l'eau issue de la fusion continue d'échanger de la chaleur : elle se réchauffe de $0^\circ$C à $\theta_f$ $\rightarrow$ terme $m_g c_e \theta_f$ indispensable.
\item Unités : $L$ en J/kg, $m$ en kg $\rightarrow$ $Q$ en J. Ne pas mélanger g et kg.
\end{itemize}
\end{attention}

\coursec{Le calorimètre et le principe des échanges de chaleur}

\begin{experience}[Le calorimètre à mélange (appareil de cours)]
\begin{popfigure}
\begin{tikzpicture}[scale=0.9]
% Vase extérieur (isolant)
\draw[fill=poporange!20, thick] (0,0) rectangle (4,5);
% Vase intérieur (métal)
\draw[fill=popblue!10, thick] (0.5,0.5) rectangle (3.5,4.5);
% Eau
\draw[fill=popblue!30] (0.5,0.5) rectangle (3.5,3);
% Agitateur
\draw[thick] (2,0.5) -- (2,4.2);
\draw[fill=popdark] (1.8,4.2) rectangle (2.2,4.5);
% Thermomètre
\draw[thick] (3.2,0.5) -- (3.2,4.7);
\draw[fill=poppink!60] (3.2,0.5) circle (0.15);
% Couvercle
\draw[fill=popgreen!30, thick] (0.3,4.5) rectangle (3.7,4.9);
% Annotations
\node[font=\small, align=left] at (5,4.2) {Couvercle (limite convection)};
\draw[->, popmuted] (3.7,4.7) -- (4.5,4.6);
\node[font=\small, align=left] at (5,3) {Vase intérieur (cuivre/aluminium)};
\draw[->, popmuted] (3.5,3) -- (4.5,3);
\node[font=\small, align=left] at (5,1.5) {Eau / mélange ($m$, $c$)};
\draw[->, popmuted] (3.5,1.5) -- (4.5,1.5);
\node[font=\small, align=left] at (5,0.2) {Isolant (liège, mousse, vide)};
\draw[->, popmuted] (0.5,0.2) -- (-0.5,0.2);
\node[font=\small] at (2,5.4) {Thermomètre};
\node[font=\small] at (2,4.8) {Agitateur};
\end{tikzpicture}
\legende{Calorimètre à mélange : isolant, vase conducteur, agitateur (uniformise $\theta$), thermomètre, couvercle.}
\end{popfigure}
\end{experience}

\begin{propriete}[Principe des échanges de chaleur (système calorifugé)]
Dans un système \textbf{calorifugé} (aucun échange thermique avec l'extérieur), la somme algébrique des quantités de chaleur reçues par tous les sous-systèmes est nulle :
\[ \boxed{\sum_i Q_i = 0} \qquad \text{avec la convention : } Q_i > 0 \text{ (reçu), } Q_i < 0 \text{ (cédé).} \]
C'est une conséquence du premier principe ($\Delta U = Q + W$, $W=0$, $\Delta U_{\text{total}}=0$).
\end{propriete}

\begin{methode}[Résolution d'un problème de calorimétrie]
\begin{enumerate}
\item \textbf{Lister} les corps en présence : masses $m_i$, états initiaux ($\theta_i$, phase), valeur en eau $\mu$ du calorimètre.
\item \textbf{Identifier} l'état final commun : température $\theta_f$, phases finales (tout liquide ? mélange glace/eau ?).
\item \textbf{Écrire} chaque terme $Q_i$ :
   \begin{itemize}
   \item Sans changement d'état : $Q_i = m_i c_i (\theta_f - \theta_i)$.
   \item Avec changement d'état (fusion à $0^\circ$C) : $Q_{\text{fus}} = m_g L_f$ (absorbé $>0$) + $Q_{\text{réchauffement eau}} = m_g c_e (\theta_f - 0)$.
   \item Calorimètre : $Q_{\text{cal}} = \mu c_e (\theta_f - \theta_i)$.
   \item Apport extérieur (si non calorifugé) : $Q_{\text{ext}}$ (reçu $>0$).
   \end{itemize}
\item \textbf{Écrire le bilan} : $\sum Q_i = 0$ (si calorifugé) ou $\sum Q_{\text{internes}} = -Q_{\text{ext}}$ (si apports extérieurs).
\item \textbf{Résoudre} l'inconnue (masse, $\theta_f$, $c$, $L$...).
\item \textbf{Vérifier} : homogénéité dimensionnelle, ordre de grandeur, signe physique ($\theta_f$ entre $\theta_{\min}$ et $\theta_{\max}$ initiaux).
\end{enumerate}
\end{methode}

\coursec{Synthèse et résolution de problèmes : La glacière du marché Mokolo}

\courssub{Objectifs de l'activité}
À l'issue de cette activité d'intégration, tu seras capable de :
\begin{itemize}
\item identifier, dans une situation réelle de conservation du froid, les phénomènes qui mettent en jeu des échanges thermiques ;
\item dire, pour chaque échange, le mode de transfert de chaleur mis en jeu (conduction, convection, rayonnement) ;
\item réaliser le schéma annoté d'une glacière assimilée à un calorimètre, avec ses accessoires ;
\item calculer une quantité de chaleur échangée sans changement d'état ($Q = m\,c\,\Delta\theta$) et lors d'un changement d'état ($Q = m\,L_f$) ;
\item appliquer le bilan énergétique (premier principe) dans un système non calorifugé en tenant compte de la valeur en eau du calorimètre ;
\item proposer et justifier des améliorations technologiques simples à partir de l'analyse physique des transferts thermiques.
\end{itemize}

\courssub{Situation : La glacière du marché Mokolo}
Mama Élise vend des boissons fraîches (jus de foléré, bissap, eaux minérales) au marché Mokolo à Yaoundé. Chaque matin à 6 h, elle remplit sa glacière de \SI{5}{kg} de boissons prises à la température ambiante et ajoute de la glace pilée sortant du congélateur. Elle constate qu'en fin de journée, vers 17 h, la glace a entièrement fondu et que les boissons sont redevenues tièdes : les clients se plaignent. Elle se demande : \emph{quelle masse de glace faudrait-il mettre le matin pour que les boissons restent fraîches toute la journée ?}

On modélise la situation de la façon suivante :
\begin{itemize}
\item la glacière est assimilée à un \cle{calorimètre} imparfait : ses parois participent aux échanges thermiques et sa \cle{valeur en eau} vaut $\mu = \SI{300}{g}$ ;
\item les boissons, assimilées à de l'eau, ont une masse totale $m_b = \SI{5}{kg}$ ; on veut les maintenir à la température finale $\theta_f = \SI{10}{\celsius}$ ;
\item la température initiale des boissons et des parois de la glacière, le matin, est $\theta_i = \SI{30}{\celsius}$ (chaleur du matin à Yaoundé) ;
\item la glace est introduite à $\theta_g = \SI{0}{\celsius}$ ; elle fond en absorbant de la chaleur, et l'eau de fusion obtenue reste à \SI{0}{\celsius} tant qu'il reste de la glace (mélange glace--eau en équilibre à \SI{0}{\celsius}) ;
\item malgré l'isolation, la glacière reçoit de l'extérieur, par ses parois, une chaleur totale estimée à $Q_{\text{ext}} = \num{2,5e5}\ \si{J}$ sur la journée ;
\item on considère que la température finale d'équilibre de l'ensemble \{boissons + parois + eau de fusion\} est $\theta_f = \SI{10}{\celsius}$.
\end{itemize}

\textbf{Données numériques :}
\begin{center}
\begin{tabularx}{\linewidth}{|l|X|}\hline
\rowcolor{popblueL} Grandeur & Valeur \\ \hline
Capacité thermique massique de l'eau (et des boissons) & $c_e = \SI{4185}{J.kg^{-1}.K^{-1}}$ \\ \hline
Chaleur latente de fusion de la glace à \SI{0}{\celsius} & $L_f = \SI{3,34e5}{J.kg^{-1}}$ \\ \hline
Valeur en eau de la glacière & $\mu = \SI{0,300}{kg}$ \\ \hline
Chaleur reçue de l'extérieur sur la journée & $Q_{\text{ext}} = \num{2,5e5}\ \si{J}$ \\ \hline
\end{tabularx}
\end{center}

\courssub{Consignes}
Travail en groupes de quatre élèves. Chaque groupe rédige un rapport et un rapporteur présente les résultats au tableau.

\begin{enumerate}
\item \textbf{Schéma.} Réalise le schéma annoté en coupe de la glacière assimilée à un calorimètre : tu feras apparaître le vase intérieur, les parois isolantes, le couvercle, les boissons, la glace et le thermomètre. Indique par des flèches colorées les trois transferts thermiques (rayonnement, conduction, convection).
\item \textbf{Modes de transfert.} La glacière est posée au soleil sur un étal en bois. Identifie les trois modes de transfert de chaleur par lesquels la chaleur extérieure peut atteindre l'intérieur de la glacière, et précise pour chacun le trajet suivi.
\item \textbf{Masse de glace minimale (fusion seule).} Dans un premier raisonnement simplifié, on suppose que la glace sert uniquement à absorber la chaleur $Q_{\text{ext}}$ venant de l'extérieur par sa fusion à \SI{0}{\celsius} (on néglige le refroidissement des boissons et le réchauffement de l'eau de fusion). Calcule la masse de glace $m_1$ nécessaire.
\item \textbf{Bilan complet.} En réalité, la glace doit aussi refroidir les boissons et les parois de la glacière de \SI{30}{\celsius} à \SI{10}{\celsius}, et l'eau de fusion se réchauffe de \SI{0}{\celsius} à \SI{10}{\celsius}. Écris le bilan énergétique complet en distinguant chaque terme, puis calcule la masse de glace totale $m_g$ à mettre le matin.
\item \textbf{Comparaison et conclusion.} Compare $m_1$ et $m_g$. Explique pourquoi la glace de Mama Élise fond avant la fin de la journée si elle n'a prévu que la masse $m_1$.
\item \textbf{Améliorations.} Propose deux améliorations concrètes et peu coûteuses de la glacière ou de son utilisation, en justifiant chacune par le mode de transfert de chaleur qu'elle permet de réduire.
\end{enumerate}

\courssub{Ressources à mobiliser}

\begin{aretenir}[Boîte à outils de l'activité]
\begin{itemize}
\item Quantité de chaleur échangée par un corps de masse $m$ dont la température varie de $\Delta\theta$, \textbf{sans changement d'état} :
\[ Q = m\,c\,\Delta\theta \qquad (Q \text{ en J},\; m \text{ en kg},\; c \text{ en }\si{J.kg^{-1}.K^{-1}},\; \Delta\theta \text{ en K ou }\si{\celsius}) \]
\item Quantité de chaleur absorbée lors de la \textbf{fusion} d'une masse $m$ de glace à \SI{0}{\celsius} :
\[ Q = m\,L_f \]
\item \textbf{Valeur en eau} $\mu$ du calorimètre : masse d'eau ayant la même capacité thermique que le calorimètre et ses accessoires ; dans les bilans, on remplace le calorimètre par une masse $\mu$ d'eau.
\item \textbf{Bilan énergétique (1er principe)} pour un système non calorifugé recevant une chaleur $Q_{\text{ext}}$ de l'extérieur :
\[ \sum_{\text{corps internes}} Q_i = Q_{\text{ext}} \quad \text{ou équivalentement} \quad \text{Chaleurs gagnées} = \text{Chaleurs perdues} + Q_{\text{ext}} \]
\item Les trois modes de transfert : \cle{conduction} (contact, solides), \cle{convection} (mouvements de fluides), \cle{rayonnement} (ondes EM, sans support).
\end{itemize}
\end{aretenir}

\courssub{Démarche et corrigé commenté}

\begin{exoresolu}[La glacière du marché Mokolo]
\textbf{1.} Schéma annoté avec flèches de transfert. \quad \textbf{2.} Identifier les modes de transfert. \quad \textbf{3.} Calculer $m_1$ (fusion seule pour $Q_{\text{ext}}$). \quad \textbf{4.} Bilan complet et calcul de $m_g$. \quad \textbf{5.} Comparer et conclure. \quad \textbf{6.} Proposer deux améliorations justifiées.

\tcblower

\textbf{1. Schéma annoté de la glacière (coupe) avec transferts thermiques.}

\begin{popfigure}
\begin{tikzpicture}[scale=1.0, every node/.style={font=\small}]
% Paroi extérieure
\draw[fill=popblue!15, thick] (0,0) rectangle (6,3.4);
% Isolant
\draw[fill=poporange!25] (0.35,0.35) rectangle (5.65,3.05);
% Vase intérieur
\draw[fill=popcream, thick] (0.65,0.65) rectangle (5.35,2.75);
% Boissons (eau)
\draw[fill=popblue!45] (0.7,0.7) rectangle (5.3,1.5);
% Glace (triangles stylisés)
\draw[fill=white, draw=popblue, thick] (1.2,1.5) -- (1.9,1.5) -- (1.75,2.1) -- (1.35,2.1) -- cycle;
\draw[fill=white, draw=popblue, thick] (2.3,1.5) -- (3.0,1.5) -- (2.85,2.1) -- (2.45,2.1) -- cycle;
\draw[fill=white, draw=popblue, thick] (3.4,1.5) -- (4.1,1.5) -- (3.95,2.1) -- (3.55,2.1) -- cycle;
% Couvercle
\draw[fill=popgreen!30, thick] (-0.1,3.4) rectangle (6.1,3.75);
% Thermomètre
\draw[thick] (4.6,1.0) -- (4.6,3.55);
\draw[fill=poppink!60] (4.6,1.0) circle (0.12);
% Flèches RAYONNEMENT (soleil -> couvercle)
\draw[poporange, thick, ->, decorate, decoration={snake, amplitude=1.5pt, segment length=4pt}] (3,4.5) -- (3,3.75);
\node[poporange, align=center] at (3,4.8) {Rayonnement\\solaire};
% Flèches CONDUCTION (parois)
\draw[popred, thick, ->] (5.8,1.8) -- (5.35,1.8);
\node[popred, align=left] at (6.5,1.8) {Conduction\\parois};
\draw[popred, thick, ->] (3,0.1) -- (3,0.35);
\node[popred, align=center] at (3,-0.2) {Conduction\\fond (étal)};
% Flèches CONVECTION (air chaud -> parois, ouverture)
\draw[popgreen!70!black, thick, ->] (6.2,2.5) .. controls (6.5,2.8) and (6.5,3.2) .. (5.8,3.4);
\node[popgreen!70!black, align=left] at (6.8,2.8) {Convection\\air chaud};
\draw[popgreen!70!black, thick, ->] (2.5,3.8) -- (2.5,3.55);
\node[popgreen!70!black, align=center] at (2.5,4.1) {Convection\\ouverture};
% Annotations statiques
\node[align=left] at (7.5,3.5) {Couvercle};
\node[align=left] at (7.5,2.8) {Paroi isolante};
\node[align=left] at (7.5,1.2) {Boissons ($m_b$, $c_e$)};
\node[align=left] at (7.5,0.4) {Vase intérieur};
\node[align=left] at (2.5,2.4) {Glace à \SI{0}{\celsius}};
\node[align=left] at (4.6,3.9) {Thermomètre};
\end{tikzpicture}
\legende{Glacière assimilée à un calorimètre : identification des trois transferts thermiques parasites.}
\end{popfigure}

\textbf{2. Modes de transfert à bloquer.}
\begin{itemize}
\item \textbf{Rayonnement} : le Soleil chauffe le couvercle et les faces exposées de la glacière, même sans contact (ondes infrarouges/visibles).
\item \textbf{Conduction} : la chaleur traverse les parois solides (couvercle, parois latérales, fond) de l'extérieur vers l'intérieur ; elle remonte aussi de l'étal en bois chaud par le fond de la glacière.
\item \textbf{Convection} : l'air chaud ambiant circule autour de la glacière et chauffe ses parois externes (convection naturelle) ; à chaque ouverture du couvercle, de l'air chaud entre directement en contact avec les boissons (convection forcée).
\end{itemize}

\textbf{3. Masse de glace pour absorber $Q_{\text{ext}}$ par fusion seule (modèle simplifié).}
On suppose que les boissons sont déjà à \SI{10}{\celsius} et que l'eau de fusion ne se réchauffe pas. La glace fond à \SI{0}{\celsius} pour absorber l'apport extérieur $Q_{\text{ext}}$.
\[ Q_{\text{ext}} = m_1 L_f \quad \Rightarrow \quad m_1 = \frac{Q_{\text{ext}}}{L_f} = \frac{\num{2,5e5}\ \si{J}}{\num{3,34e5}\ \si{J.kg^{-1}}} \approx \SI{0,75}{kg} \]
\textit{Vérification dimensionnelle :} $\dfrac{\si{J}}{\si{J.kg^{-1}}} = \si{kg}$. \textbf{Résultat :} $m_1 \approx \SI{0,75}{kg}$.

\textbf{4. Bilan énergétique complet et masse totale $m_g$.}
Le système \{boissons + parois + glace + eau de fusion\} n'est \textbf{pas calorifugé} : il reçoit $Q_{\text{ext}}$ de l'extérieur. On applique le premier principe : $\Delta U = Q_{\text{ext}}$ (travail nul). Comme l'état final est à $\theta_f = \SI{10}{\celsius}$, on écrit l'égalité entre \textbf{chaleurs gagnées} (par la glace) et \textbf{chaleurs perdues} (par les corps chauds) \textbf{augmentées de l'apport extérieur} :

\begin{align*}
\textbf{Corps qui CÈDENT de la chaleur (refroidissement \SI{30}{\celsius} $\rightarrow$ \SI{10}{\celsius}) :} \\
Q_{\text{boissons}} &= m_b c_e (\theta_f - \theta_i) = 5 \times 4185 \times (10 - 30) = \SI{-4,185e5}{J} \\
Q_{\text{parois}} &= \mu c_e (\theta_f - \theta_i) = 0,300 \times 4185 \times (-20) = \SI{-2,511e4}{J} \\
\text{Valeurs absolues (chaleurs perdues) :} \quad |Q_{\text{boissons}}| &= \SI{4,185e5}{J}, \quad |Q_{\text{parois}}| = \SI{2,51e4}{J} \\[1em]
\textbf{Apport EXTÉRIEUR (chaleur REÇUE par le système) :} \\
Q_{\text{ext}} &= \num{2,5e5}\ \si{J} \\[1em]
\textbf{Corps qui REÇOIVENT de la chaleur (glace $\rightarrow$ eau à \SI{10}{\celsius}) :} \\
Q_{\text{glace}} &= m_g L_f + m_g c_e (\theta_f - 0) = m_g (L_f + c_e \theta_f) \\[1em]
\textbf{Bilan : Chaleurs gagnées = Chaleurs perdues + Apport extérieur} \\
m_g (L_f + c_e \theta_f) &= |Q_{\text{boissons}}| + |Q_{\text{parois}}| + Q_{\text{ext}} \\
m_g &= \frac{m_b c_e (\theta_i - \theta_f) + \mu c_e (\theta_i - \theta_f) + Q_{\text{ext}}}{L_f + c_e \theta_f}
\end{align*}

Calcul numérique :
\begin{align*}
\text{Numérateur} &= 5 \times 4185 \times 20 \;+\; 0,300 \times 4185 \times 20 \;+\; 2,5\times10^5 \\
&= \SI{418500}{J} + \SI{25110}{J} + \SI{250000}{J} = \SI{693610}{J} \\
\text{Dénominateur} &= 3,34\times10^5 + 4185 \times 10 = \SI{334000}{J.kg^{-1}} + \SI{41850}{J.kg^{-1}} = \SI{375850}{J.kg^{-1}} \\
m_g &= \frac{693610}{375850} \approx \SI{1,85}{kg}
\end{align*}
\textbf{Résultat :} Il faut environ \textbf{\SI{1,85}{kg}} de glace le matin (arrondi à \SI{1,9}{kg} avec 2 chiffres significatifs).

\textbf{5. Comparaison et conclusion.}
\begin{itemize}
\item $m_1 \approx \SI{0,75}{kg}$ (fusion seule pour compenser les fuites).
\item $m_g \approx \SI{1,85}{kg}$ (bilan réel avec refroidissement des boissons).
\item $m_g > m_1$ : la masse de glace nécessaire est \textbf{2,5 fois plus grande} que l'estimation naïve.
\item \textbf{Explication :} Les boissons chaudes (\SI{30}{\celsius}) cèdent une énorme quantité de chaleur (\SI{4,4e5}{J}) en refroidissant à \SI{10}{\celsius}. Cette chaleur \emph{s'ajoute} à la fuite extérieure (\SI{2,5e5}{J}) pour faire fondre la glace. Si Mama Élise ne met que \SI{0,75}{kg} de glace, elle fondra en absorbant seulement \SI{2,5e5}{J} (fuites) + une partie du refroidissement, mais il restera trop de chaleur dans les boissons : elles ne descendront pas à \SI{10}{\celsius} et la glace disparaîtra bien avant 17 h.
\item \textbf{Leçon :} Pré-refroidir les boissons (les mettre au frigo la veille ou les laisser dans l'eau courante) réduit drastiquement la masse de glace nécessaire. Si $\theta_i = \theta_f = \SI{10}{\celsius}$, le bilan donne $m_g = Q_{\text{ext}} / (L_f + c_e\theta_f) \approx \SI{0,67}{kg}$.
\end{itemize}

\textbf{6. Deux améliorations concrètes et peu coûteuses.}
\begin{enumerate}
\item \textbf{Peindre la glacière en blanc (ou la couvrir d'un tissu clair humide).} \emph{Justification :} Une surface claire a une faible émissivité/absorptivité dans le visible et l'infrarouge proche. Elle réfléchit le \textbf{rayonnement} solaire au lieu de l'absorber, réduisant fortement l'échauffement du couvercle et des parois.
\item \textbf{Surélever la glacière sur des cales (briques, bois) et limiter les ouvertures.} \emph{Justification :} Le contact direct avec l'étal en bois chaud (chauffé au soleil) crée un flux de \textbf{conduction} par le fond. La surélever crée une couche d'air isolant. Ouvrir le couvercle le moins souvent possible et refermer vite limite l'entrée d'air chaud par \textbf{convection}.
\end{enumerate}
\end{exoresolu}

\begin{attention}[Erreurs fréquentes au Probatoire — Check-list de relecture]
\begin{itemize}
\item \textbf{Signe de $\Delta\theta$ :} $Q = m c (\theta_f - \theta_i)$. Si $\theta_f < \theta_i$, $Q<0$ (cédé). Ne pas écrire $m c (\theta_i - \theta_f)$ en gardant le signe $+$ dans le bilan $\sum Q=0$ sans expliciter « chaleur perdue ».
\item \textbf{Eau de fusion oubliée :} Le terme $m_g c_e \theta_f$ (réchauffement de l'eau de $0^\circ$C à $\theta_f$) est \textbf{obligatoire} tant que $\theta_f > 0^\circ$C.
\item \textbf{Valeur en eau $\mu$ oubliée :} Les parois du calorimètre (vase, couvercle, agitateur, thermomètre) ont une capacité thermique. On \textbf{doit} inclure $\mu c_e \Delta\theta$.
\item \textbf{Unités mélangées :} Masse en \textbf{kg}, $L_f$ en \textbf{J/kg}, $c$ en \textbf{J$\cdot$kg$^{-1}$$\cdot$K$^{-1}$}. Convertir \SI{300}{g} en \SI{0,300}{kg} avant calcul.
\item \textbf{Système non calorifugé :} Ne pas écrire $\sum Q_i = 0$ si $Q_{\text{ext}} \neq 0$. Écrire « Chaleurs gagnées = Chaleurs perdues + $Q_{\text{ext}}$ » ou $\Delta U = Q_{\text{ext}}$.
\end{itemize}
\end{attention}

\begin{savaistu}[Refroidissement sans électricité : le « Pot-in-Pot » nigérien (Zeer)]
Mohammed Bah Abba (Nigeria, 2000) a modernisé la jarre traditionnelle : deux pots en terre cuite imbriqués, sable humide entre les deux, couvercle humide. L'évaporation de l'eau dans le sable puise la chaleur latente $L_v$ à l'intérieur $\rightarrow$ température chute de 15 à 20°C en dessous de l'ambiant. Conservation tomates/poivrons : 2-3 jours $\rightarrow$ 3 semaines. Prix Nobel alternatif 2000. Au Cameroun, des ONG diffusent ce modèle dans l'Extrême-Nord pour la conservation des vaccins (chaîne du froid) et des légumes.
\end{savaistu}

\courssub{Exercice d'entraînement : Le thé glacé du lycéen}
Un élève verse \SI{200}{mL} de thé bouillant (\SI{100}{\celsius}, $c \approx c_e$) dans un gobelet en plastique (valeur en eau $\mu = \SI{10}{g}$) et ajoute des glaçons à \SI{0}{\celsius}. Il veut obtenir du thé à \SI{5}{\celsius}. On néglige les échanges avec l'extérieur.
\begin{enumerate}
\item Quelle masse de glace $m_g$ faut-il ajouter pour que toute la glace fonde et que le mélange atteigne \SI{5}{\celsius} ?
\item Si l'élève met \SI{50}{g} de glace seulement, quel est l'état final (température, phases présentes) ?
\end{enumerate}

\begin{corrige}[Exercice : Le thé glacé]
\textbf{Données :} $m_{thé} = 0,200$ kg, $\theta_i = 100^\circ$C, $\mu = 0,010$ kg, $\theta_f = 5^\circ$C, $c_e = 4185$ J/kg/K, $L_f = 3,34\times10^5$ J/kg.
\textbf{1.} Bilan : Chaleur perdue par thé + gobelet = Chaleur gagnée par glace (fusion + réchauffement).
\[ (m_{thé}+\mu) c_e (100-5) = m_g L_f + m_g c_e (5-0) \]
\[ 0,210 \times 4185 \times 95 = m_g (3,34\times10^5 + 4185 \times 5) \]
\[ 83474 = m_g \times 354925 \quad \Rightarrow \quad m_g \approx \SI{0,235}{kg} = \SI{235}{g} \]
\textbf{2.} Avec $m_g = \SI{50}{g} = 0,050$ kg.
Chaleur max absorbable par la glace (fusion + réchauffement à 100°C) : $0,050 \times (3,34\times10^5 + 4185 \times 100) = 0,050 \times 752500 = 37625$ J.
Chaleur que thé+gobelet peuvent céder en allant à $0^\circ$C : $0,210 \times 4185 \times 100 = 87885$ J $> 37625$ J.
$\rightarrow$ La glace fond \textbf{entièrement} mais le mélange final est à une température $\theta_f > 0^\circ$C.
Calcul $\theta_f$ : $(m_{thé}+\mu) c_e (100-\theta_f) = m_g L_f + m_g c_e \theta_f$
$87885 - 878,85 \theta_f = 16700 + 209,25 \theta_f$
$71185 = 1088,1 \theta_f \Rightarrow \theta_f \approx \SI{65,4}{\celsius}$.
\textbf{État final :} Tout liquide (eau + thé), $\theta_f \approx \SI{65}{\celsius}$. Pas de glace résiduelle.
\end{corrige}

\courssub{Bilan de la leçon}
Cette leçon t'a permis de maîtriser les \textbf{trois modes de transfert de chaleur} (conduction, convection, rayonnement) et de les identifier dans des situations concrètes (glacière, habitat, industrie). Tu sais désormais :
\begin{itemize}
\item Calculer une quantité de chaleur \textbf{sans changement d'état} ($Q = mc\Delta\theta$) et \textbf{avec changement d'état} ($Q = mL$).
\item Utiliser la \textbf{valeur en eau} $\mu$ pour intégrer le calorimètre dans les bilans.
\item Appliquer rigoureusement le \textbf{premier principe} (bilan énergétique) à des systèmes réels, \textbf{calorifugés ou non} ($\sum Q_{\text{internes}} = Q_{\text{ext}}$).
\item Modéliser un objet du quotidien (glacière) en calorimètre, écrire le bilan, calculer, et \textbf{proposer des améliorations technologiques justifiées physiquement}.
\end{itemize}
Ces compétences sont au cœur du programme de Première C et constituent la base de la thermodynamique (cycles moteurs, pompes à chaleur, climatisation) que tu approfondiras en Terminale. Continue t'entraîner avec des bilans variés : mélange eau/eau, métal/eau, glace/eau, avec ou sans pertes !

\newpage

\coursec{Méthodes et exercices résolus}

\coursec{Modes de transfert de chaleur et calorimétrie}

\courssub{Chaleur, température et échanges thermiques}

\begin{definition}[Chaleur et température]
La \cle{température} $\theta$ (en \si{\celsius} ou \si{K}) est une grandeur macroscopique qui traduit le degré d'agitation thermique microscopique des particules (atomes, molécules) d'un corps. La \cle{chaleur} $Q$ (en \si{J}) est une \textbf{énergie} transférée d'un corps chaud vers un corps froid \emph{uniquement} à cause d'une différence de température. Ce n'est pas une propriété que le corps « possède » (on ne dit pas « la chaleur du corps » mais « son énergie interne »).
\end{definition}

\begin{propriete}[Équilibre thermique]
Deux corps mis en contact thermique évoluent vers un état d'\cle{équilibre thermique} où leurs températures deviennent égales. Pendant ce processus, le corps chaud \textbf{cède} de la chaleur ($Q<0$), le corps froid \textbf{reçoit} de la chaleur ($Q>0$).
\end{propriete}

\begin{aretenir}[Distinction fondamentale]
\begin{itemize}
\item \textbf{Température} : intensité de l'agitation microscopique (grandeur intensive, ne s'additionne pas).
\item \textbf{Chaleur} : quantité d'énergie transférée (grandeur extensive, s'additionne).
\item \textbf{Énergie interne} $U$ : énergie microscopique totale (agitation + liaisons). Pour un corps pur sans changement d'état, $\Delta U = Q + W$ (1er principe). En calorimétrie usuelle ($W=0$), $\Delta U = Q$.
\end{itemize}
\end{aretenir}

\begin{savaistu}[Le feu et l'homme au Cameroun]
Depuis la nuit des temps, la maîtrise du feu (combustion bois, charbon) repose sur les transferts thermiques : \textbf{rayonnement} pour se chauffer autour du foyer, \textbf{convection} pour faire monter l'air chaud dans les cases traditionnelles à toit conique, \textbf{conduction} pour cuire dans les marmites en aluminium ou en terre cuite. Aujourd'hui, les barrages de \textbf{Lom Pangar}, \textbf{Nachtigal} et \textbf{Edéa} transforment l'énergie potentielle de l'eau en électricité ; les lignes \textbf{SONATREL/ENEO} transportent cette énergie qui, par effet Joule, devient chaleur dans nos foyers (fers à repasser, bouilloires, climatiseurs).
\end{savaistu}

\courssub{Les trois modes de transfert de chaleur}

\begin{definition}[Conduction]
Transfert de chaleur dans un \textbf{solide} (ou fluide au repos) par choc successif entre particules voisines, \textbf{sans déplacement macroscopique de matière}. Les métaux (Cu, Al, Fe) sont de bons conducteurs ($\lambda \sim \SI{50}{W.m^{-1}.K^{-1}}$ à $\SI{400}{W.m^{-1}.K^{-1}}$) ; le bois, le plastique, l'air immobile, l'eau stagnante sont de mauvais conducteurs (isolants).
\end{definition}

\begin{definition}[Convection]
Transfert de chaleur dans un \textbf{fluide} (liquide ou gaz) \textbf{en mouvement} (courants de convection). L'échauffement local diminue la densité $\rho$ ; le fluide chaud monte (poussée d'Archimède), le fluide froid descend : c'est la \cle{convection naturelle}. Un ventilateur ou une pompe impose la \cle{convection forcée} (radiateurs de voiture, refroidissement processeurs, séchoirs agricoles).
\end{definition}

\begin{definition}[Rayonnement thermique]
Transfert d'énergie par \textbf{ondes électromagnétiques} (surtout infrarouge $\lambda \approx \SIrange{0,7}{1000}{\micro\meter}$). \textbf{Ne nécessite aucun support matériel} : fonctionne dans le vide. Tout corps à $T>0\,\si{K}$ émet. Puissance émise : loi de Stefan-Boltzmann $P = e\sigma S T^4$ ($\sigma = \SI{5,67e-8}{W.m^{-2}.K^{-4}}$, $e$ émissivité). C'est le mode dominant au-dessus de \SI{500}{\celsius} et le \textbf{seul} mode dans le vide (Soleil $\to$ Terre).
\end{definition}

\begin{popfigure}
\begin{tikzpicture}[scale=0.9, every node/.style={font=\small}]
% Conduction
\begin{scope}[xshift=-7cm]
\draw[popblue, thick, fill=popblueL] (0,0) rectangle (2,1.5);
\foreach \y in {0.2,0.5,0.8,1.1} {
  \foreach \x in {0.2,0.6,1.0,1.4,1.8} {
    \fill[popdark] (\x,\y) circle (1.5pt);
  }
}
\draw[poporange, thick, ->] (1,1.5) -- (1,2.2) node[above] {Flux $Q$};
\draw[poporange, thick, ->] (1,-0.2) -- (1,-0.9);
\draw[decorate, decoration={snake, amplitude=1pt, segment length=3pt}, poporange] (0.5,0.5) -- (0.5,1.0);
\draw[decorate, decoration={snake, amplitude=1pt, segment length=3pt}, poporange] (1.5,0.5) -- (1.5,1.0);
\node[align=center] at (1,-1.3) {\textbf{Conduction}\\(solide immobile)};
\end{scope}

% Convection
\begin{scope}[xshift=0cm]
\draw[popblue, thick] (0,0) rectangle (2,2);
\draw[popblue, thick, fill=popblueL] (0,0) rectangle (2,0.2);
\draw[popblue, thick, fill=popblueL] (0,1.8) rectangle (2,2);
\draw[->, thick, poporange] (0.5,0.2) -- (0.5,1.8) node[midway, right] {Chaud $\uparrow$};
\draw[->, thick, popgreen] (1.5,1.8) -- (1.5,0.2) node[midway, left] {Froid $\downarrow$};
\draw[poporange, thick] (0.5,0.2) .. controls (0.8,1.0) .. (1.5,1.8);
\draw[popgreen, thick] (1.5,1.8) .. controls (1.2,1.0) .. (0.5,0.2);
\node[align=center] at (1,-0.5) {\textbf{Convection}\\(fluide en mouvement)};
\end{scope}

% Rayonnement
\begin{scope}[xshift=7cm]
\draw[popgold, thick, fill=popgoldL] (0.5,1) circle (0.6);
\node at (0.5,1) {$\odot$};
\foreach \angle/\len in {30/1.5, 60/1.8, 90/2.0, 120/1.8, 150/1.5} {
  \draw[poporange, thick, ->] ({0.5+0.6*cos(\angle)}, {1+0.6*sin(\angle)}) -- ++(\angle:\len);
}
\draw[popblue, thick, fill=popblueL] (3.5,0.5) rectangle (4.5,1.5);
\node at (4,1) {Terre};
\node[align=center] at (2,-0.5) {\textbf{Rayonnement}\\(vide, ondes EM)};
\end{scope}
\end{tikzpicture}
\legende{Les trois modes de transfert de chaleur : conduction (solide), convection (fluide), rayonnement (vide).}
\end{popfigure}

\begin{methode}[Reconnaître conduction, convection et rayonnement]
Pour identifier le(s) mode(s) de transfert dans une situation :
\begin{enumerate}
\item Y a-t-il un \cle{support matériel} (solide ou fluide) entre la source et le récepteur ?
\item \textbf{Si oui} : la matière se déplace-t-elle macroscopiquement ?
  \begin{itemize}
  \item Matière \textbf{immobile} (solide, ou fluide sans courant) $\to$ \cle{conduction}.
  \item Matière en \textbf{mouvement organisé} (courants, vent, pompe) $\to$ \cle{convection}.
  \end{itemize}
\item \textbf{Si non} (vide, ou transfert à distance sans contact direct) $\to$ \cle{rayonnement}.
\item \textbf{Réalité} : plusieurs modes coexistent souvent. Cite-les tous et précise lequel \textbf{domine} (ex: poêle à bois : conduction dans la fonte, convection dans la pièce, rayonnement vers les murs).
\end{enumerate}
\textbf{Réflexe Probatoire} : Le Soleil chauffe la Terre \emph{uniquement} par rayonnement (le vide spatial interdit conduction et convection).
\end{methode}

\begin{exoresolu}[Modes de transfert dans une cuisine camerounaise]
Une marmite en aluminium contenant de l'eau est posée sur un foyer à bois. Une cuisinière tient la queue métallique de la marmite et sent la chaleur monter dans sa main ; elle sent aussi la chaleur du feu sur son visage, sans le toucher. Pour chaque situation soulignée, indiquer le mode de transfert de chaleur en justifiant.
\tcblower
\textbf{Raisonnement.} On applique la grille : support matériel ? Mouvement de matière ?
\begin{itemize}
\item \textbf{Chauffage de l'eau dans la marmite} : La chaleur traverse le fond métallique par \textbf{conduction} (solide immobile). L'eau au fond s'échauffe, sa densité diminue, elle monte ; l'eau froide descend : ce sont des courants de \textbf{convection naturelle} qui brassent l'eau et uniformisent la température.
\item \textbf{Chaleur ressentie dans la main via la queue métallique} : Le métal est un solide, les atomes transmettent l'agitation thermique de proche en proche sans déplacement de matière : c'est la \textbf{conduction}. (C'est pourquoi on utilise un manche en bois ou plastique, mauvais conducteurs).
\item \textbf{Chaleur ressentie sur le visage} : Pas de contact avec le feu ; le rayonnement infrarouge émis par les braises traverse l'air (qu'il chauffe peu) : c'est le \textbf{rayonnement}.
\end{itemize}
\textbf{Conclusion.} Conduction à travers le métal, convection dans l'eau, rayonnement entre le feu et le visage : les trois modes coexistent dans cette situation quotidienne.
\end{exoresolu}

\courssub{Quantité de chaleur sans changement d'état}

\begin{definition}[Chaleur massique (capacité thermique massique) $c$]
Quantité de chaleur nécessaire pour élever de \SI{1}{K} (ou \SI{1}{\celsius}) la température de \SI{1}{kg} d'un corps pur, \textbf{sans changement d'état}.
Unité SI : \si{J.kg^{-1}.K^{-1}}. Valeurs usuelles : $c_{\text{eau}} \approx \SI{4185}{J.kg^{-1}.K^{-1}}$, $c_{\text{Al}} \approx \SI{897}{J.kg^{-1}.K^{-1}}$, $c_{\text{Cu}} \approx \SI{385}{J.kg^{-1}.K^{-1}}$, $c_{\text{Fe}} \approx \SI{460}{J.kg^{-1}.K^{-1}}$.
\end{definition}

\begin{definition}[Capacité calorifique $C$]
Quantité de chaleur nécessaire pour élever de \SI{1}{K} la température d'un \textbf{système donné} (masse $m$ de substance, ou objet composite).
$C = m\,c$ \quad Unité SI : \si{J.K^{-1}}.
\end{definition}

\begin{propriete}[Quantité de chaleur échangée sans changement d'état]
Pour un corps de masse $m$, chaleur massique $c$, dont la température varie de $\theta_i$ à $\theta_f$ sans changement de phase :
\[ Q = m\,c\,(\theta_f - \theta_i) = m\,c\,\Delta\theta \quad\text{ou}\quad Q = C\,\Delta\theta. \]
\textbf{Signe} : $Q>0$ si le corps \textbf{reçoit} de la chaleur ($\theta_f > \theta_i$, échauffement) ; $Q<0$ s'il en \textbf{cède} ($\theta_f < \theta_i$, refroidissement).
\emph{Analyse dimensionnelle} : $\si{kg}\times\si{J.kg^{-1}.K^{-1}}\times\si{K} = \si{J}$. \checkmark
\end{propriete}

\begin{methode}[Appliquer $Q = m\,c\,\Delta\theta$]
\begin{enumerate}
\item Vérifier qu'il n'y a \textbf{aucun changement d'état} pendant le processus (sinon voir Méthode 3).
\item Relever les données et convertir en unités SI : $m$ en \si{kg}, $\theta$ en \si{\celsius} ou \si{K} (un \emph{écart} $\Delta\theta$ est identique dans les deux échelles), $c$ en \si{J.kg^{-1}.K^{-1}}.
\item Calculer $\Delta\theta = \theta_f - \theta_i$.
\item Écrire $Q = m\,c\,\Delta\theta$ et faire l'application numérique.
\item Conclure avec le \textbf{signe} et l'unité : « Le corps reçoit $Q = \dots\ \si{J}$ » ou « Le corps cède $|Q| = \dots\ \si{J}$ ».
\end{enumerate}
\end{methode}

\begin{exoresolu}[Chauffage de l'eau pour le thé]
Pour préparer le thé, on chauffe $m = \num{1,5}\ \si{kg}$ d'eau de $\theta_i = \SI{25}{\celsius}$ à $\theta_f = \SI{80}{\celsius}$. On donne $c_{\text{eau}} = \SI{4185}{J.kg^{-1}.K^{-1}}$.
\begin{enumerate}
\item Calculer la quantité de chaleur reçue par l'eau.
\item Même question pour une casserole en aluminium de masse $\num{0,40}\ \si{kg}$ chauffée dans les mêmes conditions ($c_{\text{Al}} = \SI{897}{J.kg^{-1}.K^{-1}}$). Commenter.
\end{enumerate}
\tcblower
\textbf{1.} Pas de changement d'état entre \SI{25}{\celsius} et \SI{80}{\celsius} : $Q = m\,c\,\Delta\theta$ s'applique.
\[ \Delta\theta = \theta_f - \theta_i = 80 - 25 = \SI{55}{K} \]
\[ Q_{\text{eau}} = \num{1,5}\times 4185\times 55 = \SI{345262,5}{J} \approx \SI{3,45e5}{J} = \SI{345}{kJ}. \]
L'eau \textbf{reçoit} $Q_{\text{eau}} \approx \SI{3,5e5}{J}$ ($Q>0$).

\textbf{2.} Pour la casserole en aluminium :
\[ Q_{\text{Al}} = \num{0,40}\times 897\times 55 = \SI{19734}{J} \approx \SI{1,97e4}{J} = \SI{19,7}{kJ}. \]
\textbf{Commentaire.} Bien que de masse comparable, la casserole demande environ 17 fois moins de chaleur que l'eau : la chaleur massique de l'eau est exceptionnellement élevée. C'est pourquoi l'eau est un excellent « réservoir » de chaleur (calorifères, refroidissement moteurs, bouteilles d'eau chaude).
\textbf{Conclusion.} $Q_{\text{eau}} \approx \SI{3,45e5}{J}$ reçus ; $Q_{\text{Al}} \approx \SI{1,97e4}{J}$ reçus.
\end{exoresolu}

\courssub{Chaleur latente et changements d'état}

\begin{definition}[Chaleur latente massique $L$]
Quantité de chaleur nécessaire pour faire changer d'état \SI{1}{kg} d'un corps pur \textbf{à température constante}.
Unité SI : \si{J.kg^{-1}}.
\begin{itemize}
\item \textbf{Fusion} (solide $\to$ liquide) à $\theta_f$ : $L_f$ (ex: eau $L_f = \SI{3,34e5}{J.kg^{-1}}$ à \SI{0}{\celsius}).
\item \textbf{Vaporisation} (liquide $\to$ gaz) à $\theta_v$ : $L_v$ (ex: eau $L_v = \SI{2,26e6}{J.kg^{-1}}$ à \SI{100}{\celsius}).
\item Solidification et liquéfaction : même valeur, chaleur \textbf{cédée} ($Q<0$).
\end{itemize}
\end{definition}

\begin{popfigure}
\begin{tikzpicture}
\begin{axis}[
    width=\linewidth, height=7cm,
    axis lines=middle,
    xlabel={Énergie fournie $Q$ (\si{kJ}) pour $m=\SI{1}{kg}$ d'eau}, ylabel={Température $\theta$ (\si{\celsius})},
    xmin=0, xmax=3000, ymin=-20, ymax=120,
    xtick={0, 334, 752, 3012}, xticklabels={0, $L_f$, $c\Delta\theta$, $L_f+c\Delta\theta+L_v$},
    ytick={0, 100}, yticklabels={\SI{0}{\celsius} (fusion), \SI{100}{\celsius} (vaporisation)},
    domain=0:3000, samples=200,
    popcurve, thick,
    clip=false
]
% Glace chauffage
\addplot[domain=0:334] { -10 + (10/334)*x };
% Fusion plateau
\addplot[domain=334:752] { 0 };
% Eau chauffage
\addplot[domain=752:3012] { (100/(3012-752))*(x-752) };
% Vaporisation plateau
\addplot[domain=3012:3000] { 100 }; % just a point
% Annotations
\node[align=center, popblue] at (167, -15) {Glace\\$c_{\text{glace}}$};
\node[align=center, poporange] at (543, 15) {Fusion\\$L_f$};
\node[align=center, popgreen] at (1882, 50) {Eau\\$c_{\text{eau}}$};
\node[align=center, poppurple] at (3100, 105) {Vaporisation\\$L_v$};
\draw[dashed, thin] (334, -20) -- (334, 120);
\draw[dashed, thin] (752, -20) -- (752, 120);
\draw[dashed, thin] (3012, -20) -- (3012, 120);
\end{axis}
\end{tikzpicture}
\legende{Courbe de chauffage de \SI{1}{kg} d'eau de \SI{-10}{\celsius} à la vapeur. Les paliers correspondent aux changements d'état (température constante).}
\end{popfigure}

\begin{methode}[Appliquer $Q = m\,L$ et enchaîner les étapes]
\begin{enumerate}
\item Repérer le \textbf{changement d'état} : fusion, solidification, vaporisation, liquéfaction. Pendant toute la durée du changement d'état, la température reste \textbf{constante} : \textbf{interdit} d'utiliser $mc\Delta\theta$ pendant cette phase.
\item Pendant le changement d'état : $Q = m\,L$ ($L$ en \si{J.kg^{-1}}). $Q>0$ si le corps reçoit (fusion, vaporisation) ; $Q<0$ s'il cède (solidification, liquéfaction).
\item Si le problème comporte un chauffage \textbf{puis} un changement d'état (ou l'inverse), découper en étapes successives :
\[ Q_{\text{totale}} = \sum m\,c\,\Delta\theta \quad(\text{phases sans changement d'état}) + \sum m\,L \quad(\text{changements d'état}). \]
\item Vérifier la cohérence des unités : $m$ en \si{kg}, $L$ en \si{J.kg^{-1}} $\to$ $Q$ en \si{J}.
\end{enumerate}
Valeurs utiles pour l'eau : $L_f = \SI{3,34e5}{J.kg^{-1}}$ (fusion à \SI{0}{\celsius}), $L_v = \SI{2,26e6}{J.kg^{-1}}$ (vaporisation à \SI{100}{\celsius}).
\end{methode}

\begin{exoresolu}[De la glace à la vapeur]
Quelle quantité de chaleur faut-il fournir à $m = \num{200}\ \si{g}$ de glace prise à $\SI{-10}{\celsius}$ pour la transformer entièrement en vapeur à $\SI{100}{\celsius}$ ?
Données : $c_{\text{glace}} = \SI{2090}{J.kg^{-1}.K^{-1}}$, $c_{\text{eau}} = \SI{4185}{J.kg^{-1}.K^{-1}}$, $L_f = \SI{3,34e5}{J.kg^{-1}}$, $L_v = \SI{2,26e6}{J.kg^{-1}}$.
\tcblower
\textbf{Raisonnement.} Le processus se décompose en quatre étapes successives ; on additionne les chaleurs de chaque étape. Conversion : $m = \num{0,200}\ \si{kg}$.

\textbf{Étape 1 : chauffage de la glace de $-10\ \si{\celsius}$ à $0\ \si{\celsius}$.}
\[ Q_1 = m\,c_{\text{glace}}\,\Delta\theta_1 = \num{0,200}\times 2090\times 10 = \SI{4,18e3}{J}. \]

\textbf{Étape 2 : fusion complète à $0\ \si{\celsius}$.}
\[ Q_2 = m\,L_f = \num{0,200}\times \num{3,34e5} = \SI{6,68e4}{J}. \]

\textbf{Étape 3 : chauffage de l'eau de $0\ \si{\celsius}$ à $100\ \si{\celsius}$.}
\[ Q_3 = m\,c_{\text{eau}}\,\Delta\theta_3 = \num{0,200}\times 4185\times 100 = \SI{8,37e4}{J}. \]

\textbf{Étape 4 : vaporisation complète à $100\ \si{\celsius}$.}
\[ Q_4 = m\,L_v = \num{0,200}\times \num{2,26e6} = \SI{4,52e5}{J}. \]

\textbf{Bilan.}
\begin{align*}
Q &= Q_1 + Q_2 + Q_3 + Q_4 \\
&= \num{4,18e3} + \num{6,68e4} + \num{8,37e4} + \num{4,52e5} \\
&= \SI{606680}{J} \approx \SI{6,07e5}{J} = \SI{607}{kJ}.
\end{align*}
\textbf{Conclusion.} Il faut fournir environ $\SI{6,1e5}{J}$. On remarque que la vaporisation ($Q_4$) représente à elle seule environ $75\,\%$ de l'énergie totale : c'est l'étape la plus coûteuse.
\end{exoresolu}

\courssub{Le calorimètre et le principe des échanges de chaleur}

\begin{experience}[Le calorimètre à mélange — Schéma annoté et fonctionnement]
\textbf{Matériel.} Calorimètre (vase en aluminium ou inox, parois doubles pour isolation), agitateur (mélangeur), thermomètre (précision \SI{0,1}{\celsius}), couvercle isolant (bouchon liège/plastique avec trous pour thermomètre et agitateur), balance, étuve ou bain-marie, corps d'essai (métal, glace), eau.
\textbf{Schéma.}
\begin{popfigure}
\begin{tikzpicture}[scale=0.8, every node/.style={font=\small}]
% Vase double paroi
\draw[popblue, thick, fill=popblueL!30] (0,0) -- (0,5) -- (4,5) -- (4,0) -- cycle;
\draw[popblue, thick, fill=popblueL!30] (0.3,0) -- (0.3,4.5) -- (4.2,4.5) -- (4.2,0.5) -- cycle;
% Eau
\draw[popblue, fill=popblueL!50] (0.3,0.5) rectangle (4.2,3.5);
% Agitateur
\draw[popdark, thick] (2.25,4.5) -- (2.25,1.5);
\draw[popdark, thick] (1.75,1.5) -- (2.75,1.5);
\draw[popdark, thick] (1.75,1.5) -- (1.9,0.8);
\draw[popdark, thick] (2.75,1.5) -- (2.6,0.8);
% Thermomètre
\draw[poporange, thick] (3.5,4.5) -- (3.5,1.0);
\draw[poporange, thick, fill=poporangeL] (3.3,1.0) rectangle (3.7,0.5);
% Couvercle
\draw[popgreen, thick, fill=popgreenL!30] (0,4.5) rectangle (4.5,5.0);
% Corps d'essai (boule)
\draw[poppurple, thick, fill=poppurpleL] (2.25,4.0) circle (0.4);
% Légendes
\node[align=left, popblue] at (-1.5, 2) {Vase double paroi\\isolant};
\node[align=left, popdark] at (5.5, 3) {Agitateur\\(homogénéisation)};
\node[align=left, poporange] at (5.5, 1.5) {Thermomètre\\$\theta_{eq}$};
\node[align=left, popgreen] at (5.5, 4.7) {Couvercle isolant};
\node[align=left, poppurple] at (-1.5, 4) {Corps d'essai\\chaud/froid};
\node[align=center] at (2.25, -0.5) {Eau de masse $m_f$};
\end{tikzpicture}
\legende{Schéma annoté d'un calorimètre à mélange. La valeur en eau $\mu$ caractérise la capacité calorifique du vase + accessoires (couvercle, agitateur, thermomètre).}
\end{popfigure}
\textbf{Protocole type.} 1) Mesurer $\mu$ (Méthode 5). 2) Mettre eau froide ($m_f, \theta_f$). 3) Chauffer corps d'essai ($m_c, \theta_c$). 4) Plonger corps d'essai, refermer, agiter. 5) Lire $\theta_{eq}$ stable. 6) Bilan $\sum Q = 0$.
\end{experience}

\begin{propriete}[Principe des échanges de chaleur (Calorimètre isolé)]
Un calorimètre est un système \textbf{quasi isolé thermiquement} (parois doubles, couvercle, durée courte). La somme algébrique des quantités de chaleur échangées entre tous les composants du système est nulle :
\[ \sum_i Q_i = 0 \quad\text{avec}\quad Q_i = m_i\,c_i\,(\theta_{eq} - \theta_i) \;\text{(sans chgt d'état)}\quad\text{ou}\quad Q_i = \pm m_i L \;\text{(chgt d'état)}. \]
Pour le calorimètre (vase + accessoires), on utilise sa \cle{valeur en eau} $\mu$ (masse d'eau équivalente) : $Q_{\text{cal}} = \mu\,c_{\text{eau}}\,(\theta_{eq} - \theta_f)$.
\end{propriete}

\begin{methode}[Bilan thermique dans un calorimètre (sans changement d'état)]
\begin{enumerate}
\item Schématiser : identifier le corps \textbf{chaud} (température initiale $\theta_c$) et l'ensemble \textbf{froid} (eau $m_f, \theta_f$ + calorimètre $\mu, \theta_f$).
\item Écrire le principe : $\sum Q = 0$.
\item Exprimer chaque terme avec $\theta_{eq}$ (inconnue ou connue) :
  \begin{align*}
  Q_{\text{chaud}} &= m_c\,c_c\,(\theta_{eq} - \theta_c) \quad (\text{négatif car } \theta_{eq} < \theta_c) \\
  Q_{\text{froid}} &= (m_f\,c_{\text{eau}} + \mu\,c_{\text{eau}})\,(\theta_{eq} - \theta_f) \quad (\text{positif car } \theta_{eq} > \theta_f)
  \end{align*}
\item Résoudre l'équation (inconnue : $\theta_{eq}$, $c_c$, $m_c$ ou $\mu$).
\item \textbf{Vérifier} : $\theta_f < \theta_{eq} < \theta_c$ ; $c_c > 0$ et ordre de grandeur cohérent (métaux : quelques centaines \si{J.kg^{-1}.K^{-1}}).
\end{enumerate}
\end{methode}

\begin{exoresolu}[Détermination de la chaleur massique d'un métal]
Un calorimètre de valeur en eau $\mu = \SI{40}{g}$ contient $m_f = \SI{250}{g}$ d'eau à $\theta_f = \SI{20}{\celsius}$. On y plonge un bloc de cuivre de masse $m_c = \SI{150}{g}$ sortant d'une étuve à $\theta_c = \SI{95}{\celsius}$. La température d'équilibre est $\theta_{eq} = \SI{23,5}{\celsius}$. Donnée : $c_{\text{eau}} = \SI{4185}{J.kg^{-1}.K^{-1}}$. Déterminer $c_{\text{Cu}}$.
\tcblower
\textbf{Raisonnement.} Calorimètre isolé : cuivre cède, eau + calorimètre reçoivent. $\sum Q = 0$.
Masses en \si{kg} : $m_c = \num{0,150}$, $m_f = \num{0,250}$, $\mu = \num{0,040}$.

\textbf{Chaleur cédée par le cuivre :}
\[ Q_{\text{Cu}} = m_c\,c_{\text{Cu}}\,(\theta_{eq} - \theta_c) = \num{0,150}\,c_{\text{Cu}}\,(23,5 - 95) = -\num{10,725}\,c_{\text{Cu}}. \]

\textbf{Chaleur reçue par l'eau et le calorimètre :}
\[ Q_{\text{reçue}} = (m_f\,c_{\text{eau}} + \mu\,c_{\text{eau}})\,(\theta_{eq} - \theta_f) = (\num{0,250}+\num{0,040})\times 4185\times(23,5-20). \]
\[ Q_{\text{reçue}} = \num{0,290}\times 4185\times 3,5 = \SI{4247,775}{J} \approx \SI{4,25e3}{J}. \]

\textbf{Résolution :}
\[ -\num{10,725}\,c_{\text{Cu}} + \num{4247,775} = 0 \quad\Longrightarrow\quad c_{\text{Cu}} = \frac{4247,775}{10,725} \approx \SI{396}{J.kg^{-1}.K^{-1}}. \]

\textbf{Vérification.} $\theta_{eq} = \SI{23,5}{\celsius} \in [20; 95]$ \checkmark. Valeur positive, proche de la valeur tabulée du cuivre ($\SI{385}{J.kg^{-1}.K^{-1}}$), l'écart ($\approx 3\,\%$) provenant des fuites thermiques inévitables.
\textbf{Conclusion.} $c_{\text{Cu}} \approx \SI{3,96e2}{J.kg^{-1}.K^{-1}}$.
\end{exoresolu}

\begin{methode}[Mesure de la valeur en eau $\mu$ d'un calorimètre (Méthode des mélanges)]
\begin{enumerate}
\item Verser une masse $m_1$ d'eau \textbf{froide} à $\theta_1$ dans le calorimètre vide ; attendre l'équilibre (le vase prend $\theta_1$).
\item Ajouter une masse $m_2$ d'eau \textbf{chaude} à $\theta_2$ ; refermer, agiter, mesurer $\theta_{eq}$.
\item Bilan : eau chaude cède $Q_2 = m_2\,c_{\text{eau}}(\theta_{eq}-\theta_2)$ ; eau froide + calorimètre reçoivent $Q_1 = (m_1 + \mu)\,c_{\text{eau}}(\theta_{eq}-\theta_1)$.
\item $\sum Q = 0 \;\Rightarrow\; m_2(\theta_{eq}-\theta_2) + (m_1+\mu)(\theta_{eq}-\theta_1) = 0$ ($c_{\text{eau}}$ se simplifie).
\item Résoudre pour $\mu$ (en \si{kg} ou \si{g} selon les masses utilisées) :
\[ \mu = \frac{m_2(\theta_2 - \theta_{eq})}{\theta_{eq} - \theta_1} - m_1. \]
\end{enumerate}
\end{methode}

\begin{exoresolu}[Mesure de la valeur en eau]
Un calorimètre contient $m_1 = \SI{200}{g}$ d'eau à $\theta_1 = \SI{18}{\celsius}$. On y verse $m_2 = \SI{150}{g}$ d'eau à $\theta_2 = \SI{60}{\celsius}$. L'équilibre s'établit à $\theta_{eq} = \SI{33,2}{\celsius}$. Calculer $\mu$.
\tcblower
\textbf{Raisonnement.} Système isolé : $\sum Q = 0$. $c_{\text{eau}}$ se simplifie.
\[ m_2(\theta_{eq}-\theta_2) + (m_1+\mu)(\theta_{eq}-\theta_1) = 0. \]
Masses en \si{g} (résultat en \si{g}) :
\[ 150\times(33,2-60) + (200+\mu)\times(33,2-18) = 0 \]
\[ 150\times(-26,8) + (200+\mu)\times 15,2 = 0 \]
\[ -4020 + 15,2\,(200+\mu) = 0 \]
\[ 200+\mu = \frac{4020}{15,2} \approx 264,47 \quad\Longrightarrow\quad \mu \approx \SI{64,5}{g}. \]
\textbf{Vérification.} $\mu > 0$ \checkmark. Sans calorimètre ($\mu=0$), $\theta_{eq} = \frac{200\times18 + 150\times60}{350} = \SI{36}{\celsius}$. La présence du vase ($\mu>0$) abaisse l'équilibre à \SI{33,2}{\celsius} : cohérent.
\textbf{Conclusion.} La valeur en eau du calorimètre est $\mu \approx \SI{65}{g}$.
\end{exoresolu}

\begin{methode}[Bilan avec fusion de glace dans le calorimètre]
\begin{enumerate}
\item Acteurs : eau chaude ($m_1, \theta_1$), calorimètre ($\mu, \theta_1$), glace ($m_g, \SI{0}{\celsius}$).
\item La glace suit \textbf{deux} phases obligatoires :
  \begin{enumerate}
  \item Fusion à \SI{0}{\celsius} : $Q_{\text{fusion}} = m_g L_f$ (chaleur reçue, $>0$).
  \item Chauffage de l'eau de fusion de \SI{0}{\celsius} à $\theta_{eq}$ : $Q_{\text{chauf}} = m_g\,c_{\text{eau}}\,(\theta_{eq} - 0)$ (chaleur reçue, $>0$).
  \end{enumerate}
  \textbf{Ne jamais oublier la seconde phase !}
\item Écrire $\sum Q = 0$ :
\[ (m_1+\mu)\,c_{\text{eau}}\,(\theta_{eq}-\theta_1) + m_g L_f + m_g\,c_{\text{eau}}\,\theta_{eq} = 0. \]
\item \textbf{Vérification a priori} : toute la glace fond-elle ? Calculer $Q_{\text{dispo}} = (m_1+\mu)c_{\text{eau}}\theta_1$ (chaleur max cédée en allant à \SI{0}{\celsius}). Si $Q_{\text{dispo}} < m_g L_f$, l'équilibre est à \SI{0}{\celsius} avec mélange eau-glace (problème à deux inconnues : $\theta_{eq}=0$ et masse fondue).
\end{enumerate}
\end{methode}

\begin{exoresolu}[Fusion de glace dans un calorimètre]
Un calorimètre de valeur en eau $\mu = \SI{30}{g}$ contient $m_1 = \SI{300}{g}$ d'eau à $\theta_1 = \SI{40}{\celsius}$. On y introduit $m_g = \SI{50}{g}$ de glace à $\SI{0}{\celsius}$. Données : $c_{\text{eau}} = \SI{4185}{J.kg^{-1}.K^{-1}}$, $L_f = \SI{3,34e5}{J.kg^{-1}}$.
\begin{enumerate}
\item Vérifier que toute la glace fond.
\item Déterminer la température d'équilibre $\theta_{eq}$.
\end{enumerate}
\tcblower
\textbf{1. Vérification de l'hypothèse « toute la glace fond ».}
Chaleur maximale disponible (eau + calorimètre refroidis de \SI{40}{\celsius} à \SI{0}{\celsius}) :
\[ Q_{\text{dispo}} = (m_1+\mu)\,c_{\text{eau}}\,\theta_1 = \num{0,330}\times 4185\times 40 = \SI{55242}{J} \approx \SI{5,52e4}{J}. \]
Chaleur nécessaire pour fondre toute la glace :
\[ Q_{\text{fusion}} = m_g L_f = \num{0,050}\times \num{3,34e5} = \SI{16700}{J} = \SI{1,67e4}{J}. \]
$Q_{\text{dispo}} > Q_{\text{fusion}}$ \checkmark : toute la glace fond, $\theta_{eq} > \SI{0}{\celsius}$.

\textbf{2. Bilan thermique complet.}
L'eau chaude et le calorimètre cèdent ; la glace reçoit (fusion + chauffage eau de fusion) :
\[ (m_1+\mu)\,c_{\text{eau}}\,(\theta_{eq}-\theta_1) + m_g L_f + m_g\,c_{\text{eau}}\,(\theta_{eq}-0) = 0. \]
Masses en \si{kg} :
\[ \num{0,330}\times 4185\times(\theta_{eq}-40) + \num{16700} + \num{0,050}\times 4185\,\theta_{eq} = 0. \]
\[ \num{1381,05}\,\theta_{eq} - \num{55242} + \num{16700} + \num{209,25}\,\theta_{eq} = 0 \]
\[ \num{1590,3}\,\theta_{eq} = \num{38542} \quad\Longrightarrow\quad \theta_{eq} \approx \SI{24,23}{\celsius}. \]
\textbf{Vérification.} $0 < \theta_{eq} < \SI{40}{\celsius}$ \checkmark.
\textbf{Conclusion.} Toute la glace fond et la température finale du mélange est $\theta_{eq} \approx \SI{24,2}{\celsius}$.
\end{exoresolu}

\courssub{Synthèse et résolution de problèmes de calorimétrie}

\begin{methode}[Algorithme général de résolution (Probatoire)]
\begin{enumerate}
\item \textbf{Lire, lister, convertir} : $m$ en \si{kg}, $\theta$ en \si{\celsius}, $c$ en \si{J.kg^{-1}.K^{-1}}, L$ en \si{J.kg^{-1}}, \mu$ en \si{kg}.
\item \textbf{Schéma} : Dessiner le calorimètre, indiquer $\theta_i$ de chaque corps, $\theta_{eq}$ (inconnue ?).
\item \textbf{Phases} : Pour chaque corps, lister les étapes (chauffage/refroidissement $\to mc\Delta\theta$ ; changement d'état $\to mL$).
\item \textbf{Bilan} : Écrire $\sum Q_i = 0$ avec $Q_i = m_i c_i (\theta_{eq} - \theta_i)$ pour chaque étape sans changement d'état, et $Q_j = \pm m_j L_j$ pour chaque changement d'état. \textbf{Ne pas décider du signe à l'avance}, la formule $(\theta_{eq}-\theta_i)$ le donne.
\item \textbf{Résoudre} l'équation (linéaire en $\theta_{eq}$ ou en l'inconnue cherchée).
\item \textbf{Valider} : $\theta_{eq}$ entre les $\theta_i$ extrêmes ? $c>0$ ? Ordre de grandeur ? Unités ? Chiffres significatifs (2 ou 3) ?
\end{enumerate}
\end{methode}

\begin{attention}[Les 5 erreurs qui coûtent des points au Probatoire]
\begin{enumerate}
\item \textbf{Oublier la valeur en eau $\mu$ (ou la capacité $C_{\text{cal}}$).} Dès que l'énoncé donne $\mu$ ou $C_{\text{cal}}$, le vase participe aux échanges : $Q_{\text{cal}} = \mu\,c_{\text{eau}}\,\Delta\theta$ ou $C_{\text{cal}}\,\Delta\theta$. L'omettre fausse tout le bilan.
\item \textbf{Se tromper de signe dans le bilan.} \textbf{Astuce infaillible} : écris $\sum Q_i = 0$ avec $Q_i = m_i c_i(\theta_{eq}-\theta_i)$ pour \emph{chaque} corps/étape, sans te soucier de qui cède ou reçoit. Les signes sortent automatiquement. La méthode « chaleur cédée = chaleur reçue » (valeurs positives) est acceptée mais exige de prendre les $\Delta\theta$ en valeur absolue dans le bon sens.
\item \textbf{Utiliser $Q = mc\Delta\theta$ pendant un changement d'état.} Pendant la fusion/vaporisation, $\theta = \text{cste} \Rightarrow \Delta\theta = 0$. Il faut $Q = mL$. Réciproquement, n'utilise pas $mL$ hors changement d'état.
\item \textbf{Oublier le chauffage de l'eau de fusion.} Après la fonte de la glace à \SI{0}{\celsius}, l'eau obtenue continue de recevoir de la chaleur jusqu'à $\theta_{eq}$ : le terme $m_g c_{\text{eau}}\theta_{eq}$ est \textbf{obligatoire}.
\item \textbf{Négliger les unités et la vérification finale.} Convertis \textbf{toutes} les masses en \si{kg} (ou toutes en \si{g} avec $c$ en \si{J.g^{-1}.K^{-1}}, mais sois cohérent). Donne $Q$ en \si{J} (ou \si{kJ}) avec 2-3 chiffres significatifs. Vérifie \textbf{systématiquement} que $\theta_{eq}$ est comprise entre les températures initiales. Une $\theta_{eq}$ hors intervalle = erreur de calcul certaine.
\end{enumerate}
\end{attention}

\begin{exercice}[Entraînement 1 : Chauffage huile moteur]
On chauffe $m = \num{4,0}\ \si{kg}$ d'huile moteur ($c_{\text{huile}} = \SI{1900}{J.kg^{-1}.K^{-1}}$) de $\SI{20}{\celsius}$ à $\SI{100}{\celsius}$ dans un carter en aluminium ($m_{\text{Al}} = \num{2,0}\ \si{kg}$, $c_{\text{Al}} = \SI{897}{J.kg^{-1}.K^{-1}}$). Calculer la quantité de chaleur totale fournie. (Négliger les pertes).
\end{exercice}

\begin{exercice}[Entraînement 2 : Refroidissement pièce métallique]
Une pièce en acier ($c_{\text{Fe}} = \SI{460}{J.kg^{-1}.K^{-1}}$) de masse $\num{500}{g}$ à $\SI{300}{\celsius}$ est trempée dans un bac contenant $\num{10}{kg}$ d'eau à $\SI{20}{\celsius}$. Le bac a une capacité calorifique $C_{\text{bac}} = \SI{500}{J.K^{-1}}$. Déterminer la température d'équilibre $\theta_{eq}$.
\end{exercice}

\begin{exercice}[Entraînement 3 : Glace dans le jus]
Un verre (négligeable) contient $m_j = \SI{250}{g}$ de jus (assimilé à de l'eau, $c_{\text{eau}}$) à $\theta_j = \SI{25}{\celsius}$. On y ajoute $m_g = \SI{30}{g}$ de glace à $\SI{-10}{\celsius}$. Données : $c_{\text{glace}} = \SI{2090}{J.kg^{-1}.K^{-1}}$, $L_f = \SI{3,34e5}{J.kg^{-1}}$. Toute la glace fond-elle ? Si oui, $\theta_{eq}$ ? Sinon, masse de glace fondue ?
\end{exercice}

\begin{exercice}[Entraînement 4 : Valeur en eau par condensation]
Pour mesurer $\mu$ d'un calorimètre, on y introduit de la vapeur d'eau à $\SI{100}{\celsius}$ qui se condense. Initialement : $m_1 = \SI{200}{g}$ eau à $\SI{20}{\celsius}$. Après condensation de $m_v = \SI{10}{g}$ de vapeur, $\theta_{eq} = \SI{45}{\celsius}$. $L_v = \SI{2,26e6}{J.kg^{-1}}$. Calculer $\mu$.
\end{exercice}

\courssub{Corrigés des exercices d'entraînement}

\begin{corrige}[Exercice 1]
\textbf{Données :} $m_h = 4,0\,\si{kg}$, $c_h = 1900\,\si{J.kg^{-1}.K^{-1}}$ ; $m_{Al} = 2,0\,\si{kg}$, $c_{Al} = 897\,\si{J.kg^{-1}.K^{-1}}$ ; $\Delta\theta = 80\,\si{K}$.
\textbf{Bilan :} $Q_{\text{tot}} = Q_h + Q_{Al} = (m_h c_h + m_{Al} c_{Al})\Delta\theta$.
$Q_{\text{tot}} = (4,0\times 1900 + 2,0\times 897)\times 80 = (7600 + 1794)\times 80 = 9394\times 80 = \SI{751520}{J} \approx \SI{7,5e5}{J}$.
\textbf{Réponse :} $\SI{752}{kJ}$ fournis.
\end{corrige}

\begin{corrige}[Exercice 2]
\textbf{Inconnue :} $\theta_{eq}$. Masses en kg : $m_{Fe}=0,500$, $m_{eau}=10,0$. $C_{bac}=500\,\si{J.K^{-1}}$.
$\sum Q = 0 \Rightarrow m_{Fe}c_{Fe}(\theta_{eq}-300) + m_{eau}c_{eau}(\theta_{eq}-20) + C_{bac}(\theta_{eq}-20) = 0$.
$0,500\times 460(\theta_{eq}-300) + (10,0\times 4185 + 500)(\theta_{eq}-20) = 0$.
$230(\theta_{eq}-300) + 42350(\theta_{eq}-20) = 0$.
$230\theta_{eq} - 69000 + 42350\theta_{eq} - 847000 = 0$.
$42580\theta_{eq} = 916000 \Rightarrow \theta_{eq} \approx \SI{21,5}{\celsius}$.
\textbf{Vérification :} $20 < 21,5 < 300$ \checkmark. L'eau est très massive, $\theta_{eq}$ proche de $\theta_{eau}$.
\textbf{Réponse :} $\theta_{eq} \approx \SI{21,5}{\celsius}$.
\end{corrige}

\begin{corrige}[Exercice 3]
\textbf{1. Vérification fonte totale.}
$Q_{\text{dispo}}$ (jus de $25\to 0$) $= m_j c_{eau} \times 25 = 0,250\times 4185\times 25 = \SI{26156}{J}$.
$Q_{\text{nec}}$ (glace $-10\to 0$ + fusion) $= m_g c_{glace}\times 10 + m_g L_f = 0,030\times 2090\times 10 + 0,030\times 3,34e5 = 627 + 10020 = \SI{10647}{J}$.
$Q_{\text{dispo}} > Q_{\text{nec}}$ \checkmark $\to$ toute la glace fond, $\theta_{eq} > 0$.
\textbf{2. Bilan complet :}
$m_j c_{eau}(\theta_{eq}-25) + m_g c_{glace}(0 - (-10)) + m_g L_f + m_g c_{eau}(\theta_{eq}-0) = 0$.
$0,250\times 4185(\theta_{eq}-25) + 627 + 10020 + 0,030\times 4185\,\theta_{eq} = 0$.
$1046,25\theta_{eq} - 26156 + 10647 + 125,55\theta_{eq} = 0$.
$1171,8\theta_{eq} = 15509 \Rightarrow \theta_{eq} \approx \SI{13,2}{\celsius}$.
\textbf{Réponse :} Oui, toute la glace fond. $\theta_{eq} \approx \SI{13,2}{\celsius}$.
\end{corrige}

\begin{corrige}[Exercice 4]
\textbf{Acteurs :} Eau initiale ($m_1, \theta_1$), Calorimètre ($\mu, \theta_1$), Vapeur ($m_v, \SI{100}{\celsius}$).
La vapeur : 1) Condense à $\SI{100}{\celsius}$ ($Q_1 = -m_v L_v$), 2) Eau condensée refroidit $\num{100}\to \theta_{eq}$ ($Q_2 = m_v c_{eau}(\theta_{eq}-100)$).
Bilan $\sum Q = 0$ :
$m_1 c_{eau}(\theta_{eq}-\theta_1) + \mu c_{eau}(\theta_{eq}-\theta_1) - m_v L_v + m_v c_{eau}(\theta_{eq}-100) = 0$.
Masses en kg : $m_1=0,200$, $m_v=0,010$, $\theta_1=20$, $\theta_{eq}=45$.
$c_{eau}[(0,200+\mu)(45-20) + 0,010(45-100)] - 0,010\times 2,26e6 = 0$.
$4185[25(0,200+\mu) - 0,55] - 22600 = 0$.
$4185[5,0 + 25\mu - 0,55] = 22600$.
$4185[4,45 + 25\mu] = 22600$.
$4,45 + 25\mu = 22600/4185 \approx 5,400$.
$25\mu = 0,950 \Rightarrow \mu = \SI{0,038}{kg} = \SI{38}{g}$.
\textbf{Vérification :} $\mu > 0$ \checkmark.
\textbf{Réponse :} $\mu \approx \SI{38}{g}$.
\end{corrige}

\newpage

\coursec{Exercices}

\courssub{Je vérifie mes connaissances}

\begin{exercice}[QCM : les modes de transfert de chaleur]
Pour chaque situation, indiquer le mode de transfert de chaleur \textbf{dominant} : conduction (C), convection (V) ou rayonnement (R).
\begin{enumerate}
\item Le manche métallique d'une casserole posée sur le feu devient brûlant.
\item La chaleur du Soleil parvient jusqu'à la côte de Kribi.
\item L'eau chauffe uniformément dans une marmite placée sur un réchaud.
\item Une brise marine rafraîchit les habitants de Limbé en fin de journée.
\item Un mur en banco chauffé au soleil restitue sa chaleur à la pièce le soir.
\end{enumerate}
\end{exercice}

\begin{exercice}[Vrai ou faux justifié]
Répondre par vrai ou faux, en justifiant chaque réponse par une phrase.
\begin{enumerate}
\item La température d'un corps est une énergie.
\item La chaleur se transfère toujours spontanément du corps le plus chaud vers le corps le plus froid.
\item La chaleur latente de fusion de la glace s'exprime en \si{\joule\per\kilogram\per\kelvin}.
\item Deux corps en équilibre thermique sont nécessairement à la même température.
\item Une bouteille isotherme (thermos) supprime totalement les trois modes de transfert de chaleur.
\end{enumerate}
\end{exercice}

\begin{exercice}[Définitions et unités]
\begin{enumerate}
\item Définir la \cle{capacité thermique massique} (chaleur massique) d'un corps et donner son unité SI.
\item Définir la \cle{valeur en eau} d'un calorimètre.
\item Définir la \cle{chaleur latente} de changement d'état et donner son unité SI.
\item Écrire l'expression de la quantité de chaleur $Q$ reçue par un corps de masse $m$, de chaleur massique $c$, dont la température passe de $\theta_i$ à $\theta_f$ (sans changement d'état). Vérifier l'homogénéité de la relation par analyse dimensionnelle.
\end{enumerate}
\end{exercice}

\begin{exercice}[Calculs directs]
Données : $c_{\text{eau}} = \SI{4185}{\joule\per\kilogram\per\kelvin}$ ; $c_{\text{fer}} = \SI{460}{\joule\per\kilogram\per\kelvin}$ ; $L_f(\text{glace}) = \SI{3,34e5}{\joule\per\kilogram}$.
\begin{enumerate}
\item Calculer la quantité de chaleur nécessaire pour élever la température de $\SI{2,0}{kg}$ d'eau de $\SI{20}{\celsius}$ à $\SI{60}{\celsius}$.
\item Calculer la quantité de chaleur cédée par une barre de fer de masse $\SI{500}{g}$ qui se refroidit de $\SI{200}{\celsius}$ à $\SI{50}{\celsius}$.
\item Calculer la quantité de chaleur nécessaire pour fondre entièrement $\SI{250}{g}$ de glace prise à $\SI{0}{\celsius}$.
\end{enumerate}
\end{exercice}

\courssub{Je m'entraîne}

\begin{exercice}[Chauffer puis vaporiser de l'eau]
Une ménagère de Bafoussam veut faire bouillir $\SI{3}{L}$ d'eau prise à $\SI{25}{\celsius}$.
Données : $c_{\text{eau}} = \SI{4185}{\joule\per\kilogram\per\kelvin}$ ; $L_v(\text{eau}) = \SI{2,26e6}{\joule\per\kilogram}$ ; masse volumique de l'eau : $\rho = \SI{1,0}{kg/L}$.
\begin{enumerate}
\item Calculer la quantité de chaleur $Q_1$ nécessaire pour porter cette eau de $\SI{25}{\celsius}$ à $\SI{100}{\celsius}$.
\item Calculer la quantité de chaleur $Q_2$ nécessaire pour vaporiser entièrement cette eau à $\SI{100}{\celsius}$.
\item En déduire la chaleur totale $Q$ à fournir. Comparer $Q_1$ et $Q_2$ et commenter.
\end{enumerate}
\end{exercice}

\begin{exercice}[Bloc de fer plongé dans l'eau]
On plonge un bloc de fer de masse $m_1 = \SI{400}{g}$ sorti d'un four à $\theta_1 = \SI{150}{\celsius}$ dans une masse $m_2 = \SI{500}{g}$ d'eau à $\theta_2 = \SI{20}{\celsius}$, contenue dans un calorimètre de valeur en eau $\mu = \SI{25}{g}$.
Données : $c_{\text{eau}} = \SI{4185}{\joule\per\kilogram\per\kelvin}$ ; $c_{\text{fer}} = \SI{460}{\joule\per\kilogram\per\kelvin}$.
\begin{enumerate}
\item Faire le bilan des échanges de chaleur en précisant quel corps cède de la chaleur et quels corps en reçoivent.
\item Écrire l'équation traduisant le principe des échanges de chaleur.
\item Déterminer la température d'équilibre $\theta_e$ du mélange.
\end{enumerate}
\end{exercice}

\begin{exercice}[Glace dans l'eau chaude]
On introduit $m_g = \SI{100}{g}$ de glace à $\SI{0}{\celsius}$ dans $m_e = \SI{400}{g}$ d'eau à $\theta_e = \SI{30}{\celsius}$, dans un calorimètre de valeur en eau négligeable.
Données : $c_{\text{eau}} = \SI{4185}{\joule\per\kilogram\per\kelvin}$ ; $L_f = \SI{3,34e5}{\joule\per\kilogram}$.
\begin{enumerate}
\item Calculer la quantité de chaleur maximale que l'eau chaude peut céder en se refroidissant jusqu'à $\SI{0}{\celsius}$.
\item Calculer la quantité de chaleur nécessaire à la fusion totale de la glace. Toute la glace fond-elle ?
\item Déterminer la température finale $\theta_f$ du mélange.
\end{enumerate}
\end{exercice}

\begin{exercice}[Refroidissement d'un moteur de moto]
Le moteur d'une mototaxi de Douala dégage une chaleur de $\SI{8,0e4}{\joule}$ par minute de fonctionnement. On veut absorber cette chaleur avec de l'eau dont la température ne doit pas dépasser $\SI{90}{\celsius}$, l'eau étant initialement à $\SI{20}{\celsius}$.
Donnée : $c_{\text{eau}} = \SI{4185}{\joule\per\kilogram\per\kelvin}$.
\begin{enumerate}
\item Calculer la chaleur totale dégagée par le moteur en $10$ minutes de fonctionnement.
\item En déduire la masse minimale d'eau nécessaire pour absorber cette chaleur.
\item Proposer une raison pour laquelle, en pratique, on utilise un radiateur avec circulation d'air plutôt qu'un simple réservoir d'eau.
\end{enumerate}
\end{exercice}

\courssub{Je me prépare au Probatoire}

\begin{exercice}[Détermination de la valeur en eau d'un calorimètre]
Au laboratoire du lycée de Ngaoundéré, un groupe d'élèves de Première C veut déterminer la valeur en eau $\mu$ du calorimètre de l'établissement. Ils versent dans le calorimètre une masse $m_1 = \SI{200}{g}$ d'eau froide à $\theta_1 = \SI{18}{\celsius}$, puis ajoutent une masse $m_2 = \SI{250}{g}$ d'eau chaude à $\theta_2 = \SI{62}{\celsius}$. Après agitation, la température d'équilibre vaut $\theta_e = \SI{40,5}{\celsius}$.
Donnée : $c_{\text{eau}} = \SI{4185}{\joule\per\kilogram\per\kelvin}$.
\begin{enumerate}
\item Rappeler le rôle d'un calorimètre et citer deux de ses accessoires.
\item Faire le schéma annoté du calorimètre utilisé (vase intérieur, paroi isolante, couvercle, thermomètre, agitateur).
\item Exprimer la chaleur $Q_c$ cédée par l'eau chaude et la chaleur $Q_r$ reçue par l'ensemble \{eau froide + calorimètre\} en fonction des données.
\item En appliquant le principe des échanges de chaleur, établir l'expression littérale de $\mu$, puis calculer sa valeur en grammes.
\item Les élèves mesurent en réalité $\theta_e = \SI{40,5}{\celsius}$ alors que le calcul sans le calorimètre prévoyait une température plus élevée. Expliquer l'origine de cet écart.
\end{enumerate}
\end{exercice}

\begin{exercice}[Identification d'un métal par la méthode des mélanges]
Un bijoutier de Yaoundé veut identifier le métal constituant un échantillon de masse $m_m = \SI{150}{g}$. L'échantillon est chauffé dans un bain-marie jusqu'à $\theta_m = \SI{95}{\celsius}$, puis plongé rapidement dans un calorimètre de valeur en eau $\mu = \SI{30}{g}$ contenant $m_e = \SI{300}{g}$ d'eau à $\theta_i = \SI{22}{\celsius}$. La température d'équilibre se stabilise à $\theta_f = \SI{27,4}{\celsius}$.
Données : $c_{\text{eau}} = \SI{4185}{\joule\per\kilogram\per\kelvin}$. Chaleurs massiques de quelques métaux (en \si{\joule\per\kilogram\per\kelvin}) : aluminium $900$ ; fer $460$ ; cuivre $385$ ; plomb $130$.
\begin{enumerate}
\item Indiquer deux précautions expérimentales à prendre pour limiter les pertes thermiques pendant l'opération.
\item Écrire le bilan des échanges de chaleur entre l'échantillon, l'eau et le calorimètre.
\item Établir l'expression littérale de la chaleur massique $c_m$ du métal.
\item Calculer $c_m$ et identifier le métal constituant l'échantillon.
\item La valeur trouvée est légèrement inférieure à la valeur tabulée. Proposer une explication.
\end{enumerate}
\end{exercice}

\begin{exercice}[Préparation d'un bain tiède à l'hôpital]
Dans un centre de santé de l'Extrême-Nord, on prépare un bain tiède en mélangeant de l'eau chaude et de la glace fondante. On dispose de $m_1 = \SI{2,0}{kg}$ d'eau à $\theta_1 = \SI{70}{\celsius}$ dans une cuve supposée adiabatique (valeur en eau négligeable), et on y introduit une masse $m_2$ de glace à $\SI{0}{\celsius}$. On souhaite obtenir un mélange final à $\theta_f = \SI{37}{\celsius}$ (température du corps humain).
Données : $c_{\text{eau}} = \SI{4185}{\joule\per\kilogram\per\kelvin}$ ; $L_f = \SI{3,34e5}{\joule\per\kilogram}$.
\begin{enumerate}
\item Rappeler les trois modes de transfert de chaleur et préciser celui qui est supprimé par les parois adiabatiques de la cuve.
\item Calculer la quantité de chaleur $Q_1$ que l'eau chaude doit céder pour passer de $\SI{70}{\celsius}$ à $\SI{37}{\celsius}$.
\item Exprimer, en fonction de $m_2$, la quantité de chaleur $Q_2$ reçue par la glace pour fondre totalement à $\SI{0}{\celsius}$, puis pour que l'eau de fusion atteigne $\SI{37}{\celsius}$.
\item En appliquant le principe des échanges de chaleur, calculer la masse $m_2$ de glace nécessaire.
\item Vérifier a posteriori que l'hypothèse « toute la glace fond » est cohérente avec le résultat obtenu.
\item En l'absence de glace, quelle masse d'eau froide à $\SI{15}{\celsius}$ faudrait-il mélanger aux $\SI{2,0}{kg}$ d'eau à $\SI{70}{\celsius}$ pour obtenir la même température finale de $\SI{37}{\celsius}$ ?
\end{enumerate}
\end{exercice}

\newpage

\coursec{Corrigés des exercices}

\begin{corrige}[Exercice 1 — QCM : les modes de transfert de chaleur]
\begin{enumerate}
\item \textbf{Conduction.} Le manche métallique est en contact direct avec la casserole chaude ; la chaleur se propage à l'intérieur du métal par agitation des électrons libres et vibration des ions (conduction thermique).
\item \textbf{Rayonnement.} Le vide interplanétaire ne permet ni conduction ni convection ; l'énergie solaire voyage sous forme d'ondes électromagnétiques (rayonnement infrarouge et visible).
\item \textbf{Convection naturelle.} L'eau au fond de la marmite se dilate, devient moins dense et monte, tandis que l'eau plus froide descend : c'est un transfert de chaleur par mouvement de matière (convection naturelle).
\item \textbf{Convection naturelle.} La brise marine est un mouvement d'air (fluide) plus frais venant de la mer vers la terre, dû à une différence de température ; elle transporte de l'énergie thermique par convection naturelle (et non forcée, car aucun dispositif mécanique ne crée le flux).
\item \textbf{Rayonnement.} Le mur en banco, une fois chauffé, émet un rayonnement infrarouge vers l'intérieur de la pièce ; l'air étant un mauvais conducteur, le rayonnement domine l'échange avec les parois et les occupants.
\end{enumerate}
\end{corrige}

\begin{corrige}[Exercice 2 — Vrai ou faux justifié]
\begin{enumerate}
\item \textbf{Faux.} La température est une grandeur intensive qui caractérise l'état thermique d'un corps ; elle n'est pas une énergie. L'énergie thermique interne est une énergie, mesurée en joules (\si{\joule}).
\item \textbf{Vrai.} C'est une conséquence du second principe de la thermodynamique : les échanges thermiques spontanés se font toujours du corps le plus chaud vers le corps le plus froid.
\item \textbf{Faux.} La chaleur latente de fusion s'exprime en \si{\joule\per\kilogram} (\si{J/kg}). L'unité \si{\joule\per\kilogram\per\kelvin} (\si{J/(kg\cdot K)}) correspond à la capacité thermique massique.
\item \textbf{Vrai.} L'équilibre thermique entre deux systèmes signifie qu'ils ont la même température ; sinon, un flux de chaleur persisterait.
\item \textbf{Faux.} Un thermos \emph{réduit fortement} les trois modes (paroi à double vitrage sous vide pour la conduction et la convection, paroi argentée pour le rayonnement), mais ne les supprime pas totalement : des pertes résiduelles persistent (par le bouchon, par rayonnement imparfait, etc.).
\end{enumerate}
\end{corrige}

\begin{corrige}[Exercice 3 — Définitions et unités]
\begin{enumerate}
\item La \cle{capacité thermique massique} (ou \cle{chaleur massique}) $c$ d'un corps est la quantité de chaleur qu'il faut fournir à \SI{1}{kg} de ce corps pour élever sa température de \SI{1}{K} (ou \SI{1}{\degreeCelsius}). Son unité SI est le \si{\joule\per\kilogram\per\kelvin} (\si{J/(kg\cdot K)}).
\item La \cle{valeur en eau} $\mu$ d'un calorimètre est la masse d'eau qui aurait la même capacité calorifique que le calorimètre (vase + accessoires). Elle s'exprime en kilogrammes (\si{kg}) ou en grammes (\si{g}).
\item La \cle{chaleur latente} $L$ de changement d'état est la quantité de chaleur nécessaire pour faire changer de phase \SI{1}{kg} d'une substance pure à température constante. Son unité SI est le \si{\joule\per\kilogram} (\si{J/kg}).
\item La quantité de chaleur reçue (sans changement d'état) est $Q = m\,c\,(\theta_f - \theta_i)$.
  \textbf{Analyse dimensionnelle :}
  \[
  [Q] = \si{J},\quad [m] = \si{kg},\quad [c] = \si{J/(kg\cdot K)},\quad [\theta_f-\theta_i] = \si{K}.
  \]
  \[
  [m\,c\,\Delta\theta] = \si{kg} \times \si{J/(kg\cdot K)} \times \si{K} = \si{J} = [Q].
  \]
  La relation est homogène.
\end{enumerate}
\end{corrige}

\begin{corrige}[Exercice 4 — Calculs directs]
Données : $c_{\text{eau}} = \SI{4185}{J/(kg\cdot K)}$, $c_{\text{fer}} = \SI{460}{J/(kg\cdot K)}$, $L_f = \SI{3,34e5}{J/kg}$.
\begin{enumerate}
\item $m = \SI{2,0}{kg}$, $\Delta\theta = \SI{60}{\degreeCelsius} - \SI{20}{\degreeCelsius} = \SI{40}{K}$.
  \[
  Q = m\,c_{\text{eau}}\,\Delta\theta = 2,0 \times 4185 \times 40 = \SI{334800}{J} \approx \SI{3,35e5}{J}.
  \]
\item $m = \SI{500}{g} = \SI{0,500}{kg}$, $\Delta\theta = \SI{50}{\degreeCelsius} - \SI{200}{\degreeCelsius} = -\SI{150}{K}$.
  \[
  Q = m\,c_{\text{fer}}\,\Delta\theta = 0,500 \times 460 \times (-150) = -\SI{34500}{J}.
  \]
  La barre de fer \textbf{cède} une chaleur de \SI{3,45e4}{J} (valeur absolue).
\item $m = \SI{250}{g} = \SI{0,250}{kg}$, glace à \SI{0}{\degreeCelsius} $\rightarrow$ eau à \SI{0}{\degreeCelsius}.
  \[
  Q = m\,L_f = 0,250 \times 3,34 \times 10^5 = \SI{83500}{J} = \SI{8,35e4}{J}.
  \]
\end{enumerate}
\end{corrige}

\begin{corrige}[Exercice 5 — Chauffer puis vaporiser de l'eau]
Données : $V = \SI{3}{L}$, $\rho = \SI{1,0}{kg/L} \Rightarrow m = \SI{3,0}{kg}$ ; $c = \SI{4185}{J/(kg\cdot K)}$ ; $L_v = \SI{2,26e6}{J/kg}$.
\begin{enumerate}
\item Échauffement de \SI{25}{\degreeCelsius} à \SI{100}{\degreeCelsius} : $\Delta\theta = \SI{75}{K}$.
  \[
  Q_1 = m\,c\,\Delta\theta = 3,0 \times 4185 \times 75 = \SI{941625}{J} \approx \SI{9,42e5}{J}.
  \]
\item Vaporisation à \SI{100}{\degreeCelsius} :
  \[
  Q_2 = m\,L_v = 3,0 \times 2,26 \times 10^6 = \SI{6,78e6}{J}.
  \]
\item Chaleur totale :
  \[
  Q = Q_1 + Q_2 = 9,41625 \times 10^5 + 6,78 \times 10^6 = \SI{7,72e6}{J}.
  \]
  \textbf{Comparaison :} $Q_2 \approx 7,2 \times Q_1$. La vaporisation demande beaucoup plus d'énergie que l'échauffement car il faut vaincre les forces cohésives entre molécules pour passer de l'état liquide à l'état gazeux (travail de dilatation contre la pression atmosphérique + augmentation de l'énergie potentielle intermoléculaire).
\end{enumerate}
\end{corrige}

\begin{corrige}[Exercice 6 — Bloc de fer plongé dans l'eau]
Données : $m_1 = \SI{400}{g} = \SI{0,400}{kg}$, $\theta_1 = \SI{150}{\degreeCelsius}$ ; $m_2 = \SI{500}{g} = \SI{0,500}{kg}$, $\theta_2 = \SI{20}{\degreeCelsius}$ ; $\mu = \SI{25}{g} = \SI{0,025}{kg}$ ; $c_{\text{eau}} = \SI{4185}{J/(kg\cdot K)}$, $c_{\text{fer}} = \SI{460}{J/(kg\cdot K)}$.
\begin{enumerate}
\item \textbf{Bilan des échanges :}
  \begin{itemize}
  \item Le bloc de \textbf{fer cède} de la chaleur (il se refroidit).
  \item L'\textbf{eau} et le \textbf{calorimètre} (sa valeur en eau) \textbf{reçoivent} de la chaleur (ils s'échauffent).
  \end{itemize}
\item \textbf{Équation du principe des échanges de chaleur} (système isolé) :
  \[
  Q_{\text{fer}} + Q_{\text{eau}} + Q_{\text{cal}} = 0
  \]
  \[
  m_1 c_{\text{fer}} (\theta_e - \theta_1) + m_2 c_{\text{eau}} (\theta_e - \theta_2) + \mu c_{\text{eau}} (\theta_e - \theta_2) = 0.
  \]
\item \textbf{Résolution pour $\theta_e$ :}
  \begin{align*}
  \theta_e (m_1 c_{\text{fer}} + m_2 c_{\text{eau}} + \mu c_{\text{eau}}) &= m_1 c_{\text{fer}} \theta_1 + (m_2 + \mu) c_{\text{eau}} \theta_2 \\
  \theta_e &= \frac{m_1 c_{\text{fer}} \theta_1 + (m_2 + \mu) c_{\text{eau}} \theta_2}{m_1 c_{\text{fer}} + (m_2 + \mu) c_{\text{eau}}}.
  \end{align*}
  Calcul des termes :
  \begin{align*}
  m_1 c_{\text{fer}} &= 0,400 \times 460 = \SI{184}{J/K}, \\
  m_2 c_{\text{eau}} &= 0,500 \times 4185 = \SI{2092,5}{J/K}, \\
  \mu c_{\text{eau}} &= 0,025 \times 4185 = \SI{104,625}{J/K}, \\
  \text{Numérateur} &= 184 \times 150 + (2092,5 + 104,625) \times 20 \\
  &= 27600 + 2197,125 \times 20 = 27600 + 43942,5 = \SI{71542,5}{J}, \\
  \text{Dénominateur} &= 184 + 2092,5 + 104,625 = \SI{2381,125}{J/K}, \\
  \theta_e &= \frac{71542,5}{2381,125} \approx \SI{30,0}{\degreeCelsius}.
  \end{align*}
  La température d'équilibre est \textbf{$\theta_e \approx \SI{30,0}{\degreeCelsius}$}.
\end{enumerate}
\end{corrige}

\begin{corrige}[Exercice 7 — Glace dans l'eau chaude]
Données : $m_g = \SI{100}{g} = \SI{0,100}{kg}$ (glace à \SI{0}{\degreeCelsius}) ; $m_e = \SI{400}{g} = \SI{0,400}{kg}$ (eau à \SI{30}{\degreeCelsius}) ; calorimètre négligeable ; $c = \SI{4185}{J/(kg\cdot K)}$ ; $L_f = \SI{3,34e5}{J/kg}$.
\begin{enumerate}
\item Chaleur maximale que l'eau chaude peut céder en refroidissant jusqu'à \SI{0}{\degreeCelsius} :
  \[
  Q_{\text{max}} = m_e c (30 - 0) = 0,400 \times 4185 \times 30 = \SI{50220}{J}.
  \]
\item Chaleur nécessaire pour fondre toute la glace :
  \[
  Q_{\text{fusion}} = m_g L_f = 0,100 \times 3,34 \times 10^5 = \SI{33400}{J}.
  \]
  Comme $Q_{\text{max}} > Q_{\text{fusion}}$, \textbf{toute la glace fond}.
\item Après fusion, on a $m_{\text{total}} = m_e + m_g = \SI{0,500}{kg}$ d'eau. Soit $\theta_f$ la température finale.
  Bilan : chaleur cédée par l'eau initialement chaude = chaleur absorbée par la glace (fusion + échauffement de l'eau de fusion).
  \[
  m_e c (30 - \theta_f) = m_g L_f + m_g c (\theta_f - 0).
  \]
  \begin{align*}
  0,400 \times 4185 \times (30 - \theta_f) &= 33400 + 0,100 \times 4185 \times \theta_f \\
  1674 (30 - \theta_f) &= 33400 + 418,5 \theta_f \\
  50220 - 1674 \theta_f &= 33400 + 418,5 \theta_f \\
  50220 - 33400 &= (1674 + 418,5) \theta_f \\
  16820 &= 2092,5 \theta_f \\
  \theta_f &= \frac{16820}{2092,5} \approx \SI{8,04}{\degreeCelsius}.
  \end{align*}
  La température finale du mélange est \textbf{$\theta_f \approx \SI{8,0}{\degreeCelsius}$}.
\end{enumerate}
\end{corrige}

\begin{corrige}[Exercice 8 — Refroidissement d'un moteur de moto]
Données : puissance thermique $P = \SI{8,0e4}{J/min}$ ; $\Delta t = \SI{10}{min}$ ; $c = \SI{4185}{J/(kg\cdot K)}$ ; $\theta_i = \SI{20}{\degreeCelsius}$, $\theta_{\text{max}} = \SI{90}{\degreeCelsius}$.
\begin{enumerate}
\item Chaleur totale dégagée en 10 min :
  \[
  Q_{\text{tot}} = P \times \Delta t = 8,0 \times 10^4 \times 10 = \SI{8,0e5}{J}.
  \]
\item Masse minimale d'eau pour absorber cette chaleur sans dépasser \SI{90}{\degreeCelsius} :
  \[
  Q_{\text{tot}} = m_{\text{eau}} c (\theta_{\text{max}} - \theta_i) \Rightarrow m_{\text{eau}} = \frac{Q_{\text{tot}}}{c \Delta\theta} = \frac{8,0 \times 10^5}{4185 \times 70}.
  \]
  \[
  m_{\text{eau}} = \frac{800000}{292950} \approx \SI{2,73}{kg}.
  \]
  Il faut au minimum \textbf{$\approx \SI{2,7}{kg}$ d'eau} (2 chiffres significatifs imposés par les données).
\item \textbf{Raison pratique :} Une simple réserve d'eau de \SI{2,7}{kg} atteindrait rapidement \SI{90}{\degreeCelsius} et ne pourrait plus absorber de chaleur ; le moteur surchaufferait. Un radiateur avec circulation d'air (convection forcée) permet d'évacuer la chaleur en continu vers l'atmosphère, maintenant l'eau à une température plus basse et limitant le volume d'eau nécessaire. De plus, l'ébullition de l'eau serait évitée.
\end{enumerate}
\end{corrige}

\begin{corrige}[Exercice 9 — Détermination de la valeur en eau d'un calorimètre]
Données : $m_1 = \SI{200}{g} = \SI{0,200}{kg}$, $\theta_1 = \SI{18}{\degreeCelsius}$ ; $m_2 = \SI{250}{g} = \SI{0,250}{kg}$, $\theta_2 = \SI{62}{\degreeCelsius}$ ; $\theta_e = \SI{40,5}{\degreeCelsius}$ ; $c = \SI{4185}{J/(kg\cdot K)}$.
\begin{enumerate}
\item \textbf{Rôle :} Le calorimètre permet de réaliser des mesures de quantités de chaleur en limitant les échanges thermiques avec le milieu extérieur (système quasi adiabatique).
  \textbf{Accessoires :} thermomètre, agitateur (pour uniformiser la température), couvercle isolant, paroi à double paroi (vide d'air ou liège).
\item \textbf{Schéma annoté :} (Description) Vase intérieur métallique (cuve) contenant l'eau, entouré d'une paroi isolante (vide ou matériau poreux), couvercle percé pour le thermomètre et l'agitateur, thermomètre plongeant dans l'eau, agitateur pour homogénéiser la température.
\item Chaleur cédée par l'eau chaude (indice 2) :
  \[
  Q_c = m_2 c (\theta_2 - \theta_e).
  \]
  Chaleur reçue par l'ensemble \{eau froide + calorimètre\} :
  \[
  Q_r = m_1 c (\theta_e - \theta_1) + \mu c (\theta_e - \theta_1) = (m_1 + \mu) c (\theta_e - \theta_1).
  \]
\item Principe des échanges : $Q_c = Q_r$.
  \[
  m_2 c (\theta_2 - \theta_e) = (m_1 + \mu) c (\theta_e - \theta_1).
  \]
  Simplifier $c$ :
  \[
  m_2 (\theta_2 - \theta_e) = (m_1 + \mu) (\theta_e - \theta_1).
  \]
  Expression littérale de $\mu$ :
  \[
  \mu = \frac{m_2 (\theta_2 - \theta_e)}{\theta_e - \theta_1} - m_1.
  \]
  Calcul numérique :
  \begin{align*}
  \mu &= \frac{0,250 \times (62 - 40,5)}{40,5 - 18} - 0,200 \\
      &= \frac{0,250 \times 21,5}{22,5} - 0,200 \\
      &= \frac{5,375}{22,5} - 0,200 \\
      &= 0,2389 - 0,200 = \SI{0,0389}{kg} = \SI{38,9}{g}.
  \end{align*}
  La valeur en eau du calorimètre est \textbf{$\mu \approx \SI{38,9}{g}$}.
\item Sans tenir compte du calorimètre ($\mu = 0$), la température d'équilibre prévue serait plus élevée car toute la chaleur cédée par l'eau chaude irait uniquement à l'eau froide. La mesure réelle ($\theta_e = \SI{40,5}{\degreeCelsius}$) est \textbf{plus basse} car le calorimètre absorbe une partie de la chaleur, ce qui abaisse la température finale. L'écart observé confirme que le calorimètre a une capacité thermique non nulle.
\end{enumerate}
\textbf{Barème indicatif (sur 6 pts) :}
\begin{itemize}
\item Rôle et accessoires : 1 pt.
\item Schéma annoté (au moins 4 éléments) : 1 pt.
\item Expressions de $Q_c$ et $Q_r$ : 1 pt.
\item Équation, expression littérale et calcul de $\mu$ : 2 pts.
\item Explication de l'écart : 1 pt.
\end{itemize}
\textbf{Points de vigilance :} Ne pas oublier la valeur en eau dans le bilan ; bien distinguer $\theta_e - \theta_1$ (réchauffement) et $\theta_2 - \theta_e$ (refroidissement) ; exprimer $\mu$ en kg puis convertir en g ; justifier le signe de l'écart.
\end{corrige}

\begin{corrige}[Exercice 10 — Identification d'un métal par la méthode des mélanges]
Données : $m_m = \SI{150}{g} = \SI{0,150}{kg}$, $\theta_m = \SI{95}{\degreeCelsius}$ ; $\mu = \SI{30}{g} = \SI{0,030}{kg}$ ; $m_e = \SI{300}{g} = \SI{0,300}{kg}$, $\theta_i = \SI{22}{\degreeCelsius}$ ; $\theta_f = \SI{27,4}{\degreeCelsius}$ ; $c_{\text{eau}} = \SI{4185}{J/(kg\cdot K)}$.
\begin{enumerate}
\item \textbf{Précautions :}
  \begin{itemize}
  \item Transférer l'échantillon du bain-marie au calorimètre \textbf{le plus rapidement possible} pour limiter les pertes de chaleur vers l'air.
  \item Bien \textbf{essuyer} l'échantillon pour éviter d'introduire de l'eau supplémentaire qui fausserait les masses.
  \item Refermer immédiatement le calorimètre et \textbf{agiter} doucement pour uniformiser la température sans trop échanger avec l'extérieur.
  \end{itemize}
\item \textbf{Bilan des échanges} (système isolé) :
  Chaleur cédée par le métal = Chaleur reçue par l'eau + chaleur reçue par le calorimètre.
  \[
  m_m c_m (\theta_m - \theta_f) = m_e c_{\text{eau}} (\theta_f - \theta_i) + \mu c_{\text{eau}} (\theta_f - \theta_i).
  \]
\item \textbf{Expression littérale de $c_m$ :}
  \[
  c_m = \frac{(m_e + \mu) c_{\text{eau}} (\theta_f - \theta_i)}{m_m (\theta_m - \theta_f)}.
  \]
\item \textbf{Calcul numérique :}
  \begin{align*}
  m_e + \mu &= 0,300 + 0,030 = \SI{0,330}{kg}, \\
  \theta_f - \theta_i &= 27,4 - 22 = \SI{5,4}{K}, \\
  \theta_m - \theta_f &= 95 - 27,4 = \SI{67,6}{K}, \\
  \text{Numérateur} &= 0,330 \times 4185 \times 5,4 = \SI{7457,7}{J}, \\
  \text{Dénominateur} &= 0,150 \times 67,6 = \SI{10,14}{kg\cdot K}, \\
  c_m &= \frac{7457,7}{10,14} \approx \SI{735}{J/(kg\cdot K)}.
  \end{align*}
  La chaleur massique trouvée est \textbf{$c_m \approx \SI{735}{J/(kg\cdot K)}$}.
  \textbf{Identification :} Cette valeur ne correspond exactement à aucun des métaux purs usuels. Elle est la plus proche de l'\textbf{aluminium} (\SI{900}{J/(kg\cdot K)}), mais reste significativement plus faible. L'échantillon pourrait être un alliage d'aluminium ou présenter des impuretés.
\item \textbf{Explication de l'écart :} La valeur expérimentale est \textbf{inférieure} à la valeur tabulée de l'aluminium. Cela s'explique par des \textbf{pertes thermiques} vers l'extérieur pendant le transfert et la manipulation : une partie de la chaleur de l'échantillon est perdue dans l'air au lieu d'être transmise à l'eau, ce qui fait sous-estimer $c_m$. D'autres facteurs : isolation imparfaite du calorimètre, erreurs de mesure des masses ou des températures.
\end{enumerate}
\textbf{Barème indicatif (sur 7 pts) :}
\begin{itemize}
\item Précautions (2 pertinentes) : 1 pt.
\item Bilan des échanges : 1 pt.
\item Expression littérale de $c_m$ : 1 pt.
\item Calcul correct de $c_m$ : 2 pts.
\item Identification et commentaire sur l'écart : 2 pts.
\end{itemize}
\textbf{Points de vigilance :} Utiliser les masses en kg (ou rester cohérent en g avec $c$ en J/(g·K)) ; ne pas inverser $\theta_m - \theta_f$ et $\theta_f - \theta_i$ ; citer explicitement les pertes thermiques comme cause d'une valeur mesurée inférieure à la valeur théorique.
\end{corrige}

\begin{corrige}[Exercice 11 — Préparation d'un bain tiède à l'hôpital]
Données : $m_1 = \SI{2,0}{kg}$, $\theta_1 = \SI{70}{\degreeCelsius}$ ; glace à $\theta_g = \SI{0}{\degreeCelsius}$ ; $\theta_f = \SI{37}{\degreeCelsius}$ ; cuve adiabatique ($\mu \approx 0$) ; $c = \SI{4185}{J/(kg\cdot K)}$ ; $L_f = \SI{3,34e5}{J/kg}$.
\begin{enumerate}
\item \textbf{Trois modes :} conduction, convection, rayonnement.
  Les parois adiabatiques de la cuve suppriment (ou rendent négligeable) le mode \textbf{conduction} à travers les parois (elles sont isolantes). Elles limitent aussi la convection et le rayonnement vers l'extérieur, mais la conduction est le mode directement bloqué par un isolant thermique.
\item Chaleur cédée par l'eau chaude pour passer de \SI{70}{\degreeCelsius} à \SI{37}{\degreeCelsius} :
  \[
  Q_1 = m_1 c (\theta_1 - \theta_f) = 2,0 \times 4185 \times (70 - 37) = 2,0 \times 4185 \times 33 = \SI{276210}{J} \approx \SI{2,76e5}{J}.
  \]
\item La glace reçoit de la chaleur en deux étapes :
  \begin{enumerate}
  \item Fusion à \SI{0}{\degreeCelsius} : $Q_{\text{fusion}} = m_2 L_f$.
  \item Échauffement de l'eau de fusion de \SI{0}{\degreeCelsius} à \SI{37}{\degreeCelsius} : $Q_{\text{échauff}} = m_2 c (\theta_f - 0)$.
  \end{enumerate}
  Total :
  \[
  Q_2 = m_2 L_f + m_2 c \theta_f = m_2 (L_f + c \theta_f).
  \]
\item Principe des échanges : $Q_1 = Q_2$.
  \[
  m_1 c (\theta_1 - \theta_f) = m_2 (L_f + c \theta_f).
  \]
  \[
  m_2 = \frac{m_1 c (\theta_1 - \theta_f)}{L_f + c \theta_f}.
  \]
  Calcul :
  \begin{align*}
  c \theta_f &= 4185 \times 37 = \SI{154845}{J/kg}, \\
  L_f + c \theta_f &= 334000 + 154845 = \SI{488845}{J/kg}, \\
  m_1 c (\theta_1 - \theta_f) &= 2,0 \times 4185 \times 33 = \SI{276210}{J}, \\
  m_2 &= \frac{276210}{488845} \approx \SI{0,565}{kg} = \SI{565}{g}.
  \end{align*}
  Il faut \textbf{$m_2 \approx \SI{565}{g}$ de glace}.
\item \textbf{Vérification a posteriori :}
  Chaleur nécessaire pour fondre toute la glace : $m_2 L_f = 0,565 \times 334000 = \SI{188710}{J}$.
  Chaleur maximale que l'eau chaude peut céder en allant jusqu'à \SI{0}{\degreeCelsius} : $m_1 c (70 - 0) = 2,0 \times 4185 \times 70 = \SI{585900}{J}$.
  Comme $\SI{585900}{J} > \SI{188710}{J}$, l'eau chaude fournit largement assez de chaleur pour fondre toute la glace. De plus, la température finale \SI{37}{\degreeCelsius} est $> \SI{0}{\degreeCelsius}$, ce qui confirme que tout est fondu. L'hypothèse est \textbf{cohérente}.
\item Sans glace, on mélange $m_1 = \SI{2,0}{kg}$ à \SI{70}{\degreeCelsius} avec $m_3$ d'eau à \SI{15}{\degreeCelsius} pour obtenir \SI{37}{\degreeCelsius}.
  \[
  m_1 c (70 - 37) = m_3 c (37 - 15) \Rightarrow m_1 \times 33 = m_3 \times 22.
  \]
  \[
  m_3 = m_1 \frac{33}{22} = 2,0 \times 1,5 = \SI{3,0}{kg}.
  \]
  Il faudrait \textbf{$\SI{3,0}{kg}$ d'eau à \SI{15}{\degreeCelsius}}.
\end{enumerate}
\textbf{Barème indicatif (sur 8 pts) :}
\begin{itemize}
\item Modes de transfert et rôle paroi adiabatique : 1 pt.
\item Calcul de $Q_1$ : 1 pt.
\item Expression de $Q_2$ en fonction de $m_2$ : 1 pt.
\item Équation, expression littérale et calcul de $m_2$ : 2 pts.
\item Vérification de la cohérence (fusion totale) : 1 pt.
\item Calcul de la masse d'eau froide alternative : 1 pt.
\item Rigueur des unités, nombres significatifs, rédaction : 1 pt.
\end{itemize}
\textbf{Points de vigilance :} Ne pas oublier l'échauffement de l'eau de fusion après la fonte ; vérifier que la température finale est bien supérieure à \SI{0}{\degreeCelsius} pour valider l'hypothèse ; bien distinguer les différences de température en Kelvin (identiques aux degrés Celsius pour les écarts) ; exprimer les masses en kg pour utiliser $c$ en SI.
\end{corrige}

\newpage

\coursec{Fiche bilan}

\begin{popfigure}
\begin{tikzpicture}[mindmap, grow cyclic, every node/.style={concept}, text width=2.6cm, align=center,
  root concept/.append style={concept color=popblue, font=\large\bfseries, text=white, text width=3.2cm},
  level 1 concept/.append style={level distance=4.2cm, sibling angle=51.4, font=\small\bfseries},
  level 2 concept/.append style={level distance=2.6cm, font=\scriptsize, text width=2.2cm}]
\node[root concept]{Transferts de chaleur et calorimétrie}
  child[concept color=poporange]{ node {Chaleur et température}
    child[concept color=poporangeL]{ node[align=center]{$Q$ en joule (J)\\ $\theta$ en \si{\celsius}, $T$ en K} }
    child[concept color=poporangeL]{ node {Sens : chaud $\to$ froid} }
  }
  child[concept color=poppurple]{ node {3 modes de transfert}
    child[concept color=poppurpleL]{ node {Conduction (solides)} }
    child[concept color=poppurpleL]{ node {Convection (fluides)} }
    child[concept color=poppurpleL]{ node {Rayonnement (sans support)} }
  }
  child[concept color=popgreen]{ node {Sans changement d'état}
    child[concept color=popgreenL]{ node {$Q=mc\,\Delta\theta$} }
    child[concept color=popgreenL]{ node {$C=mc$ (J/K)} }
  }
  child[concept color=poppink]{ node {Changement d'état}
    child[concept color=poppinkL]{ node {$Q=mL$} }
    child[concept color=poppinkL]{ node[align=center]{$L$ en J/kg\\ palier de $\theta$} }
  }
  child[concept color=popgold]{ node {Calorimètre}
    child[concept color=popgoldL]{ node[align=center]{Vase Dewar\\ isolant} }
    child[concept color=popgoldL]{ node {Valeur en eau $\mu$} }
  }
  child[concept color=popblue!70]{ node {Principe des échanges}
    child[concept color=popblueL]{ node {$Q_{\text{cédée}}+Q_{\text{re\c{c}ue}}=0$} }
    child[concept color=popblueL]{ node {Système isolé : $\sum Q_i=0$} }
  };
\end{tikzpicture}
\legende{Carte mentale de la leçon : les six idées forces de la calorimétrie.}
\end{popfigure}

\begin{bilan}
\textbf{Définitions clés.}
\begin{itemize}
  \item \cle{Température} : grandeur liée à l'agitation microscopique ; se mesure au thermomètre. \cle{Chaleur} $Q$ : énergie transférée spontanément du corps chaud vers le corps froid (unité : le joule, J).
  \item \cle{Conduction} : transfert de proche en proche, sans transport de matière (métaux). \cle{Convection} : transfert avec déplacement de matière (fluides). \cle{Rayonnement} : transfert par ondes électromagnétiques, possible dans le vide.
  \item \cle{Capacité thermique massique} (chaleur massique) $c$ : chaleur nécessaire pour élever de \SI{1}{\celsius} (ou \SI{1}{K}) la température de \SI{1}{kg} de corps ; unité : \si{J.kg^{-1}.K^{-1}}. Pour l'eau : $c_{\text{eau}} = \SI{4185}{J.kg^{-1}.K^{-1}}$.
  \item \cle{Capacité calorifique} $C = mc$ : chaleur nécessaire pour élever de \SI{1}{K} la température du corps entier ; unité : \si{J.K^{-1}}.
  \item \cle{Chaleur latente} $L$ : chaleur échangée par unité de masse lors d'un changement d'état, à température constante ; unité : \si{J.kg^{-1}}. $L_{\text{fusion}} > 0$ et $L_{\text{solidification}} = -L_{\text{fusion}}$.
  \item \cle{Valeur en eau} $\mu$ du calorimètre : masse d'eau ayant la même capacité calorifique que le calorimètre et ses accessoires : $C_{\text{cal}} = \mu\, c_{\text{eau}}$.
\end{itemize}

\textbf{Formules à connaître par c\oe ur.}
\begin{center}
\begin{tabularx}{\linewidth}{|l|X|l|}\hline
\rowcolor{popblueL} \textbf{Grandeur} & \textbf{Formule} & \textbf{Unités SI} \\ \hline
Chaleur sensible & $Q = mc\,(\theta_f - \theta_i) = C\,\Delta\theta$ & J ; kg ; \si{J.kg^{-1}.K^{-1}} \\ \hline
Chaleur latente & $Q = mL$ (à $\theta$ constante) & J ; kg ; \si{J.kg^{-1}} \\ \hline
Principe des échanges & $\sum Q_i = 0$ \quad ($Q_{\text{re\c{c}ue}} = -Q_{\text{cédée}}$) & système isolé \\ \hline
Calorimètre & $C_{\text{cal}} = \mu\, c_{\text{eau}}$ & $\mu$ en kg \\ \hline
\end{tabularx}
\end{center}

\textbf{Méthode express (problème de mélange).}
\begin{enumerate}
  \item Faire l'inventaire des corps et des températures ; convertir les masses en kg.
  \item Repérer qui cède ($Q<0$) et qui re\c{c}oit ($Q>0$) ; ne pas oublier le calorimètre ($\mu c_{\text{eau}}\Delta\theta$).
  \item S'il y a changement d'état : vérifier que la température finale le permet, puis ajouter les termes $mL$.
  \item Écrire $\sum Q_i = 0$ et résoudre l'équation en l'inconnue ($\theta_f$, $c$, $\mu$ ou $m$).
  \item Vérifier : résultat entre les températures extrêmes, unité correcte, ordre de grandeur plausible.
\end{enumerate}
\end{bilan}

\begin{aretenir}[Checklist avant le Probatoire]
\begin{itemize}
  \item[$\square$] Je sais distinguer chaleur (énergie échangée, en J) et température (état du corps, en \si{\celsius} ou K).
  \item[$\square$] Je cite les trois modes de transfert et je donne un exemple concret de chacun (barre métallique, eau qui bout, Soleil).
  \item[$\square$] Je sais que le rayonnement est le seul mode possible dans le vide.
  \item[$\square$] Je connais par c\oe ur $Q = mc\,\Delta\theta$ et je sais que $\Delta\theta$ est identique en \si{\celsius} et en K.
  \item[$\square$] Je sais que $Q > 0$ si le corps re\c{c}oit de la chaleur et $Q < 0$ s'il en cède.
  \item[$\square$] Je connais $c_{\text{eau}} = \SI{4185}{J.kg^{-1}.K^{-1}}$ et je convertis toujours les masses en kg.
  \item[$\square$] Je sais qu'un changement d'état se fait à température constante avec $Q = mL$.
  \item[$\square$] Je sais dessiner et annoter un calorimètre (vase Dewar, parois argentées, bouchon, agitateur, thermomètre).
  \item[$\square$] Je sais définir et utiliser la valeur en eau $\mu$ d'un calorimètre.
  \item[$\square$] J'applique le principe des échanges : dans un système isolé, $\sum Q_i = 0$.
  \item[$\square$] Je vérifie chaque résultat : unité SI, signe cohérent, température finale comprise entre les températures initiales.
  \item[$\square$] Je sais reconnaître dans la vie courante des appareils analogues au calorimètre (thermos, glacière, chauffe-eau isolé).
\end{itemize}
\end{aretenir}

\findecours

\end{document}
